\chapter{Formal calculus and locality}\label{ch:L01}
\begingroup\setlength{\emergencystretch}{3em}

\providecommand{\loneW}{\mathcal{W}}
\providecommand{\loneV}{V}
\providecommand{\loneiota}[2]{\iota_{#1,#2}}
\providecommand{\lonedelta}{\delta(z_1,z_2)}
\providecommand{\lonecoef}[1]{\big[#1\big]}
\providecommand{\lonebinom}[2]{\binom{#1}{#2}}
\providecommand{\lonewarn}[1]{the trap box on p.~\pageref{#1}}

This chapter documents the release of the LeanVOA library named in the front matter. Every Lean quotation below is checked by \texttt{scripts/\allowbreak textbook/\allowbreak sync.py}
to occur verbatim (up to whitespace and comments) in the source files of that release. Checks that were run at an earlier state of the library and not re-run
for this edition are marked as such. The chapter covers the thirteen files assigned to it in
\texttt{Notes/\allowbreak Textbook/\allowbreak assignment.txt}, all in the directory
\leanpath{VOA/\allowbreak Foundations/}: \leanpath{FormalSeries.lean}, \leanpath{FormalCalculus.lean},
\leanpath{Locality.lean}, \leanpath{Expansion.lean}, \leanpath{IteratedLaurent.lean},
\leanpath{Substitution.lean}, \leanpath{DeltaJump.lean}, \leanpath{IotaZExpansion.lean},
\leanpath{FormalDist3.lean}, \leanpath{FormalDist4.lean}, \leanpath{FormalDist4Delta.lean},
\leanpath{FormalDist4DeltaContract.lean} and \leanpath{FormalDist4DeltaSupport.lean}.

Chapter~\ref{ch:01} (Section~\ref{ch01:sec:synonyms}) already introduced, briefly, some of the objects
that every later chapter uses: formal distributions, modes, the residue, locality and the
structure theorem for local pairs. The present chapter is the full treatment. It presents
\emph{every} definition and every substantive theorem of the thirteen files, with the complete
fidelity protocol of the Constitution, and it works out by hand the examples that a student
should have computed once in their life: the two expansions of $(z_1-z_2)^{-1}$ and
$(z_1-z_2)^{-2}$, the jump between them, the delta function and its derivatives, the free boson
and free fermion as local pairs, and the substitution $z_1\mapsto z_0+z_2$ applied to an
explicit series. Where Chapter~\ref{ch:01} already printed a Lean declaration we print it
again, because the referee should be able to check this chapter on its own.

\paragraph{What is proved, in one paragraph.} Everything in this chapter is
\emph{unconditional} formal algebra. The headline results are:
the structure theorem for distributions annihilated by a power of $z_1-z_2$
(\lean{VOA.local_iff_deltaSum}, Gannon's Proposition~5.1.2, Kac's Theorem~2.3);
the delta-function jump $\loneiota{z_1}{z_2}(z_1-z_2)^{-1}-\loneiota{z_2}{z_1}(z_1-z_2)^{-1} = \lonedelta$
(\lean{VOA.expand_jump}) and its all-order generalisation to $(z_1-z_2)^{-k-1}$
(\lean{VOA.negPowJump_eq_deltaDeriv});
multiplicativity of the expansions (\lean{VOA.expandPos_mul}, \lean{VOA.expandNeg_mul}), whose
engine is the generalised Chu--Vandermonde identity (\lean{VOA.cbz_vandermonde});
the cancellation theorem saying that $z_1-z_2$ is not a zero-divisor on lower-truncated
distributions (\lean{VOA.eq_zero_of_mulBySubPow_of_bddBelow}), which is the irreducible core of
Goddard uniqueness and of the reconstruction theorem;
the multiplicativity of the substitution $z_1\mapsto z_0+z_2$ (\lean{VOA.subst_mulBySub});
the sign-twisted (super) versions of locality (\lean{VOA.localε_iff_deltaSum},
\lean{VOA.localε_of_car}); and a four-variable double-delta calculus
(\lean{VOA.FormalDist₄Delta.kernelContract₁₂_doubleDelta_eq_zero_of_local}) intended for Li's
local-system theorem. We ran \texttt{\#print axioms} on sixty declarations of these files; every
one depends only on \lean{[propext, Classical.choice, Quot.sound]}
(Section~\ref{l01:sec:verification}). The PBW theorem \lean{VOA.pbw} (Chapter~\ref{ch:L02}; formerly a project axiom, now a theorem) does not occur. The only
\emph{hypotheses} that appear are finite-support hypotheses that make infinite sums meaningful;
each is discussed in its fidelity check, and one of them (\lean{ShiftContractFinite}) is,
at this commit, not yet discharged for the payloads it is intended for
(\lonewarn{l01:warn:shiftfinite}).

\paragraph{How to read this chapter.} Sections~\ref{l01:sec:physics}--\ref{l01:sec:binomial}
are the one-variable theory and the combinatorics of binomial coefficients. Sections
\ref{l01:sec:twovar}--\ref{l01:sec:operators} are the heart: two-variable distributions, the
structure theorem, the two expansions and their jump, and locality of operator-valued
distributions. Sections~\ref{l01:sec:truncation}--\ref{l01:sec:substitution} are the tools for
the reconstruction programme (lower truncation, cancellation, substitution). Sections
\ref{l01:sec:three}--\ref{l01:sec:four} are the three- and four-variable layers. The chapter ends
with the verification record, the referee's checklist and the notes.

%======================================================================
\section{The physical picture: fields, contours and radial ordering}\label{l01:sec:physics}
%======================================================================

Nothing in this section is used in a proof; it explains why the definitions that follow are the
right ones. Chapter~\ref{ch:01}, Section~\ref{ch01:sec:physics}, gives a longer account; here we
concentrate on the four facts that formal calculus encodes.

\subsection{Fields are Laurent series of operators}

In a two-dimensional conformal field theory on the Riemann sphere one chooses a coordinate $z$
centred at the origin and \emph{radially quantises}: circles $|z| = r$ are the equal-time slices,
$r = e^{\tau}$ with $\tau$ the Euclidean time of the cylinder. A chiral (holomorphic) field
$a(z)$ of conformal weight $h$ is expanded in modes,
\begin{equation*}
  a(z) = \sum_{n\in\Z} a_n\, z^{-n-h},\qquad a_n = \oint_{|z| = r}\frac{dz}{2\pi i}\,z^{n+h-1}a(z).
\end{equation*}
The modes $a_n$ are operators on the space of states. Vertex-algebra conventions absorb the
weight: one writes $a(z) = \sum_n a_{(n)} z^{-n-1}$, so that $a_{(n)} = \oint z^{n}a(z)\,dz/2\pi i$
for \emph{every} field, whatever its weight. The library adopts exactly this convention (the
function \lean{VOA.mode} of Definition~\ref{l01:def:mode}). The expression $a(z)$ itself is
never a convergent series of operators: it only makes sense inside correlation functions, or
applied to a state. Formal calculus keeps the series as a \emph{list of coefficients} and never
asks it to converge. This is the content of Definition~\ref{l01:def:formaldist}.

\subsection{Contour integrals are residues}

The contour integral $\oint z^{n}a(z)\,dz/2\pi i$ picks out the coefficient of $z^{-1}$ in
$z^na(z)$. In formal calculus this \emph{is} the definition: $\Res_z f(z)$ is the coefficient of
$z^{-1}$. Two facts of complex analysis survive verbatim: the integral of a total derivative over
a closed contour vanishes ($\Res_z\partial_zf = 0$, Proposition~\ref{l01:prop:moderes}), and one
may integrate by parts (Theorem~\ref{l01:thm:byparts}). The \emph{conserved charge} associated
with a current $J(z)$ is $Q = \oint J(z)\,dz/2\pi i = J_{(0)}$: a residue.

\subsection{Two orderings, one function: why there are two expansions}

The basic object of the theory is a correlation function such as
$\langle a(z_1)b(z_2)\rangle$, which is a meromorphic function of $z_1,z_2$ with poles at
$z_1 = z_2$ (and at $0,\infty$). In radial quantisation, operator products must be
\emph{radially ordered}: the product $a(z_1)b(z_2)$ of operators is defined when $|z_1|>|z_2|$,
the product $b(z_2)a(z_1)$ when $|z_2|>|z_1|$, and both are expansions of the \emph{same}
meromorphic function in two different annuli. The prototype is the function $1/(z_1-z_2)$:
\begin{align*}
  \frac{1}{z_1-z_2} &= \sum_{k\ge0} z_1^{-k-1}z_2^{k}\qquad (|z_1|>|z_2|),\\
  \frac{1}{z_1-z_2} &= -\sum_{k\ge0} z_1^{k}z_2^{-k-1}\qquad (|z_2|>|z_1|).
\end{align*}
These are the two expansions $\loneiota{z_1}{z_2}$ and $\loneiota{z_2}{z_1}$ of
Definition~\ref{l01:def:expand}. Their difference is
$\sum_{k\in\Z}z_1^{k}z_2^{-k-1}$, which is the formal delta function $\lonedelta$: in complex
analysis, the difference of the contour integrals over $|z_1| = R$ and $|z_1| = r<|z_2|<R$ is the
integral over a small circle around $z_1 = z_2$, which evaluates a test function at $z_1 = z_2$.
So \emph{the delta function is the difference of the two radial orderings}, and it is supported
at coincident points: this is the contact term. Theorem~\ref{l01:thm:jump} is the formal
statement, and Theorem~\ref{l01:thm:alljump} the version for higher poles.

\subsection{Locality is commutativity up to contact terms}

Suppose the operator product expansion of $a$ with $b$ has a pole of order at most $N$:
$a(z_1)b(z_2) = \sum_{j<N}c^j(z_2)(z_1-z_2)^{-j-1} + \text{regular}$. Then
$(z_1-z_2)^{N}a(z_1)b(z_2)$ is regular at $z_1 = z_2$, and (for bosonic fields) it is the same
function as $(z_1-z_2)^{N}b(z_2)a(z_1)$: the two orderings agree once the pole is cancelled. In
formal language,
\begin{equation*}
  (z_1-z_2)^N\,[a(z_1),b(z_2)] = 0 .
\end{equation*}
This is \emph{mutual locality}. In Minkowski signature it becomes the statement that the fields
commute at spacelike separation; in the Euclidean, radially ordered picture it says that the
commutator is a contact term. The structure theorem (Theorem~\ref{l01:thm:deltasum}) makes this
precise: a distribution killed by $(z_1-z_2)^N$ is exactly a sum
$\sum_{j<N}c^j(z_2)\,\partial_{z_2}^j\lonedelta$ of derivatives of the delta function. Taking
residues then gives the familiar commutator formula of modes, $[a_{(m)},b_{(n)}] =
\sum_j\binom mj (a_{(j)}b)_{(m+n-j)}$, which is how physicists derive the Virasoro and current
algebras from the singular parts of OPEs.

For fermionic fields the second ordering carries a sign, $b(z_2)a(z_1)\mapsto -b(z_2)a(z_1)$, and
locality is a statement about the \emph{anti}commutator; the library handles both cases with a
scalar parameter $\varepsilon$ (Section~\ref{l01:sec:super}).

\subsection{What the formal theory must avoid}

Two series in the same variable cannot in general be multiplied: $\delta(z)\cdot\delta(z)$ would
have every coefficient equal to $\sum_{k\in\Z}1$. The formal theory must therefore be
\emph{disciplined about infinite sums}: a product or a sum is allowed only when every coefficient
of the result is a finite sum. This discipline appears in the library in three forms:
(i) \lean{FormalDist} carries no multiplication at all (Definition~\ref{l01:def:formaldist});
(ii) sums are taken with Mathlib's \lean{finsum}, and every theorem that manipulates such a sum
carries a finite-support hypothesis; (iii) a definition whose coefficient is a \lean{finsum} may be
\emph{written down} for any input, but is meaningful only on inputs for which the sum is finite.
Point (iii) is a standing trap, explained in \lonewarn{l01:warn:finsum}.

%======================================================================
\section{One-variable formal distributions}\label{l01:sec:onevar}
%======================================================================

\subsection{The space of formal distributions}

\begin{definition}[Formal distribution]\label{l01:def:formaldist}
Let $\loneW$ be a set. A \emph{formal distribution} in one variable $z$ with coefficients in
$\loneW$ is a function $f:\Z\to\loneW$, written as the doubly infinite formal series
\begin{equation*}
  f(z) = \sum_{n\in\Z} f_n\,z^n,\qquad f_n := f(n)\in\loneW .
\end{equation*}
The set of formal distributions is denoted $\loneW[[z^{\pm1}]]$. When $\loneW$ is an abelian group,
respectively a module over a semiring $R$, $\loneW[[z^{\pm1}]]$ is an abelian group, respectively an
$R$-module, under coefficientwise addition and scalar multiplication. Two formal distributions are
equal iff all their coefficients agree. No product of two formal distributions is defined.
\end{definition}

\noindent File \leanpath{VOA/\allowbreak Foundations/\allowbreak FormalSeries.lean}.
\begin{leancode}
def FormalDist (𝒲 : Type*) : Type _ := ℤ → 𝒲

@[inherit_doc] notation:max 𝒲 "⟦z^±⟧" => FormalDist 𝒲

instance instAddCommGroup [AddCommGroup 𝒲] : AddCommGroup (FormalDist 𝒲) :=
  inferInstanceAs (AddCommGroup (ℤ → 𝒲))

instance instModule {R : Type*} [Semiring R] [AddCommGroup 𝒲] [Module R 𝒲] :
    Module R (FormalDist 𝒲) := inferInstanceAs (Module R (ℤ → 𝒲))

@[ext]
theorem ext {f g : FormalDist 𝒲} (h : ∀ n, f n = g n) : f = g := funext h
\end{leancode}

\begin{dictionary}
\textbf{Decoding.} \lean{Type*} means ``some type'' (a set, for our purposes). \lean{ℤ → 𝒲} is the
type of functions from the integers to $\loneW$. So \lean{f : FormalDist 𝒲} is a function and
\lean{f n} is its value at $n$, \emph{which we read as the coefficient of $z^n$}. This is the
first convention to memorise: \lean{f n} is the coefficient of $z^{n}$, not a mode. Modes are
introduced in Definition~\ref{l01:def:mode} and are shifted by one. The notation
\lean{𝒲⟦z^±⟧} is a pure abbreviation for \lean{FormalDist 𝒲}.

\textbf{The two instances.} \lean{instance ... : AddCommGroup (FormalDist 𝒲)} says: if $\loneW$ is an
abelian group, so is \lean{FormalDist 𝒲}, and \lean{inferInstanceAs (AddCommGroup (ℤ → 𝒲))} says that
the group structure is the one of the function space, i.e.\ coefficientwise. Likewise for the
module structure over any semiring $R$ (in practice $R = \C$ or $R = \Z$). The
\lean{@[ext]} theorem is the statement ``equal coefficients imply equal distributions'', registered
so that Lean's \lean{ext} tactic uses it.

\textbf{Why a \lean{def} and not the bare function type.} If $\loneW$ is a ring, the function type
\lean{ℤ → 𝒲} carries the pointwise product $(fg)_n = f_ng_n$, which has nothing to do with the
product of series. Because \lean{FormalDist} is a separate definition (Lean does not unfold it
when searching for instances), the pointwise product is never found. The file header says this
explicitly: there is ``deliberately \textbf{no} ring structure''. No multiplication of
distributions is defined anywhere in the library; products are formed by explicit, separately
justified constructions (Definitions~\ref{l01:def:lpsmul}, \ref{l01:def:convol},
\ref{l01:def:subst}).

\textbf{Fidelity.} Same as Kac (\emph{Vertex algebras for beginners}, §2.1), Gannon (§5.1.2,
the space $\mathcal W[[z^{\pm}]]$) and Lepowsky--Li (§2.1). \textbf{Could it be vacuous?} No: it is
the full space of $\Z$-indexed families.
\end{dictionary}

\begin{physics}
A formal distribution with values in $\End(V)$ is the mode expansion of a quantum field, kept as
the list of its Fourier--Laurent modes; a distribution with values in $V$ is a field applied to a
state, $a(z)v$; a distribution with values in $\C$ is a correlation function expanded around a
point. No convergence is ever claimed, because in quantum field theory none is available at the
level of operators: only matrix elements and their Laurent coefficients (OPE coefficients) are
meaningful. The refusal to define a product mirrors the physical fact that two operators at the
\emph{same} point cannot simply be multiplied: their product is singular and must be regularised
(normal ordering, Definition~\ref{l01:def:bracket}).
\end{physics}

\begin{status}
A definition with two instances and an extensionality lemma; no proof obligation beyond
type-checking. Compiled as part of the built tree at an earlier commit, when this chapter was first written (not re-run for this edition).
\end{status}

\subsection{Laurent polynomials and algebraically summable families}

The \emph{test functions} of the theory are the Laurent polynomials. Pairing a distribution with
a test function, $\Res_z\, g(z)f(z)$, is always a finite sum: this is the sense in which formal
distributions are ``distributions''.

\begin{definition}[Laurent polynomials]\label{l01:def:laurent}
Let $\loneW$ be a complex vector space. A formal distribution $f\in\loneW[[z^{\pm1}]]$ is a
\emph{Laurent polynomial} if its support $\{n\in\Z : f_n\neq0\}$ is finite. Laurent
polynomials form a $\C$-linear subspace $\loneW[z^{\pm1}]\subseteq\loneW[[z^{\pm1}]]$.
\end{definition}

\noindent File \leanpath{VOA/\allowbreak Foundations/\allowbreak FormalSeries.lean}.
\begin{leancode}
def LaurentPoly (𝒲 : Type*) [AddCommGroup 𝒲] [Module ℂ 𝒲] : Submodule ℂ (FormalDist 𝒲) where
  carrier := {f | (Function.support f).Finite}

theorem mem_laurentPoly [AddCommGroup 𝒲] [Module ℂ 𝒲] {f : FormalDist 𝒲} :
    f ∈ LaurentPoly 𝒲 ↔ (Function.support f).Finite := Iff.rfl
\end{leancode}

\begin{dictionary}
\lean{Submodule ℂ M} is Mathlib's notion of a $\C$-linear subspace of $M$; it is a structure
consisting of a carrier set together with proofs that it contains $0$ and is closed under
addition and scalar multiplication (those three proofs are in the file and are omitted here).
\lean{Function.support f} is $\{n : f(n)\neq0\}$ and \lean{.Finite} says this set is finite.
\lean{mem_laurentPoly} records that membership is literally finiteness of the support
(\lean{Iff.rfl}: the two sides are the same statement). \textbf{Fidelity:} same as Gannon §5.1.2
(the ``test functions'' $\mathcal W[z^{\pm}]$). In the library, theorems that need a Laurent
polynomial usually ask directly for \lean{(Function.support g).Finite} rather than
\lean{g ∈ LaurentPoly ℂ}; the two are interchangeable by \lean{mem_laurentPoly}.
\end{dictionary}

\begin{physics}
Laurent polynomials are the smearing functions of chiral CFT: the mode $a_{(n)} = \Res_z z^na(z)$ is
the field smeared against the test function $z^n$, and a general smeared field
$\Res_z g(z)a(z)$ is a finite linear combination of modes. Physically, one never measures a field
at a point, only smeared against a test function; for chiral fields on the circle the natural test
functions are trigonometric polynomials $e^{in\sigma}$, i.e.\ Laurent polynomials in $z = e^{i\sigma}$.
\end{physics}

\begin{status}
Definition. The closure proofs (\lean{add_mem'}, \lean{zero_mem'}, \lean{smul_mem'}) are in the file and
compile; no hypotheses. The structure \lean{LaurentPoly} is used downstream in
\leanpath{VOA/\allowbreak Examples/\allowbreak VirasoroKacLaurent.lean}.
\end{status}

\begin{definition}[Algebraically summable family; algebraic sum]\label{l01:def:algsum}
Let $\loneW$ be an abelian group and $I$ any index set. A family
$(f^{(i)})_{i\in I}$ in $\loneW[[z^{\pm1}]]$ is \emph{algebraically summable} if for every $n\in\Z$
the set $\{i\in I : f^{(i)}_n\neq0\}$ is finite. Its \emph{algebraic sum} is the distribution
\begin{equation*}
  \Big(\sum_{i\in I}f^{(i)}\Big)_n := \sum_{i\,:\,f^{(i)}_n\neq0} f^{(i)}_n ,
\end{equation*}
a finite sum for each $n$.
\end{definition}

\noindent File \leanpath{VOA/\allowbreak Foundations/\allowbreak FormalSeries.lean}.
\begin{leancode}
def AlgSummable {ι : Type*} (f : ι → FormalDist 𝒲) : Prop := ∀ n, {i | f i n ≠ 0}.Finite

noncomputable def algSum {ι : Type*} (f : ι → FormalDist 𝒲) (_h : AlgSummable f) :
    FormalDist 𝒲 := fun n => ∑ᶠ i, f i n
\end{leancode}

\begin{dictionary}
\textbf{Decoding.} \lean{ι} is the index type $I$; \lean{f : ι → FormalDist 𝒲} is the family;
\lean{f i n} is $f^{(i)}_n$. \lean{AlgSummable f} is a proposition (\lean{Prop}): ``for all $n$, the set of
$i$ with $f^{(i)}_n\neq0$ is finite''. \lean{∑ᶠ i, f i n} is Mathlib's \lean{finsum}: the sum of a
family over its support \emph{if the support is finite}, and $0$ otherwise
(\lonewarn{l01:warn:finsum}).

\textbf{Every hypothesis.} \lean{algSum} takes the summability proof \lean{_h} as an argument. It is not
used in the formula (hence the underscore), but its presence means that one cannot even write
\lean{algSum f h} without having proved that every coefficient is a finite sum. This is exactly the
discipline that rules out the junk value of \lean{finsum}: with \lean{_h} available, \lean{∑ᶠ i, f i n} is
a genuine finite sum.

\textbf{Fidelity.} Same as Gannon p.~315 (``a sum is algebraically defined iff for each fixed $n$ the
coefficient is a finite sum''). \lean{noncomputable} only means that Lean cannot run the definition as a
program, because \lean{finsum} decides finiteness by classical logic; it has no mathematical content.
\end{dictionary}

\begin{physics}
Algebraic summability is the formal version of the energy bound that makes operator products well
defined on states. For example $Y(u,z)Y(v,w)x = \sum_{m,n}u_{(m)}v_{(n)}x\,z^{-m-1}w^{-n-1}$
is a meaningful formal series because each coefficient is a single vector; but the same-point
product $\sum_m u_{(m)}v_{(s-m)}x$ needed for a normal-ordered product is an infinite sum over $m$,
and is meaningful only because, on a fixed state of bounded energy, all but finitely many terms vanish.
\end{physics}

\begin{status}
Definitions; no proof obligations. Not used outside \leanpath{FormalSeries.lean} at this commit
(later layers use \lean{finsum} directly, with explicit support hypotheses).
\end{status}

\begin{proposition}[Algebraic sums are linear]\label{l01:prop:algsum}
Every family indexed by a finite set is algebraically summable, and its algebraic sum is the
ordinary finite sum. If $(f^{(i)})$ and $(g^{(i)})$ are algebraically summable, so are
$(f^{(i)}+g^{(i)})$ and $(c f^{(i)})$ for $c\in\C$, and
$\sum_i (f^{(i)}+g^{(i)}) = \sum_if^{(i)}+\sum_ig^{(i)}$, $\sum_i cf^{(i)} = c\sum_if^{(i)}$.
\end{proposition}

\noindent File \leanpath{VOA/\allowbreak Foundations/\allowbreak FormalSeries.lean}.
\begin{leancode}
theorem algSummable_of_finite {ι : Type*} [Finite ι] (f : ι → FormalDist 𝒲) :
    AlgSummable f := fun _ => Set.toFinite _

theorem algSum_apply_fintype {ι : Type*} [Fintype ι] (f : ι → FormalDist 𝒲)
    (h : AlgSummable f) (n : ℤ) : algSum f h n = ∑ i, f i n := by

theorem AlgSummable.add {ι : Type*} {f g : ι → FormalDist 𝒲}
    (hf : AlgSummable f) (hg : AlgSummable g) : AlgSummable (fun i => f i + g i) := by

theorem algSum_add {ι : Type*} {f g : ι → FormalDist 𝒲}
    (hf : AlgSummable f) (hg : AlgSummable g) :
    algSum (fun i => f i + g i) (hf.add hg) = algSum f hf + algSum g hg := by

theorem algSum_smul {ι : Type*} (c : ℂ) {f : ι → FormalDist 𝒲} (hf : AlgSummable f) :
    algSum (fun i => c • f i) (hf.smul c) = c • algSum f hf := by
\end{leancode}

\begin{dictionary}
\lean{[Finite ι]} and \lean{[Fintype ι]} are the hypotheses ``$I$ is finite'' (the second also supplies an
enumeration, needed to write the ordinary sum \lean{∑ i, f i n}). \lean{hf.add hg} is the proof that the
sum family is summable (the theorem \lean{AlgSummable.add}); it has to be supplied because \lean{algSum}
requires a summability proof. \lean{c • f i} is scalar multiplication. All statements are the expected
ones; nothing is weaker or stronger.
\end{dictionary}

\begin{physics}
Technical bookkeeping: it says that the space of well-defined (energy-bounded) infinite sums of
fields is a vector space, so that one may add and rescale OPEs and mode expansions freely.
\end{physics}

\textbf{Proof on paper.} For each $n$ the support of $(f^{(i)}_n+g^{(i)}_n)_i$ is contained in the
union of the two finite supports; the identity is then the additivity of finite sums. Scaling can
only shrink supports. For finite $I$ there is nothing to prove. \qed

\begin{status}
Unconditional. \texttt{\#print axioms VOA.algSum\_add} $=$ \lean{[propext, Classical.choice, Quot.sound]}
(Section~\ref{l01:sec:verification}).
\end{status}

\begin{warning}\label{l01:warn:finsum}
\textbf{The junk value of \texttt{finsum}.} Mathlib defines \lean{∑ᶠ i, F i} to be the sum over the support
of \lean{F} when that support is finite, and \emph{zero} when it is infinite. This is a convention, chosen
so that \lean{finsum} is a total function; it is not a theorem about convergence. Consequently a Lean
definition whose coefficients are \lean{finsum}s can be applied to \emph{any} input, and on an input for
which the intended sum is infinite it silently returns $0$, a meaningless value. Three definitions of
this chapter are of this kind: \lean{lpSMul} (Definition~\ref{l01:def:lpsmul}), \lean{subst}
(Definition~\ref{l01:def:subst}), and the contractions of Section~\ref{l01:sec:four}. In each case the
\emph{theorems} carry the finite-support hypothesis that makes the sum genuine, and the fidelity check
below verifies this for every theorem. A reader who meets such a definition applied to an
unconstrained input should be suspicious. The authors of \leanpath{FormalDist4.lean} record the same
point in their header: finite-support hypotheses ``prevent an infinite-support \lean{finsum} from silently
collapsing to zero''.
\end{warning}

\subsection{Multiplication by a monomial}

\begin{definition}[Multiplication by $z^k$]\label{l01:def:monomial}
For $k\in\Z$, multiplication by the monomial $z^k$ is the $\C$-linear map
$\loneW[[z^{\pm1}]]\to\loneW[[z^{\pm1}]]$, $(z^kf)_n = f_{n-k}$.
\end{definition}

\noindent File \leanpath{VOA/\allowbreak Foundations/\allowbreak FormalSeries.lean}.
\begin{leancode}
def mulByMonomial (k : ℤ) : FormalDist 𝒲 →ₗ[ℂ] FormalDist 𝒲 where
  toFun f := fun n => f (n - k)
\end{leancode}

\begin{dictionary}
\lean{M →ₗ[ℂ] N} is the type of $\C$-linear maps from $M$ to $N$; a term of it is a function
(\lean{toFun}) together with proofs of additivity and homogeneity (\lean{map_add'}, \lean{map_smul'}, both
\lean{rfl} in the file and omitted here). The coefficient of $z^n$ in $z^k\sum_mf_mz^m$ is $f_{n-k}$, which is
the formula. Same as the textbook.
\end{dictionary}

\begin{physics}
Multiplying a field by $z^k$ shifts its modes: $(z^ka)_{(m)} = a_{(m+k)}$. In radial quantisation,
$z^k$ is the eigenfunction of the dilatation $L_0$ with eigenvalue $k$, so this operation shifts the
energy carried by each mode.
\end{physics}

\begin{status}
Definition; linearity is by \lean{rfl}. No hypotheses.
\end{status}

%======================================================================
\section{Delta function, residue, derivative and modes}\label{l01:sec:delta}
%======================================================================

\begin{definition}[Formal delta functions]\label{l01:def:delta}
The \emph{formal delta function} in one variable is $\delta(z) := \sum_{n\in\Z}z^n\in\C[[z^{\pm1}]]$.
The two-variable delta function in multiplicative normalisation is
\begin{equation*}
  \delta(z_1/z_2) := \sum_{k\in\Z}z_1^{k}z_2^{-k},\qquad\text{coefficient of }z_1^mz_2^n = [m+n = 0].
\end{equation*}
\emph{Notation for the rest of the chapter:} the symbol $\delta(z_1,z_2)$ (with a comma) always denotes the
\emph{additive} normalisation $z_2^{-1}\delta(z_1/z_2) = \sum_kz_1^kz_2^{-k-1}$, coefficient $[m+n+1 = 0]$, which is
Kac's $\delta(z-w)$ and the only one used in the locality theorems (Lean: \lean{deltaForm}, or
\lean{deltaDeriv 0 c}). The multiplicative $\delta(z_1/z_2)$ above (Lean: \lean{delta₂}) is not used again.
\end{definition}

\noindent File \leanpath{VOA/\allowbreak Foundations/\allowbreak FormalCalculus.lean}.
\begin{leancode}
def formalDelta : FormalDist ℂ := fun _ => 1

def delta₂ : FormalDist₂ ℂ := fun m n => if m + n = 0 then 1 else 0
\end{leancode}

\begin{dictionary}
\lean{fun _ => 1} is the constant function $1$: every coefficient of $\delta(z)$ is $1$. \lean{FormalDist₂ ℂ}
(Definition~\ref{l01:def:formaldist2}) is the type of two-variable distributions, \lean{f m n} being the
coefficient of $z_1^mz_2^n$; \lean{if m + n = 0 then 1 else 0} is the Iverson bracket $[m+n=0]$.

\textbf{Two normalisations of the two-variable delta.} The library contains \emph{two} two-variable deltas
and the referee must keep them apart. \lean{delta₂} is $\delta(z_1/z_2) = \sum_kz_1^kz_2^{-k}$
(coefficient $[m+n=0]$). The delta used in all the locality theorems is
$\lonedelta := z_2^{-1}\delta(z_1/z_2) = \sum_kz_1^kz_2^{-k-1}$ (coefficient $[m+n+1=0]$), which appears
as \lean{deltaDeriv 0 c} (Definition~\ref{l01:def:deltaderiv}) and as \lean{deltaForm}
(Definition~\ref{l01:def:expand}); it is Kac's $\delta(z-w)$. The two differ by the factor $z_2^{-1}$, a
shift of the second index by one. \lean{formalDelta} and \lean{delta₂} are not used outside
\leanpath{FormalCalculus.lean}; they are faithful to Gannon (5.1.2), where $\delta(z) = \sum_nz^n$ is ``the
Dirac delta centred at $z = 1$'' in multiplicative notation.
\end{dictionary}

\begin{physics}
On the unit circle $z = e^{i\sigma}$, $\delta(z) = \sum_n e^{in\sigma} = 2\pi\,\delta_{\mathrm{Dirac}}(\sigma)$:
the formal delta is the Fourier series of the Dirac delta at $\sigma = 0$, i.e.\ at $z = 1$. The two-variable
version $\delta(z_1/z_2)$ is supported at $z_1 = z_2$: it is the \emph{contact term} that appears in
equal-time commutators $[a(\sigma_1),b(\sigma_2)]\propto\delta(\sigma_1-\sigma_2)$ of canonical quantisation.
\end{physics}

\begin{status}
Definitions. A smoke test in the file checks \lean{res formalDelta = 1} by \lean{rfl}, and that
\lean{mulBySubPow 1 delta₂ = 0}, i.e.\ $(z_1-z_2)\delta(z_1/z_2) = 0$.
\end{status}

\begin{definition}[Modes, residue, derivative]\label{l01:def:mode}
Let $f\in\loneW[[z^{\pm1}]]$.
\begin{enumerate}[label=(\roman*)]
\item The \emph{$m$-th mode} of $f$ ($m\in\Z$) is $f_{(m)} := f_{-m-1}$, so that
  $f(z) = \sum_{m\in\Z}f_{(m)}\,z^{-m-1}$.
\item If $\loneW$ is a complex vector space, the \emph{residue} is the linear functional
  $\Res_z f := f_{-1}$, the coefficient of $z^{-1}$.
\item The \emph{formal derivative} is the linear map $\partial_z$ with
  $(\partial_zf)_n := (n+1)f_{n+1}$, i.e.\ $\partial_z\sum_n f_nz^n = \sum_n nf_nz^{n-1}$.
\end{enumerate}
\end{definition}

\noindent File \leanpath{VOA/\allowbreak Foundations/\allowbreak FormalCalculus.lean}.
\begin{leancode}
def mode (f : FormalDist 𝒲) (m : ℤ) : 𝒲 := f (-m - 1)

def res : FormalDist 𝒲 →ₗ[ℂ] 𝒲 where
  toFun f := f (-1)

def deriv : FormalDist 𝒲 →ₗ[ℂ] FormalDist 𝒲 where
  toFun f := fun n => (n + 1 : ℤ) • f (n + 1)
\end{leancode}

\begin{dictionary}
\textbf{The single most important convention of the library:} \lean{mode f m = f (-m - 1)}. If
\lean{Y u : FormalDist (Module.End ℂ V)} is the vertex operator of a state $u$, then \lean{mode (Y u) n} is
$u_{(n)}$, the coefficient of $z^{-n-1}$, whereas \lean{Y u n} is the coefficient of $z^n$, i.e.\
$u_{(-n-1)}$. Calibration: \lean{Y u 0} $= u_{(-1)}$, \lean{Y u (-1)} $= u_{(0)}$, \lean{Y u (-2)} $= u_{(1)}$.
Both forms occur in later chapters, and every fidelity check there must decode them.

\textbf{\lean{res}} is defined only when $\loneW$ is a complex vector space (the ambient section of the file
assumes \lean{[AddCommGroup 𝒲] [Module ℂ 𝒲]}), because it is packaged as a $\C$-linear map; the formula is
\lean{f (-1)}. \textbf{\lean{deriv}}: \lean{(n + 1 : ℤ) • f (n + 1)} is the integer $n+1$ acting on the vector
$f_{n+1}$ (repeated addition), which agrees with the complex scalar $n+1$. The name \lean{VOA.deriv} lives in
the namespace \lean{VOA} and is unrelated to Mathlib's analytic \lean{deriv}.

\textbf{Fidelity.} Same as Kac §2.1 ($a(z) = \sum a_{(n)}z^{-n-1}$), Gannon (5.1.3a,c), Frenkel--Ben-Zvi
§1.1 and Lepowsky--Li §2.1. In modes, $(\partial f)_{(m)} = (\partial f)_{-m-1} = -m f_{-m} = -m f_{(m-1)}$.
\end{dictionary}

\begin{physics}
$f_{(m)} = \oint z^mf(z)\,dz/2\pi i$ is the contour integral that defines a mode; for a current $J$ the zero
mode $J_{(0)} = \oint J$ is the conserved charge, and for the stress tensor $T$ the modes $T_{(m+1)} = L_m$ are
the Virasoro generators ($L_{-1} = T_{(0)}$ generates translations, $L_0 = T_{(1)}$ dilatations). For a state
$u$, the coefficients $u_{(n)}v$ of $Y(u,z)v$ are the OPE coefficients: $n\ge0$ the singular terms, $n = -1$
the normal-ordered product, $n\le-2$ the descendants. The derivative $\partial_z$ is translation of the
insertion point.
\end{physics}

\begin{status}
Definitions; linearity proofs by \lean{rfl} and by \lean{smul_add}/\lean{smul_comm}. No hypotheses.
\end{status}

\begin{proposition}[Modes are residues; residues of derivatives vanish]\label{l01:prop:moderes}
For every $f\in\loneW[[z^{\pm1}]]$ and $m\in\Z$,
\begin{equation*}
  f_{(m)} = \Res_z\big(z^mf(z)\big),\qquad \Res_z\big(\partial_zf\big) = 0 .
\end{equation*}
\end{proposition}

\noindent File \leanpath{VOA/\allowbreak Foundations/\allowbreak FormalCalculus.lean}.
\begin{leancode}
theorem mode_eq_res (f : FormalDist 𝒲) (m : ℤ) : mode f m = res (mulByMonomial m f) := by

@[simp] theorem res_deriv (f : FormalDist 𝒲) : res (deriv f) = 0 := by
\end{leancode}

\begin{dictionary}
\lean{mulByMonomial m f} is $z^mf$ (Definition~\ref{l01:def:monomial}). \lean{@[simp]} registers the second
identity as a rewriting rule for Lean's simplifier; it has no mathematical content. Same as Gannon
(5.1.3c) and the first identity of (5.1.3b). No hypotheses beyond $\loneW$ being a complex vector space.
\end{dictionary}

\begin{physics}
The first identity is the contour-integral formula for a mode. The second is Cauchy's theorem for
exact forms: the integral of a total derivative around a closed contour vanishes. In CFT it is the
reason that $\partial a$ has no conserved charge, $\oint\partial a = 0$, and that total derivatives can
be dropped from currents without changing charges.
\end{physics}

\textbf{Proof on paper.} $(z^mf)_{-1} = f_{-1-m} = f_{(m)}$. And $(\partial f)_{-1} = (-1+1)f_0 = 0$. \qed

\begin{status}
Unconditional; axioms \lean{[propext, Classical.choice, Quot.sound]} for both.
\end{status}

\begin{example}[Calibrating the conventions]\label{l01:ex:calibrate}
Let $f = \sum_{n}f_nz^n$ and $m\in\Z$.
(a) $\Res_z\delta(z) = 1$, since every coefficient of $\delta$ is $1$.
(b) $(\partial_z z^{k})_n = (n+1)[n+1 = k]$, so $\partial_zz^k = kz^{k-1}$; in particular
$\partial_zz^{0} = 0$ and $\partial_z z^{-1} = -z^{-2}$, and $\Res_z\partial_z z^{0} = 0$ although
$z^{-1}$ is the only monomial with nonzero residue: $z^{-1}$ is not a derivative, which is why the residue
is a nontrivial functional.
(c) $\partial_z\delta(z) = \sum_n(n+1)z^n$; its residue is $0\cdot1 = 0$, as
Proposition~\ref{l01:prop:moderes} demands.
(d) In modes, a field with $a(z) = \sum_na_{(n)}z^{-n-1}$ has $(\partial a)_{(n)} = -n\,a_{(n-1)}$; e.g.\ for
$a = J$ a current, $(\partial J)_{(0)} = 0$: total derivatives carry no charge.
\end{example}

\subsection{The Laurent-polynomial action and integration by parts}

\begin{definition}[Action of a Laurent polynomial]\label{l01:def:lpsmul}
For $g\in\C[[z^{\pm1}]]$ and $f\in\loneW[[z^{\pm1}]]$ define $g\cdot f$ by
\begin{equation*}
  (g\cdot f)_n := \sum_{k\in\Z}g_k\,f_{n-k}.
\end{equation*}
If $g$ is a Laurent polynomial, each coefficient is a finite sum, and $g\cdot f$ is the product of the
Laurent polynomial $g$ with the distribution $f$.
\end{definition}

\noindent File \leanpath{VOA/\allowbreak Foundations/\allowbreak FormalCalculus.lean}.
\begin{leancode}
noncomputable def lpSMul (g : FormalDist ℂ) (f : FormalDist 𝒲) : FormalDist 𝒲 :=
  fun n => ∑ᶠ k, g k • f (n - k)

theorem support_smul_left_finite (g : FormalDist ℂ) (hg : (Function.support g).Finite)
    (F : ℤ → 𝒲) : (Function.support (fun k => g k • F k)).Finite := by
\end{leancode}

\begin{dictionary}
\textbf{The junk-value trap applies here} (\lonewarn{l01:warn:finsum}). \lean{lpSMul g f} is defined for
every $g$, but if $g$ has infinite support the sum \lean{∑ᶠ k, g k • f (n - k)} may have infinite support and
then evaluates to $0$. Example: $g = f = \delta$ gives the ``product'' $\delta\cdot\delta$, which Lean
evaluates to the zero distribution, whereas the mathematical product is undefined (Gannon's Paradox~5.1).
The lemma \lean{support_smul_left_finite} says: if $g$ has finite support, then for \emph{any} $F$ the summand
$k\mapsto g_kF_k$ has finite support. Every theorem about \lean{lpSMul} below carries the hypothesis
\lean{hg : (Function.support g).Finite}, which puts us in the meaningful case. Under it, the definition is
the textbook product (Gannon p.~315: $\mathcal W[[z^{\pm}]]$ is a $\C[z^{\pm}]$-module).
\end{dictionary}

\begin{physics}
Multiplying a field by a Laurent polynomial $g(z)$ is smearing with a test function: the modes of
$g(z)a(z)$ are finite combinations of modes of $a$. This is the operation used in every contour
manipulation, e.g.\ $\oint (z-w)^k a(z)\,dz$ for a polynomial $(z-w)^k$ in $z$.
\end{physics}

\begin{status}
Definition and support lemma; the lemma is unconditional. \lean{lpSMul} is not used outside
\leanpath{FormalCalculus.lean}.
\end{status}

\begin{theorem}[Leibniz rule]\label{l01:thm:leibniz}
If $g\in\C[z^{\pm1}]$ is a Laurent polynomial and $f\in\loneW[[z^{\pm1}]]$, then
\begin{equation*}
  \partial_z(g\cdot f) = (\partial_zg)\cdot f + g\cdot(\partial_zf).
\end{equation*}
\end{theorem}

\noindent File \leanpath{VOA/\allowbreak Foundations/\allowbreak FormalCalculus.lean}.
\begin{leancode}
theorem deriv_lpSMul (g : FormalDist ℂ) (hg : (Function.support g).Finite) (f : FormalDist 𝒲) :
    deriv (lpSMul g f) = lpSMul (deriv g) f + lpSMul g (deriv f) := by
\end{leancode}

\begin{dictionary}
The one hypothesis \lean{hg} is ``$g$ is a Laurent polynomial''. It is necessary for the statement to be
meaningful (\lonewarn{l01:warn:finsum}); note that $\partial_zg$ is then also a Laurent polynomial, so all
three products are genuine. $f$ is arbitrary. Same as Gannon §5.1.2.
\end{dictionary}

\begin{physics}
The product rule for smeared fields; it is what makes the translation operator $L_{-1}$ act as a
derivation on products of fields and test functions.
\end{physics}

\textbf{Proof on paper.} Compare the coefficients of $z^n$. On the left,
$(n+1)\sum_k g_kf_{n+1-k}$. On the right, $\sum_k (k+1)g_{k+1}f_{n-k} + \sum_kg_k(n-k+1)f_{n-k+1}$;
shifting $k\mapsto k-1$ in the first sum turns it into $\sum_k k\,g_kf_{n-k+1}$, and adding gives
$\sum_k(k+n-k+1)g_kf_{n+1-k}$. All sums are finite because $g$ has finite support; the index shift is a
bijection of $\Z$. \qed

\begin{status}
Unconditional (the hypothesis is the meaningfulness condition). Axioms
\lean{[propext, Classical.choice, Quot.sound]}.
\end{status}

\begin{theorem}[Integration by parts; Gannon (5.1.3b)]\label{l01:thm:byparts}
If $g\in\C[z^{\pm1}]$ is a Laurent polynomial and $f\in\loneW[[z^{\pm1}]]$, then
\begin{equation*}
  \Res_z\big(g\cdot\partial_zf\big) = -\Res_z\big((\partial_zg)\cdot f\big).
\end{equation*}
\end{theorem}

\noindent File \leanpath{VOA/\allowbreak Foundations/\allowbreak FormalCalculus.lean}.
\begin{leancode}
theorem res_lpSMul_deriv (g : FormalDist ℂ) (hg : (Function.support g).Finite)
    (f : FormalDist 𝒲) : res (lpSMul g (deriv f)) = - res (lpSMul (deriv g) f) := by
\end{leancode}

\begin{dictionary}
Same hypothesis as Theorem~\ref{l01:thm:leibniz}. The file header writes the identity as
``$\Res_z(g\cdot\partial_zf) = -\Res_z(f\cdot\partial_zg)$''; since \lean{lpSMul} is the commutative product of a
scalar series with a vector series, these are the same statement. Exactly Gannon (5.1.3b).
\end{dictionary}

\begin{physics}
Integration by parts on a closed contour, with no boundary term. In mode language with $g = z^{m}$:
$(\partial f)_{(m)} = -m f_{(m-1)}$, the familiar rule for modes of derivative fields.
\end{physics}

\textbf{Proof on paper.} Apply $\Res_z$ to the Leibniz rule: the left side is the residue of a derivative,
which vanishes (Proposition~\ref{l01:prop:moderes}). \qed

\begin{example}
Take $g = z^m$. Then $(g\cdot\partial f)_{-1} = (\partial f)_{-1-m} = -m f_{-m}$ and
$((\partial g)\cdot f)_{-1} = m\,(z^{m-1}f)_{-1} = mf_{-m}$, confirming the theorem by direct computation.
\end{example}

\begin{status}
Unconditional; axioms \lean{[propext, Classical.choice, Quot.sound]}.
\end{status}

%======================================================================
\section{Binomial coefficients with integer upper index}\label{l01:sec:binomial}
%======================================================================

Expanding $(z_1-z_2)^n$ for negative $n$ requires binomial coefficients $\binom nk$ with $n<0$.
This section collects the library's definitions and the three identities used later: the values at
$n = -1$, Pascal's rule, and the Chu--Vandermonde convolution.

\begin{definition}[Generalised binomial coefficients, signs, the binomial convention]\label{l01:def:gbinom}
\leavevmode
\begin{enumerate}[label=(\roman*)]
\item For $n\in\Z$ and $k\in\N$,
  $\displaystyle\binom nk := \frac{n(n-1)\cdots(n-k+1)}{k!} = \frac1{k!}\prod_{i=0}^{k-1}(n-i)$
  (the empty product, $k = 0$, gives $\binom n0 = 1$).
\item For $n,k\in\Z$, $C(n,k) := \binom nk$ if $k\ge0$ and $C(n,k) := 0$ if $k<0$.
\item For $n\in\Z$, $(-1)^n := 1$ if $n$ is even and $-1$ if $n$ is odd.
\item (Gannon (5.1.4), the \emph{binomial convention}.) For $n\in\Z$,
  $\displaystyle(z_1+z_2)^n := \sum_{k\ge0}\binom nk z_1^{n-k}z_2^{k}$, always expanded in
  \emph{nonnegative powers of the second variable}.
\end{enumerate}
\end{definition}

\noindent Files (all in \leanpath{VOA/\allowbreak Foundations/}): \leanpath{FormalCalculus.lean}
(\lean{gBinom}, \lean{binomialExpand}), \leanpath{Substitution.lean} (\lean{cbz}), \leanpath{Expansion.lean}
(\lean{negOnePow}).
\begin{leancode}
noncomputable def gBinom (n : ℤ) (k : ℕ) : ℂ :=
  (∏ i ∈ Finset.range k, ((n : ℂ) - (i : ℂ))) / (k.factorial : ℂ)

noncomputable def binomialExpand (n : ℤ) : FormalDist₂ ℂ :=
  fun a b => if 0 ≤ b ∧ a = n - b then gBinom n b.toNat else 0

noncomputable def cbz (n k : ℤ) : ℂ := if 0 ≤ k then gBinom n k.toNat else 0

def negOnePow (n : ℤ) : ℂ := if Even n then 1 else -1
\end{leancode}

\begin{dictionary}
\textbf{Decoding.} \lean{Finset.range k} is $\{0,1,\dots,k-1\}$, so the product is
$\prod_{i=0}^{k-1}(n-i)$, computed in $\C$ after the casts \lean{(n : ℂ)}, \lean{(i : ℂ)}; \lean{k.factorial}
is $k!$. \lean{b.toNat} converts an integer $b$ to a natural number (it returns $0$ for $b<0$, but it is only
ever evaluated under the guard $0\le b$, so no junk value is used). \lean{binomialExpand n a b} is the
coefficient of $z_1^az_2^b$ in $(z_1+z_2)^n$: it is $\binom nb$ when $b\ge0$ and $a = n-b$, else $0$. This
is (iv) verbatim. \lean{cbz} is (ii), and \lean{negOnePow} is (iii) (it avoids the integer power
\lean{zpow}, whose sign conventions are a nuisance in Lean).

\textbf{Fidelity.} (i) agrees with the usual binomial coefficient for $n\ge0$ (e.g.\ $\binom 52 = 10$,
$\binom 27 = 0$ since the product contains the factor $2-2$), and is the standard extension to negative
$n$ (Kac §2.1, Gannon (5.1.4a), Lepowsky--Li §2.1). The convention (iv) is genuinely a
\emph{convention}: the same symbol $(z_1+z_2)^n$ with the variables written in the other order means a
different series (Example~\ref{l01:ex:paradox}).

\textbf{Usage.} \lean{binomialExpand} is not used elsewhere at this commit: the library works instead
with the expansions \lean{expandPos}/\lean{expandNeg} of $(z_1-z_2)^n$ (Definition~\ref{l01:def:expand}),
which build in the sign. \lean{cbz} is used heavily (substitution, Vandermonde, the Borcherds identity in
\leanpath{VOA/\allowbreak Core/\allowbreak Skew.lean}; the Borcherds identity is a theorem for
\lean{VertexOperatorAlgebra} only, Chapter~\ref{ch:L03}).
\end{dictionary}

\begin{physics}
Generalised binomials are the Taylor coefficients of $(1+x)^n$ for negative $n$: they are what one gets
by expanding an OPE singularity $(z-w)^{-k}$ in one of the two radial orderings. The convention
``nonnegative powers of the second variable'' fixes, once and for all, which region of the
$(z_1,z_2)$-plane a formula refers to; in a physics computation the region is fixed by the radial
ordering of the operators, and forgetting it is the most common source of sign errors in mode
calculations.
\end{physics}

\begin{status}
Definitions. No hypotheses.
\end{status}

\begin{example}[Values with negative upper index]\label{l01:ex:negbinom}
For $r\ge1$, $\binom{-r}{k} = \frac{(-r)(-r-1)\cdots(-r-k+1)}{k!} = (-1)^k\binom{r+k-1}{k}$. Thus
$\binom{-1}{k} = (-1)^k$, $\binom{-2}{k} = (-1)^k(k+1)$, $\binom{-3}{k} = (-1)^k\binom{k+2}{2}$. In
particular $\binom nk\neq0$ for every $k\ge0$ when $n<0$: binomial sums with negative upper index are
infinite, which is why the expansions of $(z_1-z_2)^{n}$, $n<0$, are infinite series and must be
controlled by truncation.
\end{example}

\begin{example}[The convention determines the series (Gannon's ``Paradox 5.2'')]\label{l01:ex:paradox}
By (iv), $(z_1+z_2)^{-1} = \sum_{k\ge0}(-1)^kz_1^{-1-k}z_2^{k}$, while
$(z_2+z_1)^{-1} = \sum_{k\ge0}(-1)^kz_2^{-1-k}z_1^{k}$. The two formal series are different although the
rational functions are equal. There is no contradiction (Gannon's word ``paradox'' is rhetorical): the
convention attaches to the \emph{ordered} expression a unique series, and the two orders are two different
expansions of one rational function. Their difference is
\begin{equation*}
  (z_1+z_2)^{-1}-(z_2+z_1)^{-1} = \sum_{k\in\Z}(-1)^kz_1^{-1-k}z_2^{k} = z_1^{-1}\delta(-z_2/z_1),
\end{equation*}
a delta function supported at $z_1 = -z_2$, the pole of the rational function. In the library this
phenomenon, for $(z_1-z_2)^{-1}$, is Theorem~\ref{l01:thm:jump}.
\end{example}

\begin{proposition}[Basic values and sign rules]\label{l01:prop:binomvals}
For all $n\in\Z$, $k\in\N$: $\binom n0 = 1$; $\binom{-1}k = (-1)^k$; $k!\binom nk = \prod_{i<k}(n-i)$;
$(-1)^{n+1} = -(-1)^n$; $(-1)^{k} $ (as an integer exponent) equals $(-1)^k$ computed as a power in $\C$;
and $(-1)^{x+y} = (-1)^x(-1)^y$ for $x,y\in\Z$.
\end{proposition}

\noindent Files \leanpath{VOA/\allowbreak Foundations/\allowbreak FormalCalculus.lean},
\leanpath{VOA/\allowbreak Foundations/\allowbreak Expansion.lean},
\leanpath{VOA/\allowbreak Foundations/\allowbreak Substitution.lean},
\leanpath{VOA/\allowbreak Foundations/\allowbreak IotaZExpansion.lean}.
\begin{leancode}
@[simp] theorem gBinom_zero (n : ℤ) : gBinom n 0 = 1 := by

theorem gBinom_neg_one (k : ℕ) : gBinom (-1) k = (-1 : ℂ) ^ k := by

theorem gBinom_mul_factorial (n : ℤ) (k : ℕ) :
    (k.factorial : ℂ) * gBinom n k = ∏ i ∈ Finset.range k, ((n : ℂ) - (i : ℂ)) := by

theorem negOnePow_succ (n : ℤ) : negOnePow (n + 1) = - negOnePow n := by

theorem negOnePow_natCast (k : ℕ) : negOnePow (k : ℤ) = (-1 : ℂ) ^ k := by

theorem negOnePow_add (x y : ℤ) : negOnePow (x + y) = negOnePow x * negOnePow y := by
\end{leancode}

\begin{dictionary}
\lean{(-1 : ℂ) ^ k} is the natural-number power in $\C$. \lean{(k : ℤ)} is $k$ viewed as an integer. All six
are the stated identities, with no hypotheses.
\end{dictionary}

\begin{physics}
Bookkeeping for signs; $\binom{-1}{k} = (-1)^k$ is the geometric series $1/(1+x) = \sum(-x)^k$ behind the
expansion of the simple pole $1/(z-w)$.
\end{physics}

\textbf{Proof on paper.} $\prod_{i<k}(-1-i) = (-1)^k k!$. The sign rules follow by induction on $y$ in
both directions. \qed

\begin{status}
All unconditional; \texttt{\#print axioms VOA.gBinom\_neg\_one} $=$
\lean{[propext, Classical.choice, Quot.sound]}.
\end{status}

\begin{proposition}[Pascal's rule]\label{l01:prop:pascal}
For all $n\in\Z$ and $j\in\N$, $\binom n{j+1} = \binom{n-1}{j+1}+\binom{n-1}{j}$. Equivalently, for all
$n,k\in\Z$, $C(n,k) = C(n-1,k)+C(n-1,k-1)$.
\end{proposition}

\noindent File \leanpath{VOA/\allowbreak Foundations/\allowbreak Substitution.lean}.
\begin{leancode}
theorem gBinom_succ_pascal (n : ℤ) (j : ℕ) :
    gBinom n (j + 1) = gBinom (n - 1) (j + 1) + gBinom (n - 1) j := by

theorem cbz_pascal (n k : ℤ) : cbz n k = cbz (n - 1) k + cbz (n - 1) (k - 1) := by
\end{leancode}

\begin{dictionary}
The integer-indexed form \lean{cbz_pascal} holds for \emph{all} $k\in\Z$, including $k = 0$ (where it
reads $1 = 1+0$) and $k<0$ (where it reads $0 = 0+0$); this uniformity is what lets later proofs shift
summation indices over all of $\Z$ without case distinctions. Same as the classical rule.
\end{dictionary}

\begin{physics}
Pascal's rule is the coefficient form of $(1+x)^n = (1+x)^{n-1}(1+x)$; in the formal calculus it is the
identity $(z_0+z_2)^{m} = (z_0+z_2)^{m-1}z_0 + (z_0+z_2)^{m-1}z_2$, used to move a factor
$z_1-z_2 = z_0$ through the substitution $z_1 = z_0+z_2$ (Theorem~\ref{l01:thm:subst}).
\end{physics}

\textbf{Proof on paper.} Multiply by $(j+1)!$:
$\prod_{i\le j}(n-i) = n\prod_{i<j}(n-1-i)$, and
$(n-1-j)\prod_{i<j}(n-1-i) + (j+1)\prod_{i<j}(n-1-i) = n\prod_{i<j}(n-1-i)$. \qed

\begin{status}
Unconditional; axioms \lean{[propext, Classical.choice, Quot.sound]} (both checked).
\end{status}

\begin{theorem}[Generalised Chu--Vandermonde identity]\label{l01:thm:vandermonde}
For all integers $m,n,k$,
\begin{equation*}
  \sum_{i\in\Z}C(m,i)\,C(n,k-i) = C(m+n,k),
\end{equation*}
the sum having finitely many nonzero terms ($0\le i\le k$).
\end{theorem}

\noindent File \leanpath{VOA/\allowbreak Foundations/\allowbreak IotaZExpansion.lean}.
\begin{leancode}
theorem cbz_convol_support_finite (m n k : ℤ) :
    (Function.support (fun i : ℤ => cbz m i * cbz n (k - i))).Finite := by

noncomputable def cbzConvol (m n k : ℤ) : ℂ := ∑ᶠ i : ℤ, cbz m i * cbz n (k - i)

theorem cbz_vandermonde (m n k : ℤ) : cbzConvol m n k = cbz (m + n) k := by
\end{leancode}

\begin{dictionary}
\lean{cbzConvol} is a \lean{finsum}, but the preceding lemma shows that its support is finite for every
$m,n,k$, so here the junk value never occurs and \lean{cbzConvol m n k} is the genuine finite sum. The
theorem holds for \emph{all} integers $m,n,k$, including negative ones; for $k<0$ both sides are $0$.
This is the full generalised Vandermonde identity, i.e.\ the coefficient form of
$(1+x)^m(1+x)^n = (1+x)^{m+n}$ in $\C[[x]]$ for integer exponents.
\end{dictionary}

\begin{physics}
Coefficientwise, this is the statement that singular OPE factors multiply as functions:
$(z-w)^m(z-w)^n = (z-w)^{m+n}$ inside a fixed radial ordering. It is the bookkeeping behind the
lattice vertex operators, whose OPEs carry the factor $(z-w)^{\langle\alpha,\beta\rangle}$, and which must
satisfy $(z-w)^{\langle\alpha,\beta\rangle}(z-w)^{\langle\alpha,\gamma\rangle} = (z-w)^{\langle\alpha,\beta+\gamma\rangle}$
(Theorem~\ref{l01:thm:mult}).
\end{physics}

\textbf{Proof on paper.} Fix $k\ge0$ (for $k<0$ both sides vanish). Both sides are polynomials in
$(m,n)$: the left is $\sum_{i=0}^{k}\binom mi\binom n{k-i}$, the right $\binom{m+n}k$. For natural
$m,n$ the identity is classical (choose $k$ objects from $m+n$ by deciding how many come from the first
$m$). Two polynomials in two variables that agree on $\N\times\N$ are equal, so the identity holds for all
integers. An alternative argument, closer to the library's, is to note that both sides satisfy the
recursion $V(m+1) = V(m)+V(m)|_{k\mapsto k-1}$ in $m$ (Pascal's rule) and agree at $m = 0$; this determines
them for $m\ge0$, and running the recursion backwards (using that both vanish for $k<0$) determines them
for $m<0$. The library formalises this second argument by integer induction on $m$ with an inner
induction on $k$. \qed

\begin{example}
$m = n = -1$, $k = 2$: $\sum_{i=0}^2\binom{-1}{i}\binom{-1}{2-i} = 1+1+1 = 3$ and $\binom{-2}{2} = 3$. And
$m = 2$, $n = -1$, $k = 3$: $\binom20\binom{-1}3+\binom21\binom{-1}2+\binom22\binom{-1}1 = -1+2-1 = 0 =
\binom13$, as it must be since $(1+x)^2(1+x)^{-1} = 1+x$.
\end{example}

\begin{status}
Unconditional; \texttt{\#print axioms VOA.cbz\_vandermonde} $=$ \lean{[propext, Classical.choice, Quot.sound]}.
\end{status}

%======================================================================
\section{Two-variable distributions and the structure theorem}\label{l01:sec:twovar}
%======================================================================

\subsection{Two-variable distributions and multiplication by \texorpdfstring{$z_1-z_2$}{z1-z2}}

\begin{definition}[Two-variable distributions; $(z_1-z_2)^N$]\label{l01:def:formaldist2}
A \emph{two-variable formal distribution} is a function $f:\Z\times\Z\to\loneW$, written
$f(z_1,z_2) = \sum_{m,n\in\Z}f_{m,n}z_1^mz_2^n$; the space is $\loneW[[z_1^{\pm1},z_2^{\pm1}]]$, a vector space
under coefficientwise operations. Multiplication by $z_1$, by $z_2$ and by $z_1-z_2$ are the linear maps
\begin{equation*}
  (z_1f)_{m,n} = f_{m-1,n},\qquad (z_2f)_{m,n} = f_{m,n-1},\qquad
  ((z_1-z_2)f)_{m,n} = f_{m-1,n}-f_{m,n-1},
\end{equation*}
and $(z_1-z_2)^N$ ($N\in\N$) is the $N$-fold iterate of the last, $(z_1-z_2)^0 = \mathrm{id}$.
\end{definition}

\noindent File \leanpath{VOA/\allowbreak Foundations/\allowbreak FormalSeries.lean}.
\begin{leancode}
def FormalDist₂ (𝒲 : Type*) : Type _ := ℤ → ℤ → 𝒲

@[inherit_doc] notation:max 𝒲 "⟦z₁^±,z₂^±⟧" => FormalDist₂ 𝒲

def shiftZ₁ : FormalDist₂ 𝒲 →ₗ[ℂ] FormalDist₂ 𝒲 where
  toFun f := fun m n => f (m - 1) n

def shiftZ₂ : FormalDist₂ 𝒲 →ₗ[ℂ] FormalDist₂ 𝒲 where
  toFun f := fun m n => f m (n - 1)

def mulBySub : FormalDist₂ 𝒲 →ₗ[ℂ] FormalDist₂ 𝒲 := shiftZ₁ - shiftZ₂

def mulBySubPow : ℕ → (FormalDist₂ 𝒲 →ₗ[ℂ] FormalDist₂ 𝒲)
  | 0 => LinearMap.id
  | (N + 1) => mulBySub.comp (mulBySubPow N)
\end{leancode}

\begin{dictionary}
\lean{ℤ → ℤ → 𝒲} is the curried form of $\Z\times\Z\to\loneW$: \lean{f m n} is $f_{m,n}$, the coefficient of
$z_1^mz_2^n$. As for one variable, \lean{FormalDist₂} is an opaque \lean{def} with only the additive and module
structure (instances in the file, analogous to Definition~\ref{l01:def:formaldist}). \lean{shiftZ₁ - shiftZ₂}
is the difference of two linear maps, and the file proves
\lean{mulBySub_apply : mulBySub f m n = f (m - 1) n - f m (n - 1)}. \lean{mulBySubPow} is defined by
recursion on $N$: \lean{mulBySubPow (N + 1) = mulBySub ∘ mulBySubPow N}. Same as Gannon p.~314. Note
that multiplication by a polynomial in two \emph{different} variables is always defined coefficientwise
by a finite formula; no truncation is needed.
\end{dictionary}

\begin{physics}
Two-variable distributions are operator products $a(z_1)b(z_2)$ and four-point-like correlators
expanded in two insertion points. Multiplying by $(z_1-z_2)^N$ is multiplying by the function that
vanishes to order $N$ where the two operators collide; it is the operation that removes the OPE
singularity of order $\le N$.
\end{physics}

\begin{status}
Definitions; no hypotheses. \lean{mulBySubPow_one : mulBySubPow 1 = mulBySub} is proved in the file.
\end{status}

\begin{lemma}[The two factorisations of $(z_1-z_2)^{N+1}$]\label{l01:lem:subpowcomm}
For every $N$ and $f$, $(z_1-z_2)^{N+1}f = (z_1-z_2)^N\big((z_1-z_2)f\big)$.
\end{lemma}

\noindent File \leanpath{VOA/\allowbreak Foundations/\allowbreak Locality.lean}.
\begin{leancode}
theorem mulBySubPow_succ' (N : ℕ) (f : FormalDist₂ 𝒲) :
    mulBySubPow (N + 1) f = mulBySubPow N (mulBySub f) := by
\end{leancode}

\begin{dictionary}
\lean{mulBySubPow} was defined by \emph{post}-composition, $(z_1-z_2)^{N+1} = (z_1-z_2)\circ(z_1-z_2)^N$; the
lemma gives the \emph{pre}-composition form. It is the associativity of iterated composition, needed to
peel factors off on either side.
\end{dictionary}

\begin{physics}
Technical: one may strip the factors $(z_1-z_2)$ one at a time from either end when cancelling an OPE
pole.
\end{physics}

\textbf{Proof on paper.} Induction on $N$. \qed

\begin{status}
Unconditional.
\end{status}

\subsection{Derivatives of the delta function}

From now on
\begin{equation*}
  \lonedelta := \sum_{k\in\Z}z_1^{k}z_2^{-k-1} = z_2^{-1}\delta(z_1/z_2),
\end{equation*}
the delta function in Kac's normalisation ($\delta(z-w)$ in his notation). It is symmetric:
$\delta(z_1,z_2) = \delta(z_2,z_1)$ (substitute $k\mapsto -k-1$). Its key property is
\begin{equation*}
  (z_1-z_2)\lonedelta = \sum_kz_1^{k+1}z_2^{-k-1}-\sum_kz_1^{k}z_2^{-k} = 0,
\end{equation*}
the two sums being equal after $k\mapsto k-1$. Consequently, for every Laurent polynomial $g$,
$g(z_1)\lonedelta = g(z_2)\lonedelta$: the delta function ``sets $z_1 = z_2$''.

\begin{definition}[$c(z_2)\,\partial_{z_2}^j\delta(z_1,z_2)$; recovered coefficients]\label{l01:def:deltaderiv}
For $j\in\N$ and $c\in\loneW[[z^{\pm1}]]$ let $c(z_2)\,\partial_{z_2}^j\lonedelta$ be the two-variable distribution
with coefficients
\begin{equation*}
  \big(c(z_2)\,\partial_{z_2}^j\lonedelta\big)_{m,n} = \Big(\prod_{i=0}^{j-1}(-m-1-i)\Big)\,c_{m+n+1+j}.
\end{equation*}
For $f\in\loneW[[z_1^{\pm1},z_2^{\pm1}]]$ and $j\in\N$, the \emph{$j$-th recovered coefficient} is the
one-variable distribution $\Res_{z_1}\big((z_1-z_2)^jf\big)$, with $z_2^n$-coefficient
$\big((z_1-z_2)^jf\big)_{-1,n}$.
\end{definition}

\noindent File \leanpath{VOA/\allowbreak Foundations/\allowbreak Locality.lean}.
\begin{leancode}
def deltaDeriv (j : ℕ) : FormalDist 𝒲 →ₗ[ℂ] FormalDist₂ 𝒲 where
  toFun c := fun m n => (∏ i ∈ Finset.range j, (-m - 1 - (i : ℤ))) • c (m + n + 1 + j)

def cCoeff (j : ℕ) (f : FormalDist₂ 𝒲) : FormalDist 𝒲 := fun n => mulBySubPow j f (-1) n
\end{leancode}

\begin{dictionary}
\textbf{Why this formula is $c(z_2)\partial_{z_2}^j\delta(z_1,z_2)$.} Differentiating $z_2^{-k-1}$ $j$ times gives
$(-k-1)(-k-2)\cdots(-k-j)\,z_2^{-k-1-j} = \prod_{i<j}(-k-1-i)\,z_2^{-k-1-j}$. So
$\partial_{z_2}^j\lonedelta = \sum_k\prod_{i<j}(-k-1-i)\,z_1^kz_2^{-k-1-j}$, whose $z_1^m$-part has $k = m$.
Multiplying by $c(z_2) = \sum_rc_rz_2^r$, the coefficient of $z_2^n$ needs $r-m-1-j = n$, i.e.\
$r = m+n+1+j$. This is the Lean formula. The product is an \emph{integer}, acting on the vector
$c_{m+n+1+j}$ by \lean{•}. The product of $c(z_2)$ with $\partial^j\delta$ is well defined because for fixed
$(m,n)$ only one coefficient of $c$ contributes. \lean{deltaDeriv j} is packaged as a $\C$-linear map in $c$.

\textbf{Plain versus divided derivatives.} Kac and Frenkel--Ben-Zvi use divided powers
$\partial^{(j)} = \partial^j/j!$; the library uses the plain $\partial^j$. The only effect is a rescaling of the
coefficients $c^j$ by $j!$; the docstring of \lean{cCoeff} records this (``up to the $j!$ of the
plain-power convention''). \lean{cCoeff j f n} is the coefficient of $z_1^{-1}z_2^{n}$ in $(z_1-z_2)^jf$, i.e.\
$\Res_{z_1}(z_1-z_2)^jf$. For $f = [a(z_1),b(z_2)]$ this is the OPE coefficient $(a_{(j)}b)(z_2)$
(Kac (2.2.4)); Chapter~\ref{ch:L03} uses it in Dong's lemma (Theorem~\ref{l03:thm:dong}).

\textbf{Which delta.} \lean{deltaDeriv 0 c} has coefficient $c_{m+n+1}$: it is $c(z_2)\delta(z_1,z_2)$ with the
$z_2^{-1}$-normalised delta, \emph{not} with \lean{delta₂} of Definition~\ref{l01:def:delta}.
\end{dictionary}

\begin{physics}
$c(w)\,\partial_w^j\delta(z,w)$ is a contact term of order $j$: paired with a test function $g(z)$ it gives
$\Res_z g(z)c(w)\partial_w^j\delta(z,w) = c(w)\,\partial_w^jg(w)$, a derivative of the test function at the
coincident point. These are precisely the terms that the singular part $c(w)\,j!/(z-w)^{j+1}$ of an OPE
contributes to a commutator (Theorem~\ref{l01:thm:alljump}). The recovered coefficient
$\Res_z(z-w)^j[a(z),b(w)]$ is the contour integral $\oint_w (z-w)^ja(z)b(w)\,dz/2\pi i$ around $w$, which
extracts the coefficient of the pole of order $j+1$: the $j$-th product $a_{(j)}b$.
\end{physics}

\begin{status}
Definitions; linearity of \lean{deltaDeriv j} proved in the file. No hypotheses.
\end{status}

\begin{example}[The first two derivatives]\label{l01:ex:deltaderiv}
Let $\mathbf 1_0\in\C[[z^{\pm1}]]$ be the constant series $1$ (coefficient $[t = 0]$; the library calls it
\lean{deltaCoeff}, Definition~\ref{l01:def:negpowjump}). Then \lean{deltaDeriv 0 deltaCoeff} has coefficient
$[m+n+1 = 0]$, which is $\lonedelta$. \lean{deltaDeriv 1 deltaCoeff} has coefficient $(-m-1)[m+n+2 = 0]$, i.e.
\begin{equation*}
  \partial_{z_2}\lonedelta = \sum_{m\in\Z}(-m-1)\,z_1^{m}z_2^{-m-2} = \sum_{k\in\Z}(k+1)z_1^{-k-2}z_2^{k}.
\end{equation*}
Check against the definition: $\partial_{z_2}z_2^{-m-1} = (-m-1)z_2^{-m-2}$. The coefficient
$(-m-1)$ vanishes at $m = -1$: $\Res_{z_1}\partial_{z_2}\lonedelta = 0$, whereas $\Res_{z_1}\lonedelta = 1$.
Similarly \lean{deltaDeriv 2 deltaCoeff} has coefficient $(m+1)(m+2)[m+n+3 = 0]$.
\end{example}

\begin{proposition}[Peeling identities]\label{l01:prop:peel}
For $j,N\in\N$ and $c\in\loneW[[z^{\pm1}]]$, writing $\partial = \partial_{z_2}$ and $\delta = \lonedelta$:
\begin{enumerate}[label=(\alph*)]
\item $(z_1-z_2)\,c\,\delta = 0$;
\item $(z_1-z_2)\,c\,\partial^{j+1}\delta = (j+1)\,c\,\partial^j\delta$;
\item $(z_1-z_2)^N\,c\,\partial^N\delta = N!\,c\,\delta$;
\item $(z_1-z_2)^N\,c\,\partial^j\delta = 0$ whenever $j+1\le N$.
\end{enumerate}
The engine is the finite-difference identity for falling factorials:
$\prod_{i\le j}(y+1-i)-\prod_{i\le j}(y-i) = (j+1)\prod_{i<j}(y-i)$ for $y\in\Z$.
\end{proposition}

\noindent File \leanpath{VOA/\allowbreak Foundations/\allowbreak Locality.lean}.
\begin{leancode}
theorem fallingDiff (y : ℤ) (j : ℕ) :
    (∏ i ∈ Finset.range (j + 1), (y + 1 - (i : ℤ))) - ∏ i ∈ Finset.range (j + 1), (y - (i : ℤ))
      = ((j : ℤ) + 1) * ∏ i ∈ Finset.range j, (y - (i : ℤ)) := by

theorem mulBySub_deltaDeriv (j : ℕ) (c : FormalDist 𝒲) :
    mulBySub (deltaDeriv (j + 1) c) = ((j : ℤ) + 1) • deltaDeriv j c := by

theorem mulBySub_deltaDeriv_zero (c : FormalDist 𝒲) : mulBySub (deltaDeriv 0 c) = 0 := by

theorem mulBySubPow_deltaDeriv_self (N : ℕ) (c : FormalDist 𝒲) :
    mulBySubPow N (deltaDeriv N c) = (N.factorial : ℤ) • deltaDeriv 0 c := by

theorem mulBySubPow_succ_deltaDeriv (j : ℕ) (c : FormalDist 𝒲) :
    mulBySubPow (j + 1) (deltaDeriv j c) = 0 := by

theorem mulBySubPow_deltaDeriv_eq_zero {j N : ℕ} (h : j + 1 ≤ N) (c : FormalDist 𝒲) :
    mulBySubPow N (deltaDeriv j c) = 0 := by
\end{leancode}

\begin{dictionary}
Statements (a)--(d) are, in order, \lean{mulBySub_deltaDeriv_zero}, \lean{mulBySub_deltaDeriv},
\lean{mulBySubPow_deltaDeriv_self} and \lean{mulBySubPow_deltaDeriv_eq_zero} (with the special case $N = j+1$
named \lean{mulBySubPow_succ_deltaDeriv}). The scalars \lean{((j : ℤ) + 1)} and \lean{(N.factorial : ℤ)} act by
integer scalar multiplication. The hypothesis \lean{h : j + 1 ≤ N} in (d) is the natural one. With divided
powers, (b) would read $(z_1-z_2)\partial^{(j+1)}\delta = \partial^{(j)}\delta$ (Kac §2.1); the library's (b) is
that identity multiplied by $(j+1)!$.
\end{dictionary}

\begin{physics}
Multiplying a contact term of order $j+1$ by $(z_1-z_2)$ lowers its order by one; $N$ such factors kill
every contact term of order $<N$. This is the formal version of the statement that $(z-w)^N$ cancels an
OPE pole of order $\le N$.
\end{physics}

\textbf{Proof on paper.} Coefficientwise,
$((z_1-z_2)F)_{m,n} = F_{m-1,n}-F_{m,n-1}$. For $F = c\,\partial^{j+1}\delta$ both terms involve the same
coefficient of $c$, namely $c_{m+n+1+j}$, because $(m-1)+n+1+(j+1) = m+(n-1)+1+(j+1) = m+n+1+j$; they are multiplied by the falling products at $y = -m-1$ shifted by one
and not shifted; \lean{fallingDiff} with $y = -m-1$ turns their difference into $(j+1)\prod_{i<j}(-m-1-i)$,
which is the coefficient of $(j+1)c\,\partial^j\delta$. For (a) the two terms are equal. (c) is (b) iterated
$N$ times, and (d) follows from (c) and (a). \qed

\begin{status}
All unconditional; axioms \lean{[propext, Classical.choice, Quot.sound]} (checked for (b), (c), (d)).
\end{status}

\begin{proposition}[The case $N = 1$]\label{l01:prop:Neq1}
If $(z_1-z_2)f = 0$, then $f = c(z_2)\lonedelta$ with $c = \Res_{z_1}f$, i.e.\ $c_n = f_{-1,n}$.
\end{proposition}

\noindent File \leanpath{VOA/\allowbreak Foundations/\allowbreak Locality.lean}.
\begin{leancode}
theorem eq_deltaDeriv_zero_of_mulBySub_eq_zero {f : FormalDist₂ 𝒲} (hf : mulBySub f = 0) :
    f = deltaDeriv 0 (fun n => f (-1) n) := by
\end{leancode}

\begin{dictionary}
\lean{fun n => f (-1) n} is the one-variable distribution $n\mapsto f_{-1,n}$, i.e.\ $\Res_{z_1}f$. The
conclusion \lean{deltaDeriv 0 (...)} has coefficient $f_{-1,m+n+1}$ at $(m,n)$. So the theorem says
$f_{m,n} = f_{-1,m+n+1}$ for all $m,n$: $f$ is constant along each antidiagonal $m+n = \text{const}$. The
hypothesis is the only one, and it is used. Same as the $N = 1$ case of Gannon Prop.~5.1.2.
\end{dictionary}

\begin{physics}
A distribution killed by $(z_1-z_2)$ is supported at $z_1 = z_2$ with no derivatives: a pure contact
term, like the commutator $[\psi(z),\psi^*(w)] = \delta(z,w)$ of a free fermion pair (Section~\ref{l01:sec:super}).
\end{physics}

\textbf{Proof on paper.} $(z_1-z_2)f = 0$ says $f_{m-1,n} = f_{m,n-1}$ for all $m,n$, i.e.\
$f_{m,n} = f_{m-1,n+1}$. Iterating in both directions, $f_{m,n} = f_{m-d,n+d}$ for all $d\in\Z$; with
$d = m+1$ this is $f_{m,n} = f_{-1,m+n+1}$. \qed

\begin{status}
Unconditional; axioms \lean{[propext, Classical.choice, Quot.sound]}.
\end{status}

\subsection{The structure theorem for local distributions}

\begin{theorem}[Structure theorem; Gannon Prop.~5.1.2, Kac Thm.~2.3]\label{l01:thm:deltasum}
Let $\loneW$ be a complex vector space, $f\in\loneW[[z_1^{\pm1},z_2^{\pm1}]]$ and $N\in\N$. Then
\begin{equation*}
  (z_1-z_2)^Nf = 0\quad\Longleftrightarrow\quad
  \exists\,c^0,\dots,c^{N-1}\in\loneW[[z^{\pm1}]]:\ f(z_1,z_2) = \sum_{j=0}^{N-1}c^j(z_2)\,\partial_{z_2}^j\lonedelta .
\end{equation*}
Moreover, in the forward direction the coefficients can be taken to be
$c^j = \frac1{j!}\Res_{z_1}(z_1-z_2)^jf$ (Corollary~\ref{l01:cor:unique}).
\end{theorem}

\noindent File \leanpath{VOA/\allowbreak Foundations/\allowbreak Locality.lean}.
\begin{leancode}
theorem exists_deltaSum_of_mulBySubPow_eq_zero (N : ℕ) (f : FormalDist₂ 𝒲)
    (hf : mulBySubPow N f = 0) :
    ∃ c : Fin N → FormalDist 𝒲, f = ∑ j : Fin N, deltaDeriv (j : ℕ) (c j) := by

theorem local_iff_deltaSum (f : FormalDist₂ 𝒲) (N : ℕ) :
    mulBySubPow N f = 0 ↔
      ∃ c : Fin N → FormalDist 𝒲, f = ∑ j : Fin N, deltaDeriv (j : ℕ) (c j) := by
\end{leancode}

\begin{dictionary}
\textbf{Decoding.} \lean{Fin N} is the type $\{0,1,\dots,N-1\}$; \lean{c : Fin N → FormalDist 𝒲} is the tuple
$(c^0,\dots,c^{N-1})$; \lean{(j : ℕ)} views $j\in$ \lean{Fin N} as a natural number;
\lean{∑ j : Fin N, deltaDeriv (j : ℕ) (c j)} is $\sum_{j<N}c^j(z_2)\partial^j_{z_2}\lonedelta$ by
Definition~\ref{l01:def:deltaderiv}. \lean{↔} is ``if and only if''.

\textbf{Quantifiers.} In Lean, arguments written before the colon of a theorem, here
\lean{(f : FormalDist₂ 𝒲) (N : ℕ)}, are universally quantified inputs: the theorem holds for \emph{every} $f$ and
every $N$ (as do the implicit ambient $\loneW$ and its vector-space structure). The tuple $c$ appears only inside the
\lean{∃} on the right-hand side of the \lean{↔}, so it is existentially quantified \emph{after} $f$ and $N$, so it may depend on $f$ and $N$, as it must. The bound $N$ is the same on
both sides: a distribution killed by $(z_1-z_2)^N$ involves derivatives of order $<N$ only, and
conversely. This is the sharp form.

\textbf{Hypotheses.} Only the ambient \lean{[AddCommGroup 𝒲] [Module ℂ 𝒲]}: $\loneW$ is a complex vector space.
The division by $j!$ in the construction is why the field must contain $\Q$; over $\C$ this is automatic.

\textbf{Same, stronger or weaker?} Same as Kac Thm.~2.3 (i)$\Leftrightarrow$(ii), up to the
plain-versus-divided-derivative rescaling of $c^j$. Gannon assumes $N\ge1$; the library allows $N = 0$, where
both sides say $f = 0$ (an empty sum). Kac's theorem also contains the formula for $c^j$ and uniqueness;
the Lean theorem states existence only, but uniqueness and the formula are an immediate consequence of
Proposition~\ref{l01:prop:peel} (Corollary~\ref{l01:cor:unique}), and the existence proof in the file constructs
exactly those $c^j$.

\textbf{Could it be vacuous?} No: both sides are inhabited by nonzero examples (e.g.\ $f = \lonedelta$, $N = 1$),
and the right side is not trivially implied by the left (for $N = 1$ it forces constancy along antidiagonals).
\end{dictionary}

\begin{physics}
This is the theorem behind the slogan ``the commutator of local fields is a finite sum of contact
terms''. Applied to $f = [a(z),b(w)]$ for mutually local $a,b$ it gives
$[a(z),b(w)] = \sum_{j<N}(a_{(j)}b)(w)\,\partial_w^{(j)}\delta(z,w)$ with $a_{(j)}b = \Res_z(z-w)^j[a(z),b(w)]$:
the commutator is determined by the singular part of the OPE, $a(z)b(w)\sim\sum_j(a_{(j)}b)(w)/(z-w)^{j+1}$.
Taking $\Res_z z^m$ and the $w^{-n-1}$ coefficient yields the commutator formula of modes,
$[a_{(m)},b_{(n)}] = \sum_{j\ge0}\binom mj(a_{(j)}b)_{(m+n-j)}$. For the stress tensor this is the
Virasoro algebra; for currents, the affine Kac--Moody algebra. The theorem is pure formal calculus.
\end{physics}

\textbf{Proof on paper.} ($\Leftarrow$) By Proposition~\ref{l01:prop:peel}(d), $(z_1-z_2)^N$ kills each
summand $c^j\partial^j\delta$ with $j<N$, and it is linear.

($\Rightarrow$) Induction on $N$. For $N = 0$, $f = (z_1-z_2)^0f = 0$ is the empty sum. Suppose the claim
holds for $N = k$ and let $(z_1-z_2)^{k+1}f = 0$. Put $g := (z_1-z_2)^kf$. Then $(z_1-z_2)g = 0$, so by
Proposition~\ref{l01:prop:Neq1}, $g = d(z_2)\delta$ with $d = \Res_{z_1}g$. Let
$c^k := \frac1{k!}d$. By Proposition~\ref{l01:prop:peel}(c),
$(z_1-z_2)^k\big(c^k\partial^k\delta\big) = k!\,c^k\delta = d\,\delta = g = (z_1-z_2)^kf$. Hence
$f' := f-c^k\partial^k\delta$ satisfies $(z_1-z_2)^kf' = 0$, and by induction
$f' = \sum_{j<k}c^j\partial^j\delta$. Then $f = \sum_{j\le k}c^j\partial^j\delta$. \qed

\begin{corollary}[Uniqueness and recovery of the coefficients]\label{l01:cor:unique}
If $f = \sum_{j<N}c^j(z_2)\partial^j_{z_2}\lonedelta$, then $c^j = \frac1{j!}\Res_{z_1}\big((z_1-z_2)^jf\big)$ for
each $j<N$; in particular the $c^j$ are unique.
\end{corollary}

\begin{dictionary}
This corollary is \emph{not} a separate Lean theorem; we state it because Kac's theorem includes it and a
referee should know it is a two-line consequence of results that \emph{are} in Lean. Its ingredients are
Proposition~\ref{l01:prop:peel} and the vanishing of $\Res_{z_1}$ on $c\,\partial^{r}\delta$ for $r\ge1$, which is
(for scalar-valued $c$; the one-line argument is the same for vector-valued $c$) the Lean lemma
\lean{deltaDeriv_succ_residue_zero} in \leanpath{VOA/\allowbreak Foundations/\allowbreak DeltaJump.lean}. The
recovered coefficient $\Res_{z_1}(z_1-z_2)^jf$ is \lean{cCoeff j f} (Definition~\ref{l01:def:deltaderiv}).
\end{dictionary}

\begin{physics}
The OPE coefficients of two local fields are uniquely determined by their commutator: the singular
part of $a(z)b(w)$ can be read off from $[a(z),b(w)]$ by contour integrals.
\end{physics}

\textbf{Proof on paper.} Apply $(z_1-z_2)^j$: by Proposition~\ref{l01:prop:peel}, the terms with $i<j$ die, the
term $i = j$ becomes $j!\,c^j\delta$, and the terms $i>j$ become multiples of $c^i\partial^{i-j}\delta$ with
$i-j\ge1$. Now take $\Res_{z_1}$: $\Res_{z_1}c\,\delta = c$, while $\Res_{z_1}c\,\partial^r\delta = 0$ for $r\ge1$
(the coefficient at $m = -1$ contains the factor $-m-1 = 0$). \qed

\begin{status}
Theorem~\ref{l01:thm:deltasum}: \texttt{\#print axioms VOA.local\_iff\_deltaSum} $=$
\lean{[propext, Classical.choice, Quot.sound]}; likewise for
\lean{exists_deltaSum_of_mulBySubPow_eq_zero}. Unconditional. Corollary~\ref{l01:cor:unique}: \textbf{not a Lean
theorem}; it is a paper consequence of Proposition~\ref{l01:prop:peel} and the Lean lemma
\lean{deltaDeriv_succ_residue_zero}. The existence proof of Theorem~\ref{l01:thm:deltasum} constructs exactly these
$c^j$, but the uniqueness statement itself is not formalised.
\end{status}

\begin{example}[The free boson]\label{l01:ex:boson}
Let $\alpha(z) = \sum_n\alpha_{(n)}z^{-n-1}$ be a field with $[\alpha_{(p)},\alpha_{(q)}] = p\,\delta_{p+q,0}$ (the
Heisenberg algebra; Chapter~\ref{ch:03}). In the library's indexing the coefficient of $z_1^mz_2^n$ in
$[\alpha(z_1),\alpha(z_2)]$ is $[\alpha_{(-m-1)},\alpha_{(-n-1)}] = (-m-1)\,[m+n+2 = 0]$. By
Example~\ref{l01:ex:deltaderiv} this is the coefficient of $\partial_{z_2}\lonedelta$:
\begin{equation*}
  [\alpha(z_1),\alpha(z_2)] = \partial_{z_2}\lonedelta,
\end{equation*}
which is \lean{deltaDeriv 1 deltaCoeff}. Hence $(z_1-z_2)^2[\alpha(z_1),\alpha(z_2)] = 0$ but
$(z_1-z_2)[\alpha(z_1),\alpha(z_2)] = \lonedelta\neq0$: the free boson is local of order exactly $2$, matching the
double pole of $\alpha(z)\alpha(w)\sim1/(z-w)^2$. Running the proof of Theorem~\ref{l01:thm:deltasum} with
$N = 2$: $g = (z_1-z_2)f = \delta$, $d = \Res_{z_1}\delta = \mathbf 1_0$, $c^1 = d/1! = \mathbf 1_0$, and
$f-c^1\partial\delta = 0$, so $c^0 = 0$. In OPE language: $\alpha_{(1)}\alpha = \mathbf 1$ (the double pole) and
$\alpha_{(0)}\alpha = 0$ (no simple pole).
\end{example}

%======================================================================
\section{The two expansions and the delta-function jump}\label{l01:sec:expansions}
%======================================================================

\subsection{The expansions \texorpdfstring{$\iota_{z_1,z_2}$ and $\iota_{z_2,z_1}$}{iota12 and iota21}}

\begin{definition}[The two expansions of $(z_1-z_2)^n$]\label{l01:def:expand}
For $n\in\Z$ define two elements of $\C[[z_1^{\pm1},z_2^{\pm1}]]$:
\begin{align*}
  \loneiota{z_1}{z_2}(z_1-z_2)^n &:= \sum_{k\ge0}(-1)^k\binom nk\,z_1^{n-k}z_2^{k}
  &&\text{(nonnegative powers of $z_2$; ``$|z_1|>|z_2|$'')},\\
  \loneiota{z_2}{z_1}(z_1-z_2)^n &:= \sum_{k\ge0}(-1)^{n+k}\binom nk\,z_1^{k}z_2^{n-k}
  &&\text{(nonnegative powers of $z_1$; ``$|z_2|>|z_1|$'')}.
\end{align*}
Also let $\lonedelta = \sum_kz_1^kz_2^{-k-1}$ be the delta function, with coefficient $[m+n+1 = 0]$.
\end{definition}

\noindent File \leanpath{VOA/\allowbreak Foundations/\allowbreak Expansion.lean}.
\begin{leancode}
noncomputable def expandPos (n : ℤ) : FormalDist₂ ℂ :=
  fun a b => if 0 ≤ b ∧ a = n - b then negOnePow b * gBinom n b.toNat else 0

noncomputable def expandNeg (n : ℤ) : FormalDist₂ ℂ :=
  fun a b => if 0 ≤ a ∧ b = n - a then negOnePow (n + a) * gBinom n a.toNat else 0

def deltaForm : FormalDist₂ ℂ := fun m n => if m + n + 1 = 0 then 1 else 0
\end{leancode}

\begin{dictionary}
\textbf{Decoding \lean{expandPos}.} The coefficient of $z_1^az_2^b$ is nonzero only if $b\ge0$ and $a = n-b$; then
it is $(-1)^b\binom nb$. With $k = b$ this is the first formula: $\loneiota{z_1}{z_2}(z_1-z_2)^n = z_1^n(1-z_2/z_1)^n$
expanded by the binomial convention of Definition~\ref{l01:def:gbinom}(iv) with second variable $-z_2$.
\textbf{Decoding \lean{expandNeg}.} Nonzero only if $a\ge0$ and $b = n-a$, with value $(-1)^{n+a}\binom na$. This is
$(z_1-z_2)^n = (-1)^n(z_2-z_1)^n = (-1)^n\sum_a\binom na z_2^{n-a}(-z_1)^a$, the second formula with $k = a$.
\textbf{Names.} The library's docstrings write $\iota_{12}$ for \lean{expandPos} and $\iota_{21}$ for \lean{expandNeg},
consistent with Kac's $i_{z,w}$ (expansion in the domain $|z|>|w|$) and FLM's convention of expanding
$(z_1-z_2)^n$ in nonnegative powers of the second variable. \lean{deltaForm} is $\lonedelta$ (the same series
as \lean{deltaDeriv 0 deltaCoeff}; see Example~\ref{l01:ex:deltaderiv}).

\textbf{Fidelity.} Same as Kac §2.1, FLM §8.1 and Lepowsky--Li §2.1. For $n\ge0$ the two series
are both the polynomial $(z_1-z_2)^n$ (all binomials with $k>n$ vanish); they differ only for $n<0$.
\end{dictionary}

\begin{physics}
These are the two radially ordered expansions of the propagator-type function $(z_1-z_2)^n$:
$\loneiota{z_1}{z_2}$ is the expansion valid when the operator at $z_1$ is further from the origin (later in radial time)
and is the one that appears in $a(z_1)b(z_2)$; $\loneiota{z_2}{z_1}$ is the one that appears in $b(z_2)a(z_1)$. A
correlation function is a single rational function; its two operator orderings correspond to its two
expansions.
\end{physics}

\begin{status}
Definitions; no hypotheses.
\end{status}

\begin{example}[Simple and double poles by hand]\label{l01:ex:poles}
For $n = -1$: $(-1)^b\binom{-1}b = (-1)^b(-1)^b = 1$, so
\begin{equation*}
  \loneiota{z_1}{z_2}(z_1-z_2)^{-1} = \sum_{b\ge0}z_1^{-1-b}z_2^{b},\qquad
  \loneiota{z_2}{z_1}(z_1-z_2)^{-1} = \sum_{a\ge0}(-1)^{a-1}(-1)^az_1^az_2^{-1-a} = -\sum_{a\ge0}z_1^az_2^{-1-a}.
\end{equation*}
For $n = -2$: $(-1)^b\binom{-2}b = b+1$ and $(-1)^{a-2}\binom{-2}a = a+1$, so
\begin{equation*}
  \loneiota{z_1}{z_2}(z_1-z_2)^{-2} = \sum_{b\ge0}(b+1)z_1^{-2-b}z_2^{b},\qquad
  \loneiota{z_2}{z_1}(z_1-z_2)^{-2} = \sum_{a\ge0}(a+1)z_1^{a}z_2^{-2-a}.
\end{equation*}
These are the Taylor series of $z_1^{-2}(1-z_2/z_1)^{-2}$ and $z_2^{-2}(1-z_1/z_2)^{-2}$. Note that the second is
\emph{not} the negative of anything: $(z_1-z_2)^{-2} = (z_2-z_1)^{-2}$ is even under exchange.
\end{example}

\begin{proposition}[Coefficient formulas and supports]\label{l01:prop:expandvals}
\leavevmode
\begin{enumerate}[label=(\alph*)]
\item $\big(\loneiota{z_1}{z_2}(z_1-z_2)^{-1}\big)_{a,b} = [b\ge0,\ a = -1-b]$ and
  $\big(\loneiota{z_2}{z_1}(z_1-z_2)^{-1}\big)_{a,b} = -[a\ge0,\ b = -1-a]$.
\item Both expansions of $(z_1-z_2)^n$ are supported on the antidiagonal $a+b = n$; the first vanishes for
  $b<0$, the second for $a<0$.
\end{enumerate}
\end{proposition}

\noindent Files \leanpath{VOA/\allowbreak Foundations/\allowbreak Expansion.lean} and
\leanpath{VOA/\allowbreak Foundations/\allowbreak IotaZExpansion.lean}.
\begin{leancode}
theorem expandPos_neg_one (a b : ℤ) :
    expandPos (-1) a b = if 0 ≤ b ∧ a = -1 - b then 1 else 0 := by

theorem expandNeg_neg_one (a b : ℤ) :
    expandNeg (-1) a b = if 0 ≤ a ∧ b = -1 - a then -1 else 0 := by

theorem expandPos_eq_zero_off_diag (n a b : ℤ) (h : a + b ≠ n) : expandPos n a b = 0 := by

theorem expandNeg_eq_zero_off_diag (n a b : ℤ) (h : a + b ≠ n) : expandNeg n a b = 0 := by

theorem expandPos_eq_zero_of_neg (n a : ℤ) {b : ℤ} (hb : b < 0) : expandPos n a b = 0 := by

theorem expandNeg_eq_zero_of_neg (n b : ℤ) {a : ℤ} (ha : a < 0) : expandNeg n a b = 0 := by
\end{leancode}

\begin{dictionary}
Exactly (a) and (b); the hypotheses \lean{h}, \lean{hb}, \lean{ha} are the conditions ``off the antidiagonal'',
``$b<0$'', ``$a<0$''.
\end{dictionary}

\begin{physics}
Homogeneity: $(z_1-z_2)^n$ has total degree $n$, so it lives on the antidiagonal $a+b = n$; in CFT, this is
scale covariance of a two-point function. The one-sided supports are the two radial orderings.
\end{physics}

\textbf{Proof on paper.} Read off from the definitions and Example~\ref{l01:ex:poles}. \qed

\begin{status}
Unconditional; \texttt{\#print axioms} clean for the jump theorem that uses them (below).
\end{status}

\subsection{The jump of the simple pole}

\begin{theorem}[The delta-function jump; Gannon §5.1.2, FLM §8.10]\label{l01:thm:jump}
\begin{equation*}
  \loneiota{z_1}{z_2}\frac1{z_1-z_2}-\loneiota{z_2}{z_1}\frac1{z_1-z_2} = \lonedelta,
  \qquad (z_1-z_2)\,\lonedelta = 0 .
\end{equation*}
\end{theorem}

\noindent File \leanpath{VOA/\allowbreak Foundations/\allowbreak Expansion.lean}.
\begin{leancode}
theorem expand_jump : expandPos (-1) - expandNeg (-1) = deltaForm := by

theorem mulBySub_deltaForm : mulBySub (deltaForm : FormalDist₂ ℂ) = 0 := by
\end{leancode}

\begin{dictionary}
The difference is the coefficientwise difference of two-variable distributions. No hypotheses. This is the
textbook identity (Kac §2.1; FLM §8.1; Lepowsky--Li §2.3, where it is a special case of the three-term
delta-function identity behind the Jacobi identity).
\end{dictionary}

\begin{physics}
This is the identity behind every equal-time commutator in chiral CFT. If $a(z)b(w)\sim c(w)/(z-w)$, then
$a(z)b(w) = c(w)\loneiota zw(z-w)^{-1}+\cdots$ and $b(w)a(z) = c(w)\loneiota wz(z-w)^{-1}+\cdots$, so the commutator
is $c(w)\delta(z,w)$: the contact term. On the circle it is the $\delta(\sigma_1-\sigma_2)$ of canonical commutation
relations. In contour language: $\oint_{|z|>|w|}-\oint_{|z|<|w|} = \oint_{w}$, the integral around the pole.
\end{physics}

\textbf{Proof on paper.} By Proposition~\ref{l01:prop:expandvals}(a), at $(m,n)$ the left side is
$[n\ge0, m = -1-n]+[m\ge0, n = -1-m]$. If $m+n+1\ne0$ both brackets vanish. If $m+n+1 = 0$, exactly one of
$n\ge0$, $m\ge0$ holds (since $m+n = -1$), so the sum is $1$. The second identity was shown at the start of
the subsection on derivatives of the delta function. \qed

\begin{status}
Unconditional; \texttt{\#print axioms VOA.expand\_jump} $=$ \lean{[propext, Classical.choice, Quot.sound]};
same for \lean{mulBySub_deltaForm}.
\end{status}

\subsection{Multiplication by \texorpdfstring{$z_1-z_2$}{z1-z2} and the algebra of the expansions}

\begin{proposition}[$(z_1-z_2)$ acts as expected on both expansions]\label{l01:prop:mulexpand}
For all $n\in\Z$ and $N\in\N$, and for both $\iota = \loneiota{z_1}{z_2}$ and $\iota = \loneiota{z_2}{z_1}$:
\begin{equation*}
  (z_1-z_2)\cdot\iota(z_1-z_2)^n = \iota(z_1-z_2)^{n+1},\qquad (z_1-z_2)^N\cdot\iota(z_1-z_2)^{n} = \iota(z_1-z_2)^{n+N}.
\end{equation*}
Moreover $\iota(z_1-z_2)^0 = 1$ (the series with the single coefficient $1$ at $z_1^0z_2^0$), and in particular
$(z_1-z_2)^N\,\iota(z_1-z_2)^{-N} = 1$.
\end{proposition}

\noindent Files \leanpath{VOA/\allowbreak Foundations/\allowbreak DeltaJump.lean} and
\leanpath{VOA/\allowbreak Foundations/\allowbreak IotaZExpansion.lean}.
\begin{leancode}
theorem mulBySub_expandPos (n : ℤ) : mulBySub (expandPos n) = expandPos (n + 1) := by

theorem mulBySub_expandNeg (n : ℤ) : mulBySub (expandNeg n) = expandNeg (n + 1) := by

def unitDist : FormalDist₂ ℂ := fun a b => if a = 0 ∧ b = 0 then (1 : ℂ) else 0

theorem expandPos_zero : expandPos 0 = unitDist := by

theorem expandNeg_zero : expandNeg 0 = expandPos 0 := by

theorem mulBySubPow_expandPos (N : ℕ) (m : ℤ) :
    mulBySubPow N (expandPos m) = expandPos (m + (N : ℤ)) := by

theorem mulBySubPow_expandPos_neg (N : ℕ) :
    mulBySubPow N (expandPos (-(N : ℤ))) = expandPos 0 := by

theorem mulBySubPow_expandNeg (N : ℕ) (m : ℤ) :
    mulBySubPow N (expandNeg m) = expandNeg (m + (N : ℤ)) := by

theorem mulBySubPow_expandNeg_neg (N : ℕ) :
    mulBySubPow N (expandNeg (-(N : ℤ))) = expandNeg 0 := by
\end{leancode}

\begin{dictionary}
\lean{unitDist} is the constant $1\in\C[[z_1^{\pm1},z_2^{\pm1}]]$. \lean{(N : ℤ)} casts $N$ to an integer. The
statements are exactly the displayed ones. \textbf{Why this is not trivial.} In a ring, $(z_1-z_2)\cdot
(z_1-z_2)^n = (z_1-z_2)^{n+1}$ is a tautology. Here the left side is the operation \lean{mulBySub} applied to a
particular infinite series, and the right side is a separately defined infinite series; the identity is
Pascal's rule (Proposition~\ref{l01:prop:pascal}) coefficientwise. That it holds for \emph{both} expansions is
the formal reason $(z_1-z_2)^N$ can be applied ``before choosing the ordering''. Note the contrast with
Theorem~\ref{l01:thm:jump}: $(z_1-z_2)\cdot[\iota_{12}-\iota_{21}](z_1-z_2)^{-1} = 1-1 = 0$, consistent with
$(z_1-z_2)\delta = 0$.
\end{dictionary}

\begin{physics}
Multiplying an OPE singularity by $(z-w)^N$ lowers the order of the pole by $N$ in either radial ordering;
after enough factors both orderings give the same polynomial, which is why $(z-w)^N[a(z),b(w)] = 0$ is the
right formulation of locality. \lean{mulBySubPow_expandPos_neg} is the statement that
$(z-w)^N\cdot(z-w)^{-N} = 1$ inside a fixed ordering; the lattice vertex operators use it to strip the
prefactor $(z-w)^{\langle\alpha,\beta\rangle}$.
\end{physics}

\textbf{Proof on paper.} On the antidiagonal $a+b = n+1$ with $b = j+1\ge1$, the coefficient of
$(z_1-z_2)\iota_{12}(z_1-z_2)^n$ is $(-1)^{j+1}\binom n{j+1}-(-1)^j\binom nj = (-1)^{j+1}\big(\binom n{j+1}+\binom nj\big) =
(-1)^{j+1}\binom{n+1}{j+1}$ by Pascal; at $b = 0$ it is $\binom n0 = 1 = \binom{n+1}0$; off the antidiagonal both sides
vanish. The computation for $\iota_{21}$ is the mirror image. The power version is induction on $N$. For
$\iota(z_1-z_2)^0$: only $k = 0$ contributes since $\binom0k = 0$ for $k\ge1$. \qed

\begin{status}
All unconditional; axioms \lean{[propext, Classical.choice, Quot.sound]} (checked for
\lean{mulBySub_expandPos}, \lean{mulBySub_expandNeg}, \lean{expandPos_zero}, \lean{expandNeg_zero},
\lean{mulBySubPow_expandPos}, \lean{mulBySubPow_expandPos_neg}).
\end{status}

\begin{definition}[Products of expansions]\label{l01:def:convol}
For $m,n\in\Z$ the product of the expansions in a fixed ordering is the two-variable series with
coefficients
\begin{align*}
  \big(\iota_{12}(z_1-z_2)^m\cdot\iota_{12}(z_1-z_2)^n\big)_{a,b} &:=
  \sum_{b_1\in\Z}\big(\iota_{12}(z_1-z_2)^m\big)_{m-b_1,b_1}\big(\iota_{12}(z_1-z_2)^n\big)_{a-(m-b_1),\,b-b_1},\\
  \big(\iota_{21}(z_1-z_2)^m\cdot\iota_{21}(z_1-z_2)^n\big)_{a,b} &:=
  \sum_{a_1\in\Z}\big(\iota_{21}(z_1-z_2)^m\big)_{a_1,m-a_1}\big(\iota_{21}(z_1-z_2)^n\big)_{a-a_1,\,b-(m-a_1)} .
\end{align*}
\end{definition}

\noindent File \leanpath{VOA/\allowbreak Foundations/\allowbreak IotaZExpansion.lean}.
\begin{leancode}
noncomputable def convolPos (m n : ℤ) : FormalDist₂ ℂ :=
  fun a b => ∑ᶠ b₁ : ℤ, expandPos m (m - b₁) b₁ * expandPos n (a - (m - b₁)) (b - b₁)

noncomputable def convolNeg (m n : ℤ) : FormalDist₂ ℂ :=
  fun a b => ∑ᶠ a₁ : ℤ, expandNeg m a₁ (m - a₁) * expandNeg n (a - a₁) (b - (m - a₁))
\end{leancode}

\begin{dictionary}
\textbf{Is this the product of series?} The genuine product of two two-variable series $F,G$ is
$(FG)_{a,b} = \sum_{(a_1,b_1)\in\Z^2}F_{a_1,b_1}G_{a-a_1,b-b_1}$, when this is a finite sum. The library has no
general product operation (Definition~\ref{l01:def:formaldist}); \lean{convolPos} is that double sum
\emph{already restricted} to the antidiagonal $a_1 = m-b_1$ of the first factor. Since
$F = \iota_{12}(z_1-z_2)^m$ vanishes off this antidiagonal (Proposition~\ref{l01:prop:expandvals}(b)), the
restriction loses nothing: \lean{convolPos m n} \emph{is} the genuine product. This identification is a
one-line paper argument, not a Lean theorem; a referee who wants to see it should note that the restriction
is harmless precisely because of \lean{expandPos_eq_zero_off_diag}. Likewise for \lean{convolNeg}, restricted to
$b_1 = m-a_1$.

\textbf{Junk value?} The summand is nonzero only for $0\le b_1\le b$ (the first factor needs $b_1\ge0$, the second
$b-b_1\ge0$), so the \lean{finsum} is a genuine finite sum for every $(a,b)$: no junk value occurs. (The proof of
Theorem~\ref{l01:thm:mult} uses exactly this finiteness, via \lean{cbz_convol_support_finite}.)
\end{dictionary}

\begin{physics}
The product of two singular factors in the same radial ordering, e.g.\ the prefactors
$(z-w)^{\langle\alpha,\beta\rangle}$ and $(z-w)^{\langle\alpha,\gamma\rangle}$ that arise when a lattice vertex operator
$e^{\alpha}(z)$ is moved past $e^{\beta}(w)$ and then past $e^{\gamma}(w)$.
\end{physics}

\begin{status}
Definitions. No hypotheses.
\end{status}

\begin{theorem}[Multiplicativity of the expansions]\label{l01:thm:mult}
For all $m,n\in\Z$,
\begin{equation*}
  \iota_{12}(z_1-z_2)^m\cdot\iota_{12}(z_1-z_2)^n = \iota_{12}(z_1-z_2)^{m+n},\qquad
  \iota_{21}(z_1-z_2)^m\cdot\iota_{21}(z_1-z_2)^n = \iota_{21}(z_1-z_2)^{m+n}.
\end{equation*}
\end{theorem}

\noindent File \leanpath{VOA/\allowbreak Foundations/\allowbreak IotaZExpansion.lean}.
\begin{leancode}
theorem expandPos_mul (m n : ℤ) : convolPos m n = expandPos (m + n) := by

theorem expandNeg_mul (m n : ℤ) : convolNeg m n = expandNeg (m + n) := by
\end{leancode}

\begin{dictionary}
No hypotheses: the statement holds for all integer exponents, including two negative ones. Together with
$\iota(z_1-z_2)^0 = 1$ it says that $n\mapsto\iota_{12}(z_1-z_2)^n$ is a group homomorphism from $(\Z,+)$ into the
(partial) multiplicative structure of $\C[[z_1^{\pm1},z_2^{\pm1}]]$: the expansions inside a fixed ordering form a
copy of the group of integer powers. An earlier docstring of \leanpath{IotaZExpansion.lean} listed a lemma
``\lean{expandPos_mul_zero}'' among the unit/inverse laws, although no declaration of that name existed; the
current docstring no longer mentions it. The unit laws are \lean{convolPos_zero_right}, \lean{convolPos_zero_left}
and \lean{expandPos_zero}.
\end{dictionary}

\begin{physics}
Inside a fixed radial ordering, singular factors multiply as functions. This is the formal content of the
statement that the analytic continuation of a correlation function does not depend on how its prefactors are
grouped, as long as all of them are expanded in the same region.
\end{physics}

\textbf{Proof on paper.} On the antidiagonal $a+b = m+n$ the $(a,b)$ coefficient of the left side is
$\sum_{b_1}(-1)^{b_1}C(m,b_1)(-1)^{b-b_1}C(n,b-b_1) = (-1)^b\sum_{b_1}C(m,b_1)C(n,b-b_1) = (-1)^bC(m+n,b)$ by the
generalised Vandermonde identity (Theorem~\ref{l01:thm:vandermonde}); this is the right side. Off the
antidiagonal both sides vanish. For $\iota_{21}$ the signs combine as
$(-1)^{m+a_1}(-1)^{n+a-a_1} = (-1)^{m+n+a}$. \qed

\begin{status}
Unconditional; \texttt{\#print axioms} of \lean{VOA.expandPos_mul} and \lean{VOA.expandNeg_mul} $=$
\lean{[propext, Classical.choice, Quot.sound]}. Used downstream in
\leanpath{VOA/\allowbreak Examples/\allowbreak LatticeA1.lean}.
\end{status}

\subsection{The jumps of higher poles}

\begin{definition}[The jump of $(z_1-z_2)^{-k-1}$]\label{l01:def:negpowjump}
Let $\mathbf 1_0\in\C[[z^{\pm1}]]$ be the constant series $1$, with coefficient $[t = 0]$. For $k\in\N$, the
\emph{jump} of $(z_1-z_2)^{-k-1}$ is
\begin{equation*}
  J_k := \iota_{12}(z_1-z_2)^{-k-1}-\iota_{21}(z_1-z_2)^{-k-1}.
\end{equation*}
\end{definition}

\noindent File \leanpath{VOA/\allowbreak Foundations/\allowbreak DeltaJump.lean}.
\begin{leancode}
def deltaCoeff : FormalDist ℂ := fun t => if t = 0 then 1 else 0

noncomputable def negPowJump (k : ℕ) : FormalDist₂ ℂ :=
  expandPos (-(k : ℤ) - 1) - expandNeg (-(k : ℤ) - 1)
\end{leancode}

\begin{dictionary}
\lean{deltaCoeff} is the one-variable series $1 = z^0$ (the docstring calls it ``the one-variable point mass
$[t=0]$'': as a function of the index $t$ it is a point mass; as a series it is the constant $1$).
\lean{-(k : ℤ) - 1} is $-k-1$. So \lean{negPowJump 0} is the left side of Theorem~\ref{l01:thm:jump}.
\end{dictionary}

\begin{physics}
$J_k$ is the commutator contribution of an OPE pole of order $k+1$ with unit coefficient: if
$a(z)b(w)\sim1/(z-w)^{k+1}$, then $[a(z),b(w)]\ni J_k$.
\end{physics}

\begin{status}
Definitions. No hypotheses.
\end{status}

\begin{theorem}[All-order delta jump]\label{l01:thm:alljump}
For every $k\in\N$,
\begin{equation*}
  \iota_{12}(z_1-z_2)^{-k-1}-\iota_{21}(z_1-z_2)^{-k-1} = \frac1{k!}\,\partial_{z_2}^k\lonedelta .
\end{equation*}
\end{theorem}

\noindent File \leanpath{VOA/\allowbreak Foundations/\allowbreak DeltaJump.lean}.
\begin{leancode}
theorem mulBySub_negPowJump (k : ℕ) :
    mulBySub (negPowJump (k + 1)) = negPowJump k := by

theorem negPowJump_residue_zero (k : ℕ) (n : ℤ) :
    negPowJump (k + 1) (-1) n = 0 := by

theorem deltaDeriv_succ_residue_zero (k : ℕ) (c : FormalDist ℂ) (n : ℤ) :
    deltaDeriv (k + 1) c (-1) n = 0 := by

theorem negPowJump_eq_deltaDeriv (k : ℕ) :
    negPowJump k = (((k.factorial : ℂ))⁻¹) • deltaDeriv k deltaCoeff := by
\end{leancode}

\begin{dictionary}
\lean{(((k.factorial : ℂ))⁻¹) •} is multiplication by the complex number $1/k!$; \lean{deltaDeriv k deltaCoeff} is
$1\cdot\partial_{z_2}^k\lonedelta$ (Definition~\ref{l01:def:deltaderiv}). So the main statement is the display with
the divided derivative $\partial^{(k)} = \partial^k/k!$: exactly the formula of Kac §2.1 and FLM §8.1. The three auxiliary
lemmas are: $(z_1-z_2)J_{k+1} = J_k$; $\Res_{z_1}J_{k+1} = 0$ (the $m = -1$ slice); and $\Res_{z_1}$ of
$c\,\partial^{k+1}\lonedelta$ vanishes. No hypotheses.
\end{dictionary}

\begin{physics}
This is the formula that converts an OPE into commutators at every order:
$a(z)b(w)\sim\sum_j\frac{(a_{(j)}b)(w)}{(z-w)^{j+1}}$ implies
$[a(z),b(w)] = \sum_j(a_{(j)}b)(w)\,\partial_w^{(j)}\delta(z,w)$, because each pole contributes its jump. Together with
Theorem~\ref{l01:thm:deltasum} it shows that ``local'' and ``has an OPE with finitely many poles'' are the same
thing.
\end{physics}

\textbf{Proof on paper.} Induction on $k$; $k = 0$ is Theorem~\ref{l01:thm:jump}. Suppose
$J_j = \frac1{j!}\partial^j\delta$. By Proposition~\ref{l01:prop:mulexpand}, $(z_1-z_2)J_{j+1} = J_j$. By
Proposition~\ref{l01:prop:peel}(b), $(z_1-z_2)\frac1{(j+1)!}\partial^{j+1}\delta = \frac{j+1}{(j+1)!}\partial^j\delta = J_j$.
So $D := J_{j+1}-\frac1{(j+1)!}\partial^{j+1}\delta$ satisfies $(z_1-z_2)D = 0$, hence
$D = (\Res_{z_1}D)(z_2)\,\delta$ (Proposition~\ref{l01:prop:Neq1}). But $\Res_{z_1}D = 0$: the slice $m = -1$ of
$\iota_{12}(z_1-z_2)^{-j-2}$ needs $-1 = -j-2-b$ with $b\ge0$, impossible for $j\ge0$; that of
$\iota_{21}(z_1-z_2)^{-j-2}$ needs $m = -1\ge0$, impossible; and the slice of $\partial^{j+1}\delta$ carries the factor
$-m-1 = 0$. Hence $D = 0$. \qed

\begin{example}[The double pole]\label{l01:ex:doublejump}
From Example~\ref{l01:ex:poles}, $J_1 = \sum_{b\ge0}(b+1)z_1^{-2-b}z_2^{b}-\sum_{a\ge0}(a+1)z_1^{a}z_2^{-2-a}$. Write the
second sum with $b = -2-a\le-2$: its terms are $(-b-1)z_1^{-2-b}z_2^{b}$. Hence
$J_1 = \sum_{b\in\Z}(b+1)z_1^{-2-b}z_2^{b}$ (the term $b = -1$ is zero), which is $\partial_{z_2}\lonedelta$ by
Example~\ref{l01:ex:deltaderiv}. This is Theorem~\ref{l01:thm:alljump} for $k = 1$, and it is the free-boson
commutator of Example~\ref{l01:ex:boson}: the double pole $1/(z-w)^2$ of $\alpha(z)\alpha(w)$ produces $\partial_w\delta(z,w)$.
\end{example}

\begin{status}
Unconditional; \texttt{\#print axioms VOA.negPowJump\_eq\_deltaDeriv} $=$
\lean{[propext, Classical.choice, Quot.sound]}.
\end{status}

\begin{corollary}[The jumps are local]\label{l01:cor:jumplocal}
For $k,N\in\N$ with $k+1\le N$, $(z_1-z_2)^NJ_k = 0$; and $J_k$ is a sum
$\sum_{j\le k}c^j(z_2)\,\partial^j_{z_2}\lonedelta$ in the sense of Theorem~\ref{l01:thm:deltasum}.
\end{corollary}

\noindent File \leanpath{VOA/\allowbreak Foundations/\allowbreak IotaZExpansion.lean}.
\begin{leancode}
theorem mulBySubPow_negPowJump_self (k : ℕ) :
    mulBySubPow (k + 1) (negPowJump k) = 0 := by

theorem mulBySubPow_negPowJump_eq_zero {k N : ℕ} (h : k + 1 ≤ N) :
    mulBySubPow N (negPowJump k) = 0 := by

theorem negPowJump_deltaSum (k : ℕ) :
    ∃ c : Fin (k + 1) → FormalDist ℂ,
      negPowJump k = ∑ j : Fin (k + 1), deltaDeriv (j : ℕ) (c j) :=
\end{leancode}

\begin{dictionary}
Hypothesis \lean{h : k + 1 ≤ N} is the order condition. \lean{negPowJump_deltaSum} is an existence statement; by
Theorem~\ref{l01:thm:alljump} one may take $c^k = \mathbf 1_0/k!$ and $c^j = 0$ for $j<k$. Its proof is the one-line
term \lean{(local_iff_deltaSum (negPowJump k) (k + 1)).mp (mulBySubPow_negPowJump_self k)}: apply the structure
theorem to the first statement.
\end{dictionary}

\begin{physics}
A pole of order $k+1$ is cancelled by $(z-w)^{k+1}$ in the commutator. For lattice vertex operators
$e^\alpha(z)e^\beta(w)\propto(z-w)^{\langle\alpha,\beta\rangle}$, so the (suitably signed) commutator of $e^\alpha$ and $e^\beta$ is killed by $(z-w)^N$
once $N\ge-\langle\alpha,\beta\rangle$; the sign cocycles that make this the right bracket are discussed in
Chapter~\ref{ch:03}.
\end{physics}

\textbf{Proof on paper.} Theorem~\ref{l01:thm:alljump} and Proposition~\ref{l01:prop:peel}(d); then
Theorem~\ref{l01:thm:deltasum}. \qed

\begin{status}
Unconditional; axioms clean (checked for \lean{mulBySubPow_negPowJump_eq_zero} and
\lean{negPowJump_deltaSum}).
\end{status}

%======================================================================
\section{Operator-valued distributions: commutators, locality, normal ordering}\label{l01:sec:operators}
%======================================================================

From now on $\loneV$ is a complex vector space and $\End(\loneV)$ the algebra of its linear endomorphisms.
Distributions with values in $\End(\loneV)$ are the candidates for fields.

\begin{definition}[Commutator distribution, mutual locality, normal ordering]\label{l01:def:bracket}
Let $a,b\in\End(\loneV)[[z^{\pm1}]]$, $a(z) = \sum_ma_mz^m$, $b(z) = \sum_nb_nz^n$.
\begin{enumerate}[label=(\roman*)]
\item The \emph{commutator} is $[a(z_1),b(z_2)] := \sum_{m,n}(a_mb_n-b_na_m)\,z_1^mz_2^n$.
\item $a$ and $b$ are \emph{mutually local} if $(z_1-z_2)^N[a(z_1),b(z_2)] = 0$ for some $N\in\N$.
\item The two-variable \emph{normal-ordered product} is
  ${:}a(z_1)b(z_2){:} := a(z_1)_+b(z_2)+b(z_2)a(z_1)_-$, where $a(z)_+ := \sum_{m\ge0}a_mz^m$ and
  $a(z)_- := \sum_{m<0}a_mz^m$; coefficientwise, ${:}ab{:}_{m,n} = a_mb_n$ if $m\ge0$ and $b_na_m$ if $m<0$.
\end{enumerate}
\end{definition}

\noindent File \leanpath{VOA/\allowbreak Foundations/\allowbreak Locality.lean}.
\begin{leancode}
def bracketDist (a b : FormalDist (Module.End ℂ V)) : FormalDist₂ (Module.End ℂ V) :=
  fun m n => a m * b n - b n * a m

def Local (a b : FormalDist (Module.End ℂ V)) : Prop :=
  ∃ N : ℕ, mulBySubPow N (bracketDist a b) = 0

def normalOrdered (a b : FormalDist (Module.End ℂ V)) : FormalDist₂ (Module.End ℂ V) :=
  fun m n => if 0 ≤ m then a m * b n else b n * a m
\end{leancode}

\begin{dictionary}
\textbf{Decoding.} \lean{Module.End ℂ V} is $\End_\C(\loneV)$; its multiplication \lean{*} is composition, so
\lean{a m * b n} is $a_m\circ b_n$. \lean{bracketDist a b m n} is therefore the coefficient of $z_1^mz_2^n$ in
$[a(z_1),b(z_2)]$. Each coefficient is a single composition, so (i) is always defined: the two variables are
different. \lean{Local a b} is a \lean{Prop}: ``there exists a natural number $N$ (depending on $a,b$) such that
all coefficients of $(z_1-z_2)^N[a(z_1),b(z_2)]$ vanish''.

\textbf{Normal ordering.} In \lean{normalOrdered}, the coefficients $a_m$ with $m\ge0$ (the nonnegative powers
of $z_1$, i.e.\ the modes $a_{(k)}$ with $k = -m-1\le-1$: the ``creation'' part) stand to the \emph{left} of $b$,
the others to the right. This is Kac (2.3.5) and Gannon (5.1.6b). The same-variable product ${:}a(z)b(z){:}$ is
\emph{not} defined in this file, because it requires the field property; it is \lean{nprod} in
\leanpath{VOA/\allowbreak Core/\allowbreak Dong.lean} (Chapter~\ref{ch:L03}).

\textbf{Fidelity.} Same as Kac Def.~2.3 and Gannon §5.1.2. The definition is for arbitrary distributions, not
only fields, which is the textbook generality. The library proves elsewhere that locality is symmetric
(\lean{Local.symm} in \leanpath{VOA/\allowbreak Core/\allowbreak FieldOps.lean}); the definition itself is
asymmetric only in the bookkeeping. \textbf{Could \lean{Local} be vacuous?} With $N = 0$ it says $[a,b] = 0$; it
is satisfied by many nontrivial pairs (Examples~\ref{l01:ex:boson}, \ref{l01:ex:fermion}) and fails, e.g., for
$a = b = $ a distribution whose commutator has a nonzero coefficient on infinitely many antidiagonals with growing
derivative order, so it is a genuine condition. A useful corollary of Section~\ref{l01:sec:truncation}: if
$[a(z_1),b(z_2)]$ has \emph{finite support} and $a,b$ are local, then $[a(z_1),b(z_2)] = 0$.
\end{dictionary}

\begin{physics}
\lean{Local} is the chiral-algebra form of microcausality: two chiral fields commute up to contact terms at
coincident points. It is the axiom (VA4) of the library's structure \lean{PreVertexAlgebra} in Chapter~\ref{ch:L03}
(Definition~\ref{l03:def:VA}), imposed on all pairs of vertex operators $Y(u,z)$, $Y(v,z)$. The integer $N$ is
the maximal order of the OPE pole between $a$ and $b$. Normal ordering puts annihilation operators to the
right, so that ${:}a(z)b(w){:}$ is the regular part of the OPE, $a(z)b(w) = \text{(singular part)} + {:}a(z)b(w){:}$.
\end{physics}

\begin{status}
Definitions. \lean{Local} is used as the locality axiom of the library's structures \lean{PreVertexAlgebra} and
\lean{BarePreVA}, and hence of every structure built on \lean{PreVertexAlgebra} (\lean{VertexAlgebra},
\lean{GradedVertexOperatorAlgebra}, \lean{VertexOperatorAlgebra}). (The release audit of 2026-10-27 established that
\lean{PreVertexAlgebra} and \lean{BarePreVA} have no translation-covariance axiom and are therefore \emph{not} vertex algebras in
the standard sense; the Borcherds identity fails for them. The library has since renamed them accordingly; its
\lean{VertexAlgebra} is a \lean{PreVertexAlgebra} with translation covariance added as an axiom. See Chapter~\ref{ch:L03}.)
\end{status}

\begin{proposition}[Commuting distributions are local]\label{l01:prop:localzero}
If $[a(z_1),b(z_2)] = 0$, then $a$ and $b$ are mutually local (with $N = 0$).
\end{proposition}

\noindent File \leanpath{VOA/\allowbreak Foundations/\allowbreak Locality.lean}.
\begin{leancode}
theorem local_of_bracketDist_eq_zero {a b : FormalDist (Module.End ℂ V)}
    (h : bracketDist a b = 0) : Local a b :=
\end{leancode}

\begin{dictionary}
Hypothesis: the commutator distribution is zero. Conclusion: \lean{Local a b}, witnessed by $N = 0$.
\end{dictionary}

\begin{physics}
Commuting fields, e.g.\ a field and a constant multiple of the identity, or fields from decoupled sectors of
a tensor-product theory, are trivially local.
\end{physics}

\textbf{Proof on paper.} $(z_1-z_2)^0 = \mathrm{id}$. \qed

\begin{status}
Unconditional.
\end{status}

\subsection{Fields and the same-variable normal-ordered product}\label{l01:sec:field}

The two-variable normal ordering of Definition~\ref{l01:def:bracket} keeps the variables $z_1$ and $z_2$
apart, so each of its coefficients is one composition of two operators and nothing infinite is ever summed.
Physicists, however, want the product ${:}a(z)b(z){:}$ of two fields \emph{at the same point}: the coefficient of
$z^s$ in it is an infinite sum $\sum_m(\cdots)_{m,s-m}$ over the antidiagonal $m+n=s$ of operators, and such an
infinite sum of operators has no meaning unless one says how it converges. The remedy is to apply it to a vector
and to require that, on every vector, only finitely many terms survive. This is exactly what the \emph{field}
condition (lower truncation) buys.

\begin{definition}[Field; the normal-ordered product ${:}a(z)b(z){:}$]\label{l01:def:field}
A distribution $a(z) = \sum_{n\in\Z}a_nz^n\in\End(\loneV)[[z^{\pm1}]]$ is a \emph{field} if for every vector
$x\in\loneV$ there is an integer $B$ (depending on $x$) with $a_nx = 0$ for all $n<B$. Equivalently
$a(z)x\in\loneV((z))$ has only finitely many negative powers of $z$; equivalently, in the mode indexing
$a_{(m)} = a_{-m-1}$, one has $a_{(m)}x = 0$ for $m\gg0$.
For two fields $a,b$ the \emph{normal-ordered product} is the distribution ${:}a(z)b(z){:}$ whose $z^s$-coefficient
is the operator
\begin{equation*}
 ({:}ab{:})_s\,x \;=\; \sum_{m\in\Z}{:}a_mb_{s-m}{:}\,x
 \;=\; \sum_{m\ge0}a_mb_{s-m}x \;+\; \sum_{m<0}b_{s-m}a_mx ,
\end{equation*}
the sum being finite for every $x$. In vertex-algebra language it is the $(-1)$-st product of fields,
$a_{(-1)}b$.
\end{definition}

\noindent File \leanpath{VOA/\allowbreak Core/\allowbreak Dong.lean}.
\begin{leancode}
def IsField (a : FormalDist (End ℂ V)) : Prop :=
  ∀ x : V, ∃ B : ℤ, ∀ n : ℤ, n < B → a n x = 0

noncomputable def nprodCoeff {a b : FormalDist (End ℂ V)}
    (ha : IsField a) (hb : IsField b) (s : ℤ) : V →ₗ[ℂ] V where
  toFun x := ∑ᶠ m : ℤ, normalOrdered a b m (s - m) x

noncomputable def nprod {a b : FormalDist (End ℂ V)}
    (ha : IsField a) (hb : IsField b) : FormalDist (End ℂ V) :=
  fun s => nprodCoeff ha hb s
\end{leancode}
(The two linearity fields of \lean{nprodCoeff}, \lean{map_add'} and \lean{map_smul'}, are proofs and are omitted.
\lean{normalOrdered} is Definition~\ref{l01:def:bracket}, file \leanpath{VOA/\allowbreak Foundations/\allowbreak Locality.lean}.)

\begin{dictionary}
\textbf{Decoding \lean{IsField}.} \lean{a n x} is the coefficient $a_n$ of $z^n$, an endomorphism, applied to $x$.
The statement reads: for every $x$ there exists an integer $B$ such that for every integer $n<B$ the vector
$a_nx$ vanishes. ``$n$ very negative'' in the $z^n$-indexing is ``$m=-n-1$ very positive'' in the mode
indexing, which is the familiar lower truncation $u_{(m)}x = 0$ for $m\gg0$. The bound $B$ depends on $x$ but
not on anything else; it is \emph{not} uniform in $x$ (a uniform bound would force $a$ to be a Laurent
polynomial of operators). This is Kac's Definition~1.3 and the ``truncation condition'' of Lepowsky--Li.

\textbf{Decoding \lean{nprodCoeff}.} The $z^s$-coefficient applied to $x$ is
\lean{∑ᶠ m : ℤ, normalOrdered a b m (s - m) x}, where \lean{normalOrdered a b m n} is $a_mb_n$ if $m\ge0$ and
$b_na_m$ if $m<0$ (the creation part of $a$ is placed to the left, the annihilation part to the right).
The symbol \lean{∑ᶠ} is Mathlib's \lean{finsum}: the sum over the support of the summand if that support is
finite, and the value $0$ if it is infinite. \textbf{This is the trap to be checked here.} With no hypotheses
the Lean expression would silently equal $0$ whenever the support is infinite. The library therefore makes the two
field hypotheses \lean{ha}, \lean{hb} \emph{arguments} of \lean{nprodCoeff}, and proves
(\lean{normalOrdered_diag_support_finite}, same file) that under them the support is finite for every $s$ and $x$.
So the \lean{finsum} is a genuine finite sum, and \lean{nprod ha hb} is the textbook object, Kac (2.3.5)--(2.3.6)
and Gannon (5.1.6b). The other direction of fidelity: nothing has been left out, since the only data of
\lean{nprod} are the two distributions and the proof that they are fields.

\textbf{Same, stronger or weaker?} The same as the textbook. In the library the product of two fields is again a
field (\lean{IsField.nprod}, \leanpath{VOA/\allowbreak Core/\allowbreak FieldOps.lean}), and ${:}a\,b{:}$ is mutually
local with a third field whenever the three are pairwise local (Dong's lemma, Theorem~\ref{l03:thm:dong}).

\textbf{Could it be vacuous?} The zero distribution and every constant $c\cdot\mathrm{id}$ are fields, so the
notion is satisfied; and every vertex operator $Y(u,z)$ of a \lean{BarePreVA} (a bare pre-vertex algebra: weaker than a standard vertex
algebra, since it lacks translation covariance; Chapter~\ref{ch:L03}) is a field by the regularity axiom
VA5 (\lean{BarePreVA.isField_Y}, Chapter~\ref{ch:L03}). It is a genuine condition: the distribution $\sum_{n\in\Z}\mathrm{id}_V\,z^n$ (the delta function times the identity) has $a_nx = x\neq0$ for all $n$ and $x\neq0$, so it is not a field.
\end{dictionary}

\begin{physics}
A field is a distribution that truncates on every state, $a(z)x\in V((z))$: applying a local operator to a state of
finite energy creates only finitely many states of lower energy, which is the algebraic form of ``the energy is
bounded below'', and in the OPE it says that the singularity as $w\to z$ has a pole of finite order. The
normal-ordered product ${:}a(z)b(z){:}$ is the regular term in the operator product expansion of $a(z)b(w)$ as
$w\to z$, namely $\lim_{w\to z}\big(a(z)b(w)-\text{singular terms}\big)$. Free-field examples are the stress
tensor of a free boson, $T = \tfrac12{:}J(z)J(z){:}$, and the Sugawara construction. Placing creation operators to
the left of annihilation operators is the usual prescription that makes the vacuum expectation value of
${:}a\,b{:}$ vanish for free fields.
\end{physics}

\textbf{Proof on paper that the sum is finite, and that the product is a field.} Fix $x$ and $s$, and let $B_a$ be the
truncation bound of $a$ at $x$ and, for each vector $y$, $B_b(y)$ that of $b$ at $y$. In the sum over $m\ge0$ the term
$a_mb_{s-m}x$ vanishes as soon as $s-m<B_b(x)$, i.e.\ for $m>s-B_b(x)$. In the sum over $m<0$ the term
$b_{s-m}a_mx$ vanishes as soon as $m<B_a$. Hence only $m$ in the finite interval
$[\min(B_a,0),\ \max(s-B_b(x),0)]$ can contribute. The coefficient is clearly linear in $x$, since each summand is.
To see that ${:}ab{:}$ is again a field, let $s$ be very negative. For $m\ge0$ we have $s-m\le s<B_b(x)$, so
$a_mb_{s-m}x = 0$. For $m<0$ only the finitely many $m$ with $B_a\le m<0$ matter, since $a_mx=0$ for $m<B_a$;
choose $s$ below $\min_{B_a\le m<0}\big(m+B_b(a_mx)\big)$, and then $b_{s-m}a_mx = 0$ for each of them. \qed

\begin{status}
Definitions, not theorems. \texttt{\#print axioms} (run for this chapter with
\texttt{lake env lean -DmaxSynthPendingDepth=3}) of \lean{VOA.IsField}, \lean{VOA.nprodCoeff}, \lean{VOA.nprod} is
\lean{[propext, Classical.}\allowbreak\lean{choice, Quot.}\allowbreak\lean{sound]} for each. No hypotheses beyond the
two field arguments; \lean{VOA.pbw} is not used.
\end{status}

\subsection{Super-locality: the \texorpdfstring{$\varepsilon$}{epsilon}-twisted commutator}\label{l01:sec:super}

For fermionic (odd) fields, the radially ordered product $b(z_2)a(z_1)$ acquires a sign, and locality is a
statement about the \emph{anticommutator}. Rather than duplicating the theory, the library introduces a
scalar $\varepsilon$ in front of the second ordering.

\begin{definition}[$\varepsilon$-twisted commutator and $\varepsilon$-locality]\label{l01:def:bracketeps}
For $\varepsilon\in\C$ and $a,b\in\End(\loneV)[[z^{\pm1}]]$,
$[a(z_1),b(z_2)]_\varepsilon := \sum_{m,n}(a_mb_n-\varepsilon\,b_na_m)z_1^mz_2^n$, and $a,b$ are
\emph{$\varepsilon$-local} if $(z_1-z_2)^N[a(z_1),b(z_2)]_\varepsilon = 0$ for some $N\in\N$.
\end{definition}

\noindent File \leanpath{VOA/\allowbreak Foundations/\allowbreak Locality.lean}.
\begin{leancode}
def bracketDistε (ε : ℂ) (a b : FormalDist (Module.End ℂ V)) : FormalDist₂ (Module.End ℂ V) :=
  fun m n => a m * b n - ε • (b n * a m)

def Localε (ε : ℂ) (a b : FormalDist (Module.End ℂ V)) : Prop :=
  ∃ N : ℕ, mulBySubPow N (bracketDistε ε a b) = 0
\end{leancode}

\begin{dictionary}
\lean{ε • (b n * a m)} is the complex scalar $\varepsilon$ times the composition $b_n\circ a_m$. The intended
instances are $\varepsilon = 1$ (bosonic) and $\varepsilon = (-1)^{|a||b|}$ (Koszul sign of parities $|a|,|b|\in\Z/2$;
the library's \lean{koszul} lives in \leanpath{VOA/\allowbreak Foundations/\allowbreak Parity.lean}, another chapter).
For other values of $\varepsilon$ the definition is meaningful but has no standard name; nothing in this chapter
depends on $\varepsilon\in\{\pm1\}$ except \lean{bracketDistε_swap}, which needs $\varepsilon\neq0$. The file header
stresses that no spin--statistics relation between parity and half-integer weights is assumed here.
\end{dictionary}

\begin{physics}
For two odd fields (fermions, ghosts) the physically meaningful bracket is the anticommutator
$\{a(z),b(w)\} = a(z)b(w)+b(w)a(z)$, and locality says $(z-w)^N\{a(z),b(w)\} = 0$: fermions anticommute at
spacelike separation. This is the case $\varepsilon = -1$.
\end{physics}

\begin{status}
Definitions.
\end{status}

\begin{proposition}[Basic properties of the twisted commutator]\label{l01:prop:epsbasic}
\leavevmode
\begin{enumerate}[label=(\alph*)]
\item $[a,b]_1 = [a,b]$, hence $1$-locality is ordinary locality.
\item $[a,b]_\varepsilon$ is additive in $a$ and in $b$.
\item If $\varepsilon\neq0$: $[a(z_1),b(z_2)]_\varepsilon$ has $(m,n)$-coefficient
  $-\varepsilon\cdot\big([b(z_1),a(z_2)]_{\varepsilon^{-1}}\big)_{n,m}$.
\item If $[a,b]_\varepsilon = 0$ then $a,b$ are $\varepsilon$-local; every $a$ is $\varepsilon$-local with $0$ on either side.
\end{enumerate}
\end{proposition}

\noindent File \leanpath{VOA/\allowbreak Foundations/\allowbreak Locality.lean}.
\begin{leancode}
theorem bracketDistε_one (a b : FormalDist (Module.End ℂ V)) :
    bracketDistε 1 a b = bracketDist a b := by

theorem localε_one_iff (a b : FormalDist (Module.End ℂ V)) :
    Localε 1 a b ↔ Local a b := by

theorem bracketDistε_add_left (ε : ℂ) (a₁ a₂ b : FormalDist (Module.End ℂ V)) :
    bracketDistε ε (a₁ + a₂) b = bracketDistε ε a₁ b + bracketDistε ε a₂ b := by

theorem bracketDistε_add_right (ε : ℂ) (a b₁ b₂ : FormalDist (Module.End ℂ V)) :
    bracketDistε ε a (b₁ + b₂) = bracketDistε ε a b₁ + bracketDistε ε a b₂ := by

theorem bracketDistε_swap {ε : ℂ} (hε : ε ≠ 0) (a b : FormalDist (Module.End ℂ V)) (m n : ℤ) :
    bracketDistε ε a b m n = -(ε • bracketDistε ε⁻¹ b a n m) := by

theorem localε_of_bracketDistε_eq_zero {ε : ℂ} {a b : FormalDist (Module.End ℂ V)}
    (h : bracketDistε ε a b = 0) : Localε ε a b :=

theorem localε_zero_right (ε : ℂ) (a : FormalDist (Module.End ℂ V)) : Localε ε a 0 := by

theorem localε_zero_left (ε : ℂ) (b : FormalDist (Module.End ℂ V)) : Localε ε 0 b := by
\end{leancode}

\begin{dictionary}
In \lean{bracketDistε_swap} note the swapped indices \lean{n m} on the right: the coefficient of $z_1^mz_2^n$ in
$[a(z_1),b(z_2)]_\varepsilon$ is compared with the coefficient of $z_1^nz_2^m$ in $[b(z_1),a(z_2)]_{\varepsilon^{-1}}$. For
$\varepsilon = \pm1$, $\varepsilon^{-1} = \varepsilon$ and this is the graded antisymmetry
$[a,b] = -(-1)^{|a||b|}[b,a]$. The hypothesis \lean{hε} is needed to form $\varepsilon^{-1}$ meaningfully
(in Lean $0^{-1} = 0$, and the identity would be false at $\varepsilon = 0$ unless $ab = 0$).
\end{dictionary}

\begin{physics}
(a) says the bosonic theory is the $\varepsilon = 1$ instance of the super theory; (c) is graded
antisymmetry of the (anti)commutator; (d) says that the zero field is local with everything.
\end{physics}

\textbf{Proof on paper.} Coefficientwise algebra: $a_mb_n-\varepsilon b_na_m = -\varepsilon(b_na_m-\varepsilon^{-1}a_mb_n)$. \qed

\begin{status}
All unconditional; axioms clean (checked for \lean{bracketDistε_one}, \lean{localε_one_iff},
\lean{bracketDistε_swap}).
\end{status}

\begin{theorem}[Structure theorem, $\varepsilon$-generically]\label{l01:thm:epsdeltasum}
$a,b$ are $\varepsilon$-local iff for some $N$ there are $c^0,\dots,c^{N-1}\in\End(\loneV)[[z^{\pm1}]]$ with
$[a(z_1),b(z_2)]_\varepsilon = \sum_{j<N}c^j(z_2)\partial^j_{z_2}\lonedelta$. In particular, if
$[a(z_1),b(z_2)]_\varepsilon = c(z_2)\lonedelta$, then $a,b$ are $\varepsilon$-local with $N = 1$.
\end{theorem}

\noindent File \leanpath{VOA/\allowbreak Foundations/\allowbreak Locality.lean}.
\begin{leancode}
theorem localε_iff_deltaSum (ε : ℂ) (a b : FormalDist (Module.End ℂ V)) :
    Localε ε a b ↔ ∃ (N : ℕ) (c : Fin N → FormalDist (Module.End ℂ V)),
      bracketDistε ε a b = ∑ j : Fin N, deltaDeriv (j : ℕ) (c j) := by

theorem localε_of_bracketDistε_eq_delta {ε : ℂ} {a b c : FormalDist (Module.End ℂ V)}
    (h : bracketDistε ε a b = deltaDeriv 0 c) : Localε ε a b :=
\end{leancode}

\begin{dictionary}
Here $N$ is existentially quantified \emph{inside} the iff (``$\exists N\,\exists c$''), because \lean{Localε} itself
hides $N$; compare Theorem~\ref{l01:thm:deltasum}, where $N$ is fixed outside. Both versions are equivalent;
this one is the convenient packaging for fields. The structure theorem is applied to the two-variable
distribution $[a,b]_\varepsilon$; the sign never enters its proof, which is the structural reason the super
theory is ``a refactor rather than a rewrite'' (file header).
\end{dictionary}

\begin{physics}
The super-commutator of two local fields is determined by the singular part of their OPE, exactly as in the
bosonic case: $\{\psi(z),\psi^*(w)\} = \delta(z,w)$ for a $bc$ or free-fermion system, $[\beta(z),\gamma(w)] = \delta(z,w)$
for a $\beta\gamma$ system.
\end{physics}

\textbf{Proof on paper.} Theorem~\ref{l01:thm:deltasum} applied to $f = [a,b]_\varepsilon$; the second statement is
Proposition~\ref{l01:prop:peel}(a). \qed

\begin{status}
Unconditional; axioms clean (checked for \lean{localε_iff_deltaSum}).
\end{status}

\begin{theorem}[Free-field criterion: canonical (anti)commutation relations]\label{l01:thm:car}
Let $S\in\End(\loneV)$ and suppose $a_mb_n-\varepsilon\,b_na_m = [m+n+1 = 0]\,S$ for all $m,n\in\Z$. Then $a,b$ are
$\varepsilon$-local (with $N = 1$). In particular, if $a_mb_n+b_na_m = [m+n+1 = 0]\,S$ for all $m,n$, then $a,b$ are
$(-1)$-local.
\end{theorem}

\noindent File \leanpath{VOA/\allowbreak Foundations/\allowbreak Locality.lean}.
\begin{leancode}
theorem localε_of_car {ε : ℂ} {a b : FormalDist (Module.End ℂ V)} (S : Module.End ℂ V)
    (h : ∀ m n : ℤ, a m * b n - ε • (b n * a m) = if m + n + 1 = 0 then S else 0) :
    Localε ε a b := by

theorem localε_neg_one_of_anticomm {a b : FormalDist (Module.End ℂ V)}
    (S : Module.End ℂ V) (h : ∀ m n : ℤ, a m * b n + b n * a m = if m + n + 1 = 0 then S else 0) :
    Localε (-1) a b := by
\end{leancode}

\begin{dictionary}
\textbf{Indices.} $m,n$ are \emph{coefficient} indices ($z^m$), not modes. In modes, $a_m = a_{(-m-1)}$, and the
condition $m+n+1 = 0$ becomes $p+q+1 = 0$ for $p = -m-1$, $q = -n-1$: the hypothesis reads
$[a_{(p)},b_{(q)}]_\varepsilon = \delta_{p+q+1,0}\,S$, the same shape in both indexings. $S$ is an arbitrary fixed operator
(in examples, a multiple of the identity). The hypothesis says $[a(z_1),b(z_2)]_\varepsilon = S\,\lonedelta$.
\end{dictionary}

\begin{physics}
These are the canonical commutation relations of a $\beta\gamma$ system ($\varepsilon = 1$) and the canonical
anticommutation relations of a $bc$ system or complex fermion ($\varepsilon = -1$): the OPE has a single simple pole
$S/(z-w)$. The theorem is used to verify locality of the generating fields in
\leanpath{VOA/\allowbreak Examples/\allowbreak BetaGamma.lean}, \leanpath{BetaGammaField.lean} and
\leanpath{SymplecticFermion.lean}.
\end{physics}

\textbf{Proof on paper.} The hypothesis says $[a,b]_\varepsilon = c(z_2)\lonedelta$ with $c = S\cdot\mathbf 1_0$ (since
$c_{m+n+1} = [m+n+1 = 0]S$); apply Theorem~\ref{l01:thm:epsdeltasum}. \qed

\begin{example}[The free fermion]\label{l01:ex:fermion}
Let $\psi(z) = \sum_p\psi_{(p)}z^{-p-1}$ with $\{\psi_{(p)},\psi_{(q)}\} = \delta_{p+q+1,0}$ (a neutral fermion with
integer moding; the OPE is $\psi(z)\psi(w)\sim1/(z-w)$). In coefficient indexing $\psi_m = \psi_{(-m-1)}$ and
$\{\psi_m,\psi_n\} = [m+n+1 = 0]$, so \lean{localε_neg_one_of_anticomm} with $S = \mathrm{id}$ shows
$(z_1-z_2)\{\psi(z_1),\psi(z_2)\} = 0$. The bosonic commutator $[\psi(z_1),\psi(z_2)]$, with coefficients
$\psi_m\psi_n-\psi_n\psi_m = 2\psi_m\psi_n-[m+n+1=0]$, is \emph{not} killed by any power of $z_1-z_2$ in the Fock
representation: by Theorem~\ref{l01:thm:deltasum}, along each antidiagonal $m+n = s$ the coefficients of a distribution
killed by $(z_1-z_2)^N$ are polynomial in $m$ of degree $<N$ with fixed operator coefficients, hence span a
finite-dimensional space, whereas the operators $\psi_m\psi_{s-m}$ ($m>s-m$) are linearly independent in the Fock
space. (Reason: in the Clifford algebra generated by the $\psi_m$, ordered products of distinct generators form a
basis, as for an exterior algebra; and the Fock representation of this Clifford algebra is faithful, because the
algebra of a nondegenerate form on a countable-dimensional space is a union of matrix algebras and hence simple.
So distinct ordered pairs give linearly independent operators.) So the choice of $\varepsilon$ matters.
\end{example}

\begin{warning}\label{l01:warn:finitefermion}
\textbf{A docstring claim that was heuristic (corrected in the source).} An earlier comment block before
\lean{localε_of_car} said that ``a finite-support commutator can never satisfy this'' and concluded ``consequently
there is no finite-dimensional fermionic toy example''. The current comment block keeps only the finiteness remark
about mode support and says explicitly that no statement about finite-dimensional examples is proved or implied. The
first sentence of the old comment is a correct theorem: a finitely supported two-variable
distribution is lower-truncated, so by Theorem~\ref{l01:thm:cancel} it is killed by $(z_1-z_2)^N$ only if it is $0$.
The second sentence is not a consequence of the first (a finite-dimensional $\loneV$ does not force the
commutator to have finite support); it is a true but separate fact about Clifford algebras of infinite rank
with nondegenerate form, which have no nonzero finite-dimensional representations when $S$ is a nonzero scalar.
The library proves neither the second sentence nor needs it.
\end{warning}

\begin{status}
Unconditional; axioms clean (checked for both declarations).
\end{status}

%======================================================================
\section{Lower truncation and the cancellation theorem}\label{l01:sec:truncation}
%======================================================================

Theorem~\ref{l01:thm:deltasum} shows that $(z_1-z_2)$ has a huge kernel on $\loneW[[z_1^{\pm1},z_2^{\pm1}]]$: all of the
$c(z_2)\lonedelta$. This kernel disappears on the subspace of series that are bounded below in $z_2$, which is
where the products of fields $Y(u,z_1)Y(v,z_2)x$ live. This simple fact is the irreducible core of Goddard's
uniqueness theorem and of the reconstruction theorem (Chapter~\ref{ch:L03},
Theorems~\ref{l03:thm:goddard} and \ref{l03:thm:reconstruct}).

\begin{definition}[Lower truncation in $z_2$]\label{l01:def:bddbelow}
$f\in\loneW[[z_1^{\pm1},z_2^{\pm1}]]$ is \emph{lower-truncated in $z_2$} if there is $B\in\Z$ with $f_{m,n} = 0$ for all
$m\in\Z$ and all $n<B$.
\end{definition}

\noindent File \leanpath{VOA/\allowbreak Foundations/\allowbreak IteratedLaurent.lean}.
\begin{leancode}
def IsBddBelow₂ (f : FormalDist₂ 𝒲) : Prop := ∃ B : ℤ, ∀ m n : ℤ, n < B → f m n = 0

theorem IsBddBelow₂.mulBySub {f : FormalDist₂ 𝒲} (h : IsBddBelow₂ f) :
    IsBddBelow₂ (mulBySub f) := by

theorem IsBddBelow₂.mulBySubPow (N : ℕ) {f : FormalDist₂ 𝒲} (h : IsBddBelow₂ f) :
    IsBddBelow₂ (mulBySubPow N f) := by
\end{leancode}

\begin{dictionary}
\textbf{Quantifiers.} $\exists B\,\forall m\,\forall n$: the bound is \emph{uniform in the $z_1$-power}. This is the
space $\loneW[[z_1^{\pm1}]]((z_2))$ of Laurent series in $z_2$ (finitely many negative powers) with arbitrary
doubly-infinite coefficients in $z_1$. An earlier docstring called this ``the sub-class $\mathcal W((z_1))((z_2))$''; the
current docstring of \lean{IsBddBelow₂} says ``strictly larger than'' that class, which is right (the docstring of the
cancellation theorem that follows still speaks of $\mathcal W((z_1))((z_2))$). The old wording was imprecise: membership in $\loneW((z_1))((z_2))$ would also require each $z_2^n$-coefficient to be bounded below
in $z_1$, which \lean{IsBddBelow₂} does \emph{not} require. The predicate is therefore \emph{weaker} than
membership in $\loneW((z_1))((z_2))$, and theorems with it as a hypothesis are correspondingly \emph{stronger} than
their textbook counterparts. No harm results.

\textbf{Why $Y(u,z_1)Y(v,z_2)x$ satisfies it.} Its coefficient of $z_1^mz_2^n$ is $u_m(v_nx)$ (coefficient indices).
Since $Y(v,z)$ is a field, $v_nx = 0$ for $n<B$ with $B$ depending on $v,x$ only; then $u_m(v_nx) = 0$ for every $m$.
The stability lemmas say that multiplication by $(z_1-z_2)^N$ preserves the property (with the same $B$).
\end{dictionary}

\begin{physics}
Lower truncation in $z_2$ is the energy bound for the operator at the inner insertion point: acting with a
field on a state of bounded energy produces finitely many states of lower energy. It singles out one radial
ordering, $|z_1|>|z_2|$: the product $a(z_1)b(z_2)x$ is a genuine expansion in this region and lower truncation in
$z_2$ is its formal signature.
\end{physics}

\begin{status}
Definition and two stability lemmas, unconditional.
\end{status}

\begin{theorem}[Cancellation: $z_1-z_2$ is not a zero-divisor on lower-truncated series]\label{l01:thm:cancel}
If $f\in\loneW[[z_1^{\pm1},z_2^{\pm1}]]$ is lower-truncated in $z_2$ and $(z_1-z_2)^Nf = 0$ for some $N\in\N$, then $f = 0$.
\end{theorem}

\noindent File \leanpath{VOA/\allowbreak Foundations/\allowbreak IteratedLaurent.lean}.
\begin{leancode}
theorem eq_zero_of_mulBySub_of_bddBelow {f : FormalDist₂ 𝒲} (hsub : mulBySub f = 0)
    (hbb : IsBddBelow₂ f) : f = 0 := by

theorem eq_zero_of_mulBySubPow_of_bddBelow :
    ∀ (N : ℕ) (f : FormalDist₂ 𝒲), mulBySubPow N f = 0 → IsBddBelow₂ f → f = 0
\end{leancode}

\begin{dictionary}
Two hypotheses: annihilation by $(z_1-z_2)^N$ and lower truncation. Both are needed: $\lonedelta$ satisfies the
first ($N = 1$) but not the second, and is nonzero. The conclusion is $f = 0$. The first declaration is the case
$N = 1$. Same as the cancellation step in Kac §4 (Goddard uniqueness) and Gannon §5.1.2; slightly stronger, by
the remark on the predicate above.
\end{dictionary}

\begin{physics}
If two operator products agree after multiplication by $(z_1-z_2)^N$, and both are genuine expansions in the
same region, they agree: one may divide by $(z_1-z_2)^N$ inside a fixed radial ordering. This is how locality
(``$(z-w)^N a(z)b(w) = (z-w)^N b(w)a(z)$'') is turned into identities between correlation functions, and why a
field is determined by its action on the vacuum (Goddard).
\end{physics}

\textbf{Proof on paper.} $N = 1$: by Proposition~\ref{l01:prop:Neq1}, $f_{m,n} = f_{m-d,n+d}$ for all $d$. Given
$(m,n)$, take $d$ with $n+d<B$; then $f_{m,n} = f_{m-d,n+d} = 0$. (In the file: $f_{m,n} = f_{-1,m+n+1} =
f_{m+n-(B-1),B-1}$, and then one more step down.) General $N$: induction; $g = (z_1-z_2)^{N}f$ is lower-truncated
and killed by $(z_1-z_2)$, so $g = 0$, and the induction hypothesis applies to $f$. \qed

\begin{example}
$f = \iota_{12}(z_1-z_2)^{-1} = \sum_{b\ge0}z_1^{-1-b}z_2^{b}$ is lower-truncated in $z_2$ (with $B = 0$) and
$(z_1-z_2)f = 1\neq0$, consistent with the theorem. $f = \iota_{21}(z_1-z_2)^{-1}$ is \emph{not} lower-truncated in
$z_2$ (it contains $z_2^{-1-a}$ for all $a\ge0$). Their difference $\lonedelta$ is killed by $z_1-z_2$, and it is not
lower-truncated, as the theorem demands: for any $B$, its coefficient at $(m,n) = (-B,B-1)$ is
$[-B+(B-1)+1 = 0] = 1$, although $n = B-1<B$. In the ring $\C[[z_1^{\pm1}]]((z_2))$, $z_1-z_2$ is in fact invertible with
inverse $\iota_{12}(z_1-z_2)^{-1}$.
\end{example}

\begin{status}
Unconditional; axioms \lean{[propext, Classical.choice, Quot.sound]} (both). Used downstream in
\leanpath{VOA/\allowbreak Core/\allowbreak WeakAssoc.lean}, \leanpath{VOA/\allowbreak Core/\allowbreak Goddard.lean},
\leanpath{VOA/\allowbreak Core/\allowbreak ReconstructionExistence.lean} and six further files.
\end{status}

%======================================================================
\section{The substitution \texorpdfstring{$z_1\mapsto z_0+z_2$}{z1 to z0+z2}}\label{l01:sec:substitution}
%======================================================================

Weak associativity, $(z_0+z_2)^kY(u,z_0+z_2)Y(v,z_2)x = (z_0+z_2)^kY(Y(u,z_0)v,z_2)x$, compares a product of fields at
$z_1$ and $z_2$ with an iterated field in the relative coordinate $z_0 = z_1-z_2$. Formally this requires substituting
$z_1 = z_0+z_2$ into a two-variable series, which is the subject of this section.

\begin{definition}[Substitution]\label{l01:def:subst}
For $f\in\loneW[[z_1^{\pm1},z_2^{\pm1}]]$ lower-truncated in $z_2$, $f(z_0+z_2,z_2)$ is the series in $(z_0,z_2)$ obtained by
replacing each $z_1^m$ by $(z_0+z_2)^m$ expanded in nonnegative powers of $z_2$ (Definition~\ref{l01:def:gbinom}(iv)):
\begin{equation*}
  \big(f(z_0+z_2,z_2)\big)_{a,b} = \sum_{k\ge0}\binom{a+k}{k}\,f_{a+k,\;b-k}.
\end{equation*}
\end{definition}

\noindent File \leanpath{VOA/\allowbreak Foundations/\allowbreak Substitution.lean}.
\begin{leancode}
noncomputable def subst (f : FormalDist₂ 𝒲) : FormalDist₂ 𝒲 :=
  fun a b => ∑ᶠ k : ℤ, cbz (a + k) k • f (a + k) (b - k)

theorem subst_summand_finite {f : FormalDist₂ 𝒲} (hf : IsBddBelow₂ f) (a b lo : ℤ) (c : ℤ → ℂ)
    (hc : ∀ k, k < lo → c k = 0) :
    (Function.support (fun k : ℤ => c k • f (a + k) (b - k))).Finite := by
\end{leancode}

\begin{dictionary}
\textbf{Derivation of the formula.} $f = \sum_{m,n}f_{m,n}z_1^mz_2^n$ and $(z_0+z_2)^m = \sum_{k\ge0}\binom mkz_0^{m-k}z_2^k$. The
coefficient of $z_0^az_2^b$ collects $m-k = a$, $n+k = b$, i.e.\ $m = a+k$, $n = b-k$, with weight $\binom{a+k}{k}$. The sum
is over $k\in\Z$ with $C(\cdot,k) = 0$ for $k<0$ (\lean{cbz}), which is the restriction $k\ge0$.

\textbf{The junk-value trap} (\lonewarn{l01:warn:finsum}). \lean{subst f} is defined for every $f$; it is meaningful
only if each coefficient sum is finite. \lean{subst_summand_finite} shows that this is the case when $f$ is
lower-truncated in $z_2$ (with bound $B$): the summand vanishes for $k<0$ (binomial) and for $b-k<B$ (truncation),
so only $0\le k\le b-B$ contribute. The general form with a coefficient function $c$ vanishing below \lean{lo} is
used for the shifted sums in the next proof. Every theorem about \lean{subst} assumes \lean{IsBddBelow₂ f}.

\textbf{Name.} The docstring calls \lean{subst} ``the $\iota_{z_2,z_0}$ map''; we avoid this name because it suggests
an expansion of a rational function. What it is: the substitution $z_1\mapsto z_0+z_2$ with the FLM convention of
expanding $(z_0+z_2)^m$ in nonnegative powers of $z_2$ (the region $|z_0|>|z_2|$).
\end{dictionary}

\begin{physics}
Substituting $z_1 = z_0+z_2$ re-expresses a two-point operator product in terms of the centre $z_2$ and the
separation $z_0$: this is the passage from $a(z_1)b(z_2)$ to the OPE $\sum_n(a_{(n)}b)(z_2)z_0^{-n-1}$, i.e.\ to the
iterate $Y(Y(a,z_0)b,z_2)$. The expansion region $|z_0|>|z_2|$ is the one in which the OPE converges when
both operators are near $z_2$ but the separation is not small compared with $z_2$; formally it is the region where
$Y(u,z_0+z_2)$ makes sense as a series in nonnegative powers of $z_2$.
\end{physics}

\begin{status}
Definition and finiteness lemma; the lemma is unconditional.
\end{status}

\begin{theorem}[Keystone multiplicativity of the substitution]\label{l01:thm:subst}
If $f$ is lower-truncated in $z_2$, then
\begin{equation*}
  \big((z_1-z_2)f\big)(z_0+z_2,z_2) = z_0\cdot f(z_0+z_2,z_2),\qquad
  \big((z_1-z_2)^Nf\big)(z_0+z_2,z_2) = z_0^N\cdot f(z_0+z_2,z_2)\quad(N\in\N).
\end{equation*}
\end{theorem}

\noindent File \leanpath{VOA/\allowbreak Foundations/\allowbreak Substitution.lean}.
\begin{leancode}
theorem subst_mulBySub {f : FormalDist₂ 𝒲} (hf : IsBddBelow₂ f) :
    subst (mulBySub f) = shiftZ₁ (subst f) := by

theorem subst_mulBySubPow {f : FormalDist₂ 𝒲} (hf : IsBddBelow₂ f) (N : ℕ) :
    subst (mulBySubPow N f) = (shiftZ₁^[N]) (subst f) := by
\end{leancode}

\begin{dictionary}
On the right, the first variable of \lean{subst f} is $z_0$, so \lean{shiftZ₁} (multiplication by the first
variable) is multiplication by $z_0$; \lean{shiftZ₁^[N]} is its $N$-fold iterate (\lean{f^[N]} is Lean's notation for
function iteration). The hypothesis \lean{hf} makes both sides genuine finite sums (note that
$(z_1-z_2)^Nf$ is again lower-truncated, Definition~\ref{l01:def:bddbelow}). This is ``substituting $z_1 = z_0+z_2$ turns
$z_1-z_2$ into $z_0$'', which would be a tautology in a ring of functions and is here a coefficient identity.
\end{dictionary}

\begin{physics}
The separation of the two insertion points is $z_0$: multiplying the product $a(z_1)b(z_2)$ by $(z_1-z_2)^N$ and then
re-centring is the same as re-centring and multiplying the OPE by $z_0^N$, which removes the poles $z_0^{-n-1}$,
$n<N$. This is the step that turns locality into weak associativity, and it is used in the proof of skew-symmetry
(\leanpath{VOA/\allowbreak Core/\allowbreak Skew.lean}).
\end{physics}

\textbf{Proof on paper.} The $(a,b)$ coefficient of the left side is
$\sum_kC(a+k,k)\big(f_{a-1+k,b-k}-f_{a+k,b-1-k}\big)$, and that of the right side is $\sum_kC(a-1+k,k)f_{a-1+k,b-k}$.
Subtracting the right side from the first half of the left side gives
$\sum_k\big(C(a+k,k)-C(a-1+k,k)\big)f_{a-1+k,b-k} = \sum_kC(a-1+k,k-1)f_{a-1+k,b-k}$ by Pascal's rule
(Proposition~\ref{l01:prop:pascal}); shifting $k\mapsto k+1$ turns this into $\sum_kC(a+k,k)f_{a+k,b-1-k}$, the second half of
the left side. All sums are finite by lower truncation, so the manipulations are legitimate. The power version is
induction on $N$, using that $(z_1-z_2)^kf$ is lower-truncated. \qed

\begin{example}[Substituting into $\iota_{12}(z_1-z_2)^{-1}$]\label{l01:ex:subst}
Let $f = \iota_{12}(z_1-z_2)^{-1}$, lower-truncated with $B = 0$. Then
$f(z_0+z_2,z_2)_{a,b} = \sum_{k=0}^{b}\binom{a+k}{k}[a+k = -1-(b-k)]$, which forces $a = -1-b$; for such $a$ the
coefficient is $\sum_{k=0}^b\binom{k-b-1}{k} = \sum_{k=0}^b(-1)^k\binom bk = [b = 0]$ (Example~\ref{l01:ex:negbinom} with
$r = b-k+1$). So $f(z_0+z_2,z_2) = z_0^{-1}$: substituting $z_1 = z_0+z_2$ into $1/(z_1-z_2)$ gives $1/z_0$, as it should.
Consistently with the theorem, $(z_1-z_2)f = 1$ and $1(z_0+z_2,z_2) = 1 = z_0\cdot z_0^{-1}$.
\end{example}

\begin{status}
Unconditional (the hypothesis is the meaningfulness condition); \texttt{\#print axioms} of
\lean{VOA.subst_mulBySub} and \lean{VOA.subst_mulBySubPow} $=$ \lean{[propext, Classical.choice, Quot.sound]}.
\end{status}

%======================================================================
\section{Three-variable distributions}\label{l01:sec:three}
%======================================================================

Dong's lemma (Chapter~\ref{ch:L03}, Theorem~\ref{l03:thm:dong}) is a statement about three fields $a(z_1)$, $b(z_2)$,
$c(z_3)$ and their three pairwise localities. Its proof moves powers of $z_1-z_2$, $z_1-z_3$, $z_2-z_3$ around using
the polynomial identity $z_2-z_3 = (z_2-z_1)+(z_1-z_3)$. The file \leanpath{FormalDist3.lean} provides exactly this.

\begin{definition}[Three-variable distributions and pairwise differences]\label{l01:def:fd3}
A three-variable distribution is $f:\Z^3\to\loneW$, $f = \sum f_{m,n,p}z_1^mz_2^nz_3^p$. Multiplication by $z_i$ shifts the
$i$-th index by one, and multiplication by $z_i-z_j$ is the difference of two shifts; $(z_i-z_j)^N$ is its $N$-fold
iterate.
\end{definition}

\noindent File \leanpath{VOA/\allowbreak Foundations/\allowbreak FormalDist3.lean}.
\begin{leancode}
def FormalDist₃ (𝒲 : Type*) : Type _ := ℤ → ℤ → ℤ → 𝒲

def shift₁ : FormalDist₃ 𝒲 →ₗ[ℂ] FormalDist₃ 𝒲 where
  toFun f := fun m n p => f (m - 1) n p

def mulBySub₁₂ : FormalDist₃ 𝒲 →ₗ[ℂ] FormalDist₃ 𝒲 := shift₁ - shift₂
def mulBySub₁₃ : FormalDist₃ 𝒲 →ₗ[ℂ] FormalDist₃ 𝒲 := shift₁ - shift₃
def mulBySub₂₃ : FormalDist₃ 𝒲 →ₗ[ℂ] FormalDist₃ 𝒲 := shift₂ - shift₃
def mulBySub₂₁ : FormalDist₃ 𝒲 →ₗ[ℂ] FormalDist₃ 𝒲 := shift₂ - shift₁

def mulBySub₁₂Pow : ℕ → (FormalDist₃ 𝒲 →ₗ[ℂ] FormalDist₃ 𝒲)
  | 0 => LinearMap.id
  | (N + 1) => mulBySub₁₂.comp (mulBySub₁₂Pow N)
\end{leancode}

\begin{dictionary}
\lean{f m n p} is the coefficient of $z_1^mz_2^nz_3^p$. \lean{shift₂}, \lean{shift₃} are defined like \lean{shift₁} on the
second and third index (omitted). The powers \lean{mulBySub₁₃Pow}, \lean{mulBySub₂₃Pow}, \lean{mulBySub₂₁Pow} are defined
like \lean{mulBySub₁₂Pow}. The file proves \lean{mulBySub₁₂_apply : mulBySub₁₂ f m n p = f (m - 1) n p - f m (n - 1) p} and
its analogues. Same as the textbook three-variable calculus (Kac §3.2, Lepowsky--Li §3.2).
\end{dictionary}

\begin{physics}
Three-variable distributions are three-point operator products $a(z_1)b(z_2)c(z_3)$; the three differences are the
three collision channels. Associativity of the OPE is the statement that the channels $z_1\to z_2$ and $z_2\to z_3$
give consistent answers, and Dong's lemma is its first formal consequence.
\end{physics}

\begin{status}
Definitions; no hypotheses.
\end{status}

\begin{theorem}[The binomial expansion of $(z_2-z_3)^N$]\label{l01:thm:fd3binom}
On three-variable distributions, $z_2-z_3 = (z_2-z_1)+(z_1-z_3)$, $z_2-z_1 = -(z_1-z_2)$, and for every $N\in\N$
\begin{equation*}
  (z_2-z_3)^Nf = \sum_{i=0}^{N}\binom Ni\,(z_2-z_1)^i(z_1-z_3)^{N-i}f .
\end{equation*}
\end{theorem}

\noindent File \leanpath{VOA/\allowbreak Foundations/\allowbreak FormalDist3.lean}.
\begin{leancode}
theorem mulBySub₂₃_eq : (mulBySub₂₃ : FormalDist₃ 𝒲 →ₗ[ℂ] _) = mulBySub₂₁ + mulBySub₁₃ := by

theorem mulBySub₂₁_eq_neg : (mulBySub₂₁ : FormalDist₃ 𝒲 →ₗ[ℂ] _) = -mulBySub₁₂ := by

theorem mulBySub₂₃Pow_eq_sum (N : ℕ) (f : FormalDist₃ 𝒲) :
    mulBySub₂₃Pow N f
      = ∑ i ∈ Finset.range (N + 1),
          (N.choose i) • mulBySub₂₁Pow i (mulBySub₁₃Pow (N - i) f) := by
\end{leancode}

\begin{dictionary}
\lean{N.choose i} is the ordinary binomial coefficient (a natural number, acting by repeated addition). \lean{N - i} is
truncated subtraction on $\N$, harmless because $i\le N$ in the range. The identity is an identity of linear maps
applied to $f$; it rests on the fact that all shift operators commute, so that the binomial theorem holds in the
commutative subalgebra they generate (the file proves \lean{mulBySub₂₁_comm_mulBySub₁₃}). No hypotheses.
\end{dictionary}

\begin{physics}
In Dong's lemma one must show that $(z_2-z_3)^{M}$ kills $[(a_{(k)}b)(z_2),c(z_3)]$. Writing $z_2-z_3$ as
$(z_2-z_1)+(z_1-z_3)$ and expanding distributes the factor between the channels in which $a$ is close to $b$ and in
which $a$ is close to $c$; each term is killed by one of the given localities.
\end{physics}

\textbf{Proof on paper.} The operators of multiplication by $z_1,z_2,z_3$ commute, so the binomial theorem applies
in the commutative algebra they generate. \qed

\begin{status}
Unconditional; axioms clean (checked for \lean{mulBySub₂₃_eq} and \lean{mulBySub₂₃Pow_eq_sum}). Used downstream in
\leanpath{VOA/\allowbreak Core/\allowbreak Dong.lean} and \leanpath{VOA/\allowbreak Core/\allowbreak DongSuper.lean}.
\end{status}

%======================================================================
\section{Four-variable distributions and the double-delta calculus}\label{l01:sec:four}
%======================================================================

The last four files build a four-variable calculus intended for the proof of Li's local-system theorem: a set of
mutually local fields on a space $W$ generates a vertex algebra whose vertex operator is
\begin{equation*}
  Y_W(a(z),z_0)\,b(z) = \Res_{z_1}\Big(z_0^{-1}\delta\Big(\frac{z_1-z}{z_0}\Big)a(z_1)b(z) -
  z_0^{-1}\delta\Big(\frac{z-z_1}{-z_0}\Big)b(z)a(z_1)\Big).
\end{equation*}
To prove that two such induced vertex operators are again local one needs to apply \emph{two} such delta kernels,
sharing the base variable $z$, and to move powers of differences through the resulting double residue. We present
the layer completely, but we warn at the outset (\lonewarn{l01:warn:shiftfinite}) that at this commit the layer is not
used by any other file of the library, and that its main theorem carries a finiteness hypothesis that has not yet been
discharged for the intended payloads.

\subsection{Four-variable distributions, residues and Fubini}

\begin{definition}[Four-variable distributions, differences, residues]\label{l01:def:fd4}
A four-variable distribution is $f:\Z^4\to W$, $f = \sum f_{i,j,k,l}z_1^iz_2^jz_3^kz_4^l$. For each pair $i<j$ there is
multiplication by $z_i-z_j$ and its powers. For each coordinate (resp.\ pair of coordinates) there is a residue,
the coefficient at exponent $-1$ (resp.\ $(-1,-1)$), which is a three- (resp.\ two-)variable distribution in the
remaining coordinates, in their original order.
\end{definition}

\noindent File \leanpath{VOA/\allowbreak Foundations/\allowbreak FormalDist4.lean}.
\begin{leancode}
def FormalDist₄ (W : Type*) : Type _ := ℤ → ℤ → ℤ → ℤ → W

def mulBySub₁₂ : FormalDist₄ W →ₗ[ℂ] FormalDist₄ W := shift₁ - shift₂

def mulBySub₃₄ : FormalDist₄ W →ₗ[ℂ] FormalDist₄ W := shift₃ - shift₄

def mulBySub₁₂Pow (N : ℕ) : FormalDist₄ W →ₗ[ℂ] FormalDist₄ W := mulBySub₁₂ ^ N

def res₁ : FormalDist₄ W →ₗ[ℂ] FormalDist₃ W where
  toFun f := fun j k l => f (-1) j k l

def res₁₂ : FormalDist₄ W →ₗ[ℂ] FormalDist₂ W where
  toFun f := fun k l => f (-1) (-1) k l
\end{leancode}

\begin{dictionary}
These live in the namespace \lean{VOA.FormalDist₄}, so \lean{FormalDist₄.mulBySub₁₂} is distinct from the three-variable
\lean{VOA.mulBySub₁₂} of Definition~\ref{l01:def:fd3}. All six differences \lean{mulBySub₁₂}, \dots, \lean{mulBySub₃₄}, all six
powers, the four single residues \lean{res₁}--\lean{res₄} and the six double residues \lean{res₁₂}--\lean{res₃₄} are defined in the
same way (we show representatives). Here powers are monoid powers \lean{mulBySub₁₂ ^ N} in the algebra of linear
endomorphisms, rather than the recursion used for two and three variables; the file proves the recursion
\lean{mulBySub₁₂Pow_succ}. Residues are pure coefficient extraction and need no finiteness.
\end{dictionary}

\begin{physics}
Four-point configurations: two field insertions $x_1,x_2$ that will be contracted (integrated) and two relative
coordinates $z_{o_1},z_{o_2}$ that remain; the residues are the contour integrals over $x_1$ and $x_2$.
\end{physics}

\begin{status}
Definitions; linearity by \lean{rfl}. No hypotheses.
\end{status}

\begin{definition}[Coefficientwise finite families and their sums]\label{l01:def:coeffsum}
A family $(F_a)_{a\in I}$ of four-variable distributions is \emph{coefficientwise finite} if for every multi-index the set
of $a$ with nonzero coefficient is finite; a two-parameter family $(F_{a,b})$ is \emph{jointly coefficientwise finite} if
for every multi-index the set of pairs $(a,b)$ with nonzero coefficient is finite. The (iterated) coefficientwise sums
are $\big(\sum_aF_a\big)_{ijkl} = \sum_a(F_a)_{ijkl}$ and $\sum_a\sum_bF_{a,b}$.
\end{definition}

\noindent File \leanpath{VOA/\allowbreak Foundations/\allowbreak FormalDist4.lean}.
\begin{leancode}
def CoeffFinite {ι : Type v} (F : ι → FormalDist₄ W) : Prop :=
  ∀ i j k l, (Function.support fun a => F a i j k l).Finite

def CoeffJointFinite {ι : Type v} {κ : Type w} (F : ι → κ → FormalDist₄ W) : Prop :=
  ∀ i j k l, (Function.support fun p : ι × κ => F p.1 p.2 i j k l).Finite

noncomputable def coeffFinsum {ι : Type v} (F : ι → FormalDist₄ W) : FormalDist₄ W :=
  fun i j k l => ∑ᶠ a, F a i j k l

noncomputable def coeffFinsum₂ {ι : Type v} {κ : Type w} (F : ι → κ → FormalDist₄ W) :
    FormalDist₄ W := coeffFinsum fun a => coeffFinsum (F a)
\end{leancode}

\begin{dictionary}
\lean{ι × κ} is the product type; \lean{p.1}, \lean{p.2} its components. Unlike \lean{algSum}
(Definition~\ref{l01:def:algsum}), \lean{coeffFinsum} does \emph{not} take a finiteness proof as an argument: it is
defined for all families and returns the junk value on infinite supports (\lonewarn{l01:warn:finsum}); the predicates
\lean{CoeffFinite} and \lean{CoeffJointFinite} are carried separately by the theorems. The file also proves that residues
commute with these sums (\lean{res₁_coeffFinsum_apply} etc.); those lemmas are definitional and carry an \emph{unused}
finiteness hypothesis (named \lean{_hF}), kept, according to the file, so that callers retain the certificate.
\end{dictionary}

\begin{physics}
Bookkeeping for sums over intermediate modes, which in a mode computation of a four-point function are finite on
each state.
\end{physics}

\begin{status}
Definitions.
\end{status}

\begin{theorem}[Fubini for algebraic sums]\label{l01:thm:fubini}
If the doubly indexed family $(F_{a,b})$ in an abelian group has finite joint support, then
$\sum_b\sum_aF_{a,b} = \sum_a\sum_bF_{a,b}$. Iterated sums are invariant under reindexing both variables by
bijections. Consequently coefficientwise iterated sums of a jointly coefficientwise finite family may be exchanged.
\end{theorem}

\noindent File \leanpath{VOA/\allowbreak Foundations/\allowbreak FormalDist4.lean}.
\begin{leancode}
theorem finsum_finsum_comm_of_joint_finite {ι : Type v} {κ : Type w}
    (F : ι → κ → W)
    (hF : (Function.support fun p : ι × κ => F p.1 p.2).Finite) :
    (∑ᶠ b, ∑ᶠ a, F a b) = ∑ᶠ a, ∑ᶠ b, F a b := by

theorem finsum_finsum_reindex {ι ι' : Type v} {κ κ' : Type w}
    (eι : ι ≃ ι') (eκ : κ ≃ κ') (F : ι' → κ' → W) :
    (∑ᶠ a, ∑ᶠ b, F (eι a) (eκ b)) = ∑ᶠ a', ∑ᶠ b', F a' b' := by

theorem coeffFinsum₂_comm {ι : Type v} {κ : Type w}
    (F : ι → κ → FormalDist₄ W) (hF : CoeffJointFinite F) :
    coeffFinsum₂ F = coeffFinsum₂ (fun b a => F a b) := by
\end{leancode}

\begin{dictionary}
\lean{ι ≃ ι'} is a bijection (an equivalence of types). The reindexing lemma needs no finiteness: \lean{finsum} is
invariant under bijections even in the junk case. The exchange lemma needs \emph{joint} finiteness, and this
hypothesis is genuinely necessary, not an artefact of Lean (see the warning below).
\end{dictionary}

\begin{physics}
Exchanging the order of two contour integrals, $\oint_{x_1}\oint_{x_2} = \oint_{x_2}\oint_{x_1}$, is legitimate only when the
integrand has no singularity that one contour must cross; formally, when only finitely many intermediate modes
contribute jointly.
\end{physics}

\textbf{Proof on paper.} With finite joint support $S\subseteq I\times K$, both iterated sums equal $\sum_{(a,b)\in S}F_{a,b}$. \qed

\begin{warning}\label{l01:warn:fubini}
\textbf{Why joint finiteness and not just finiteness of each inner sum.} On $\N\times\N$ let $F_{a,b} = [a = b]-[a = b+1]$.
For each $a$, $b\mapsto F_{a,b}$ has finite support and $\sum_bF_{a,b} = [a = 0]$; so $\sum_a\sum_bF_{a,b} = 1$. For each
$b$, $\sum_aF_{a,b} = 1-1 = 0$; so $\sum_b\sum_aF_{a,b} = 0$. All inner and outer sums here are genuine finite sums, yet the
two iterated sums differ. The joint support (the two diagonals) is infinite. The library's insistence on
\lean{CoeffJointFinite} and on \lean{DoubleResidueFinite} below is therefore mathematically necessary.
\end{warning}

\begin{status}
Unconditional; axioms clean (checked for \lean{finsum_finsum_comm_of_joint_finite} and \lean{coeffFinsum₂_comm}).
\end{status}

\subsection{The single and double delta kernels}

\begin{definition}[Orientations, single and double delta kernels]\label{l01:def:singledelta}
There are two \emph{orientations}, $+$ and $-$. For an orientation $o$, the three-variable kernel
$\Delta_o(x,z_b,z_o)$ has coefficient of $x^{s}z_b^{\beta}z_o^{r}$ equal to the coefficient of $x^sz_b^\beta$ in
$\iota_o(x-z_b)^{-r-1}$, where $\iota_+ = \iota_{x,z_b}$ and $\iota_- = \iota_{z_b,x}$. Thus
\begin{align*}
  \Delta_+ &= z_o^{-1}\delta\Big(\frac{x-z_b}{z_o}\Big)\ \text{expanded in nonnegative powers of }z_b,\\
  \Delta_- &= z_o^{-1}\delta\Big(\frac{x-z_b}{z_o}\Big)\ \text{expanded in nonnegative powers of }x .
\end{align*}
For orientations $o_1,o_2$ the \emph{double delta kernel} is the product sharing the base variable $z_b$:
\begin{equation*}
  K_{o_1o_2}(x_1,x_2,z_{o_1},z_{o_2};z_b) := \Delta_{o_1}(x_1,z_b,z_{o_1})\,\Delta_{o_2}(x_2,z_b,z_{o_2}),
\end{equation*}
a four-variable distribution in $(x_1,x_2,z_{o_1},z_{o_2})$ whose coefficients are one-variable series in $z_b$.
\end{definition}

\noindent File \leanpath{VOA/\allowbreak Foundations/\allowbreak FormalDist4Delta.lean}.
\begin{leancode}
inductive Orientation
  | pos
  | neg
  deriving DecidableEq

noncomputable def singleDelta (o : Orientation) : FormalDist₃ ℂ
  | source, base, outer => match o with
    | .pos => expandPos (-outer - 1) source base
    | .neg => expandNeg (-outer - 1) source base

noncomputable def doubleDelta (o₁ o₂ : Orientation) : FormalDist₄ (FormalDist ℂ) :=
  fun source₁ source₂ outer₁ outer₂ base => ∑ᶠ split : ℤ,
    singleDelta o₁ source₁ split outer₁ *
      singleDelta o₂ source₂ (base - split) outer₂
\end{leancode}

\begin{dictionary}
\textbf{Decoding \lean{singleDelta}.} The three indices are (source $s$, base $\beta$, outer $r$): exponents of $x$, $z_b$, $z_o$.
The value is the $(s,\beta)$ coefficient of $\iota(x-z_b)^{-r-1}$ (Definition~\ref{l01:def:expand}). Summing over $r$:
$\sum_rz_o^{r}(x-z_b)^{-r-1} = \sum_nz_o^{-n-1}(x-z_b)^{n}$ ($n = -r-1$) $= z_o^{-1}\delta((x-z_b)/z_o)$, so the kernel is the
FLM delta function in the stated expansion. One checks that the $-$ orientation is the second FLM term: with the binomial
convention (expansion in nonnegative powers of the second variable),
$z_0^{-1}\delta((z-x)/(-z_0)) = \sum_n(-1)^nz_0^{-n-1}(z-x)^n$ with $(z-x)^n$ expanded in nonnegative powers of $x$, and
$(-1)^n(z-x)^n$ so expanded is $\iota_{z,x}(x-z)^n$; hence $z_0^{-1}\delta((z-x)/(-z_0)) = \sum_nz_0^{-n-1}\iota_{z,x}(x-z)^n$, which
is $\Delta_-$ with base $z_b = z$.

\textbf{Decoding \lean{doubleDelta}.} \lean{FormalDist₄ (FormalDist ℂ)}: four explicit indices
$(s_1,s_2,r_1,r_2)$ and a fifth, \lean{base}, for the $z_b$-series. The value is the $z_b^{\text{base}}$ coefficient of the product
$\Delta_{o_1}\Delta_{o_2}$, a convolution in $z_b$ over \lean{split}. It is a genuine finite sum: the first factor vanishes
unless $\text{split} = -r_1-1-s_1$ (antidiagonal support), proved as \lean{singleDelta_convolution_finite}. The four
choices of $(o_1,o_2)$ are the PP/PN/NP/NN cases of the file header.
\end{dictionary}

\begin{physics}
$z_0^{-1}\delta((z_1-z_2)/z_0)$ implements the change of variables $z_1 = z_2+z_0$ inside a contour integral: it is the
kernel that converts the product $a(z_1)b(z_2)$ into the OPE in the relative coordinate $z_0$. The two orientations
are the two radial orderings $a(z_1)b(z_2)$ and $b(z_2)a(z_1)$ that appear, with a relative sign, in the definition of
the induced vertex operator $Y_W$; the double kernel performs two such re-centrings around the same base point,
as needed for the three-point configuration $a_1(x_1)\,b(z_b)\,a_2(x_2)$.
\end{physics}

\begin{status}
Definitions.
\end{status}

\begin{proposition}[Support and the defining relation of the single kernel]\label{l01:prop:singledelta}
$\Delta_o$ is supported on $s+\beta = -r-1$; the convolution defining $K_{o_1o_2}$ has finitely many terms for each
coefficient; and $(x-z_b)\Delta_o = z_o\Delta_o$ (on the support of the delta function, $x-z_b = z_o$).
\end{proposition}

\noindent File \leanpath{VOA/\allowbreak Foundations/\allowbreak FormalDist4Delta.lean}.
\begin{leancode}
theorem singleDelta_eq_zero_off_diag (o : Orientation) (source base outer : ℤ)
    (h : source + base ≠ -outer - 1) : singleDelta o source base outer = 0 := by

theorem singleDelta_convolution_finite (o₁ o₂ : Orientation)
    (source₁ source₂ outer₁ outer₂ base : ℤ) :
    (Function.support fun split : ℤ =>
      singleDelta o₁ source₁ split outer₁ *
        singleDelta o₂ source₂ (base - split) outer₂).Finite := by

theorem singleDelta_relation (o : Orientation) (source base outer : ℤ) :
    singleDelta o (source - 1) base outer - singleDelta o source (base - 1) outer =
      singleDelta o source base (outer - 1) := by
\end{leancode}

\begin{dictionary}
In \lean{singleDelta_relation}, the left side is the coefficient of $(x-z_b)\Delta_o$ (shift in $x$ minus shift in $z_b$) and
the right side that of $z_o\Delta_o$ (shift in $z_o$). It is Proposition~\ref{l01:prop:mulexpand} with $n = -r-1$.
\end{dictionary}

\begin{physics}
The delta function enforces momentum conservation in the relative coordinate: $x = z_b+z_o$.
\end{physics}

\textbf{Proof on paper.} Proposition~\ref{l01:prop:expandvals}(b) and Proposition~\ref{l01:prop:mulexpand}. \qed

\begin{status}
Unconditional.
\end{status}

\begin{theorem}[Outer-to-inner transport]\label{l01:thm:outerinner}
For all orientations $o_1,o_2$ and $N\in\N$,
\begin{equation*}
  (z_{o_1}-z_{o_2})^N\,K_{o_1o_2} = (x_1-x_2)^N\,K_{o_1o_2}.
\end{equation*}
\end{theorem}

\noindent File \leanpath{VOA/\allowbreak Foundations/\allowbreak FormalDist4Delta.lean}.
\begin{leancode}
theorem mulBySub₁₂_comm_mulBySub₃₄ :
    Commute
      (FormalDist₄.mulBySub₁₂ : FormalDist₄ (FormalDist ℂ) →ₗ[ℂ] _)
      (FormalDist₄.mulBySub₃₄ : FormalDist₄ (FormalDist ℂ) →ₗ[ℂ] _) := by

theorem mulBySub₃₄_doubleDelta (o₁ o₂ : Orientation) :
    FormalDist₄.mulBySub₃₄ (doubleDelta o₁ o₂) =
      FormalDist₄.mulBySub₁₂ (doubleDelta o₁ o₂) := by

theorem mulBySub₃₄Pow_doubleDelta (o₁ o₂ : Orientation) (N : ℕ) :
    FormalDist₄.mulBySub₃₄Pow N (doubleDelta o₁ o₂) =
      FormalDist₄.mulBySub₁₂Pow N (doubleDelta o₁ o₂) := by
\end{leancode}

\begin{dictionary}
\textbf{Which variables.} The four explicit coordinates of \lean{doubleDelta} are $(s_1,s_2,r_1,r_2)$, i.e.\ $(x_1,x_2,z_{o_1},z_{o_2})$.
So \lean{mulBySub₃₄} is $z_{o_1}-z_{o_2}$ and \lean{mulBySub₁₂} is $x_1-x_2$. An earlier file header stated the identity as
``$(z_3-z_0)^NK = (z_1-z_4)^NK$'', in a variable naming taken from Li's paper; the current header writes
``$(z_3-z_4)^NK = (z_1-z_2)^NK$'' with $(z_1,z_2,z_3,z_4) = (\text{source}_1,\text{source}_2,\text{outer}_1,\text{outer}_2)$, which
matches the Lean coordinates. It is the same statement.
\lean{Commute A B} means $AB = BA$. No hypotheses; valid uniformly in the four orientations.
\end{dictionary}

\begin{physics}
On the support of the kernel $x_i = z_b+z_{o_i}$, so $x_1-x_2 = z_{o_1}-z_{o_2}$: the separation of the two original
insertion points equals the separation of their relative coordinates. This is the identity that transfers locality of
$a_1,a_2$ to locality of the induced vertex operators.
\end{physics}

\textbf{Proof on paper.} By Proposition~\ref{l01:prop:singledelta},
$z_{o_1}\Delta_1\Delta_2 - z_{o_2}\Delta_1\Delta_2 = ((x_1-z_b)\Delta_1)\Delta_2-\Delta_1((x_2-z_b)\Delta_2) = (x_1-x_2)\Delta_1\Delta_2$, the two
$z_b$-terms cancelling. Coefficientwise the cancellation is an index shift in the convolution variable \lean{split}, legitimate
because all convolution sums are finite. The power version follows by induction, using that the two difference
operators commute. \qed

\begin{status}
Unconditional; \texttt{\#print axioms} of \lean{mulBySub₃₄_doubleDelta} and \lean{mulBySub₃₄Pow_doubleDelta} $=$
\lean{[propext, Classical.choice, Quot.sound]} (also printed by the file's own audit section).
\end{status}

\subsection{Contraction: the semantic interface}

\begin{definition}[Single contraction and residue substitution]\label{l01:def:contraction}
For an orientation $o$ and $r,k\in\Z$ let $\gamma_o(r,k)$ be the coefficient of $\Delta_o$ at $(s,\beta,r) = (-r-k-1,k,r)$. For a
two-variable payload $f(x,z_b)$, the \emph{single contraction} is
$\big(\Res_x\Delta_o(x,z_b,z_o)f(x,z_b)\big)$, whose $(r,\beta')$ coefficient is
$\sum_k\gamma_o(r,k)\,f_{r+k,\beta'-k}$ (\emph{residue substitution}).
\end{definition}

\noindent File \leanpath{VOA/\allowbreak Foundations/\allowbreak FormalDist4DeltaContract.lean}.
\begin{leancode}
noncomputable def contractionCoeff (o : Orientation) (outer k : ℤ) : ℂ :=
  singleDelta o (-outer - k - 1) k outer

noncomputable def contractionKernel (o : Orientation) : FormalDist₄ ℂ :=
  fun outer baseOut source baseIn =>
    singleDelta o (-source - 1) (baseOut - baseIn) outer

noncomputable def singleContractSummand (o : Orientation) (f : FormalDist₂ W)
    (outer baseOut k : ℤ) : W :=
  contractionKernel o outer baseOut (outer + k) (baseOut - k) •
    f (outer + k) (baseOut - k)

noncomputable def singleDeltaContract (o : Orientation) (f : FormalDist₂ W) :
    FormalDist₂ W :=
  fun outer baseOut => ∑ᶠ k : ℤ, singleContractSummand o f outer baseOut k

noncomputable def residueSubst (o : Orientation) (f : FormalDist₂ W) : FormalDist₂ W :=
  fun outer baseOut => ∑ᶠ k : ℤ,
    contractionCoeff o outer k • f (outer + k) (baseOut - k)
\end{leancode}

\begin{dictionary}
\textbf{The matrix of the contraction.} \lean{contractionKernel o r β' s β} is the coefficient with which the input
coefficient $f_{s,\beta}$ contributes to the output coefficient $(r,\beta')$: pairing $x^{s}$ in $f$ with $x^{-s-1}$ in the
kernel (the residue in $x$) and $z_b^{\beta}$ in $f$ with $z_b^{\beta'-\beta}$ in the kernel (the product in $z_b$). It is
supported on the ``residue graph'' $s = r+(\beta'-\beta)$, parametrised by $k = \beta'-\beta$; on it the coefficient is
\lean{contractionCoeff o r k}.

\textbf{Junk values.} \lean{singleDeltaContract} and \lean{residueSubst} are \lean{finsum}s over $k\in\Z$ and are meaningful only
when finitely many $k$ contribute. For $o = +$, $\gamma_+(r,k) = (-1)^k\binom{-r-1}{k} = \binom{r+k}{k}$ for $k\ge0$ and $0$ for
$k<0$ (\lean{contractionCoeff_pos_eq_zero_of_neg}), so \lean{residueSubst .pos} has exactly the coefficient formula of
\lean{subst} (Definition~\ref{l01:def:subst}) and is meaningful on payloads lower-truncated in $z_b$. For $o = -$, $\gamma_-(r,k)=0$
unless $k\le-r-1$, i.e.\ unless the source exponent $r+k$ is $\le-1$ (\lean{contractionCoeff_neg_eq_zero_of_gt}).

\textbf{Design decision recorded in the file.} The header insists that the contraction coefficient be \emph{derived} from
\lean{singleDelta}, not postulated by a separate binomial formula, and two regression checks pin values for the negative
orientation (\lean{contractionCoeff_neg_zero_neg_one}, \lean{contractionCoeff_neg_neg_one_zero}); the second docstring says it is
the value ``where the discarded binomial interface incorrectly returned zero''. See the warning below.
\end{dictionary}

\begin{physics}
$\Res_x z_o^{-1}\delta((x-z_b)/z_o)\,f(x,z_b)$ is ``$f(z_b+z_o,z_b)$'': the contour integral over $x$ against the delta kernel
re-centres the insertion point at $z_b$ with separation $z_o$. It is the formal version of Taylor-expanding a field
around a nearby point, $a(z_b+z_o) = e^{z_o\partial}a(z_b)$, in the orientation $+$.
\end{physics}

\begin{status}
Definitions.
\end{status}

\begin{proposition}[The contraction is the residue substitution; support and regression values]\label{l01:prop:contraction}
For every payload $f$ and orientation $o$, the single contraction equals the residue substitution. The kernel vanishes off
the residue graph; $\gamma_+(r,k) = 0$ for $k<0$; $\gamma_-(r,k) = 0$ for $k>-r-1$; and $\gamma_-(0,-1) = -1$, $\gamma_-(-1,0) = 1$.
\end{proposition}

\noindent File \leanpath{VOA/\allowbreak Foundations/\allowbreak FormalDist4DeltaContract.lean}.
\begin{leancode}
theorem singleDeltaContract_eq_residueSubst (o : Orientation) (f : FormalDist₂ W) :
    singleDeltaContract o f = residueSubst o f := by

theorem contractionKernel_eq_zero_off_graph (o : Orientation)
    (outer baseOut source baseIn : ℤ)
    (h : source ≠ outer + (baseOut - baseIn)) :
    contractionKernel o outer baseOut source baseIn = 0 := by

theorem contractionCoeff_pos_eq_zero_of_neg (outer : ℤ) {k : ℤ} (hk : k < 0) :
    contractionCoeff .pos outer k = 0 := by

theorem contractionCoeff_neg_eq_zero_of_gt (outer : ℤ) {k : ℤ}
    (hk : -outer - 1 < k) : contractionCoeff .neg outer k = 0 := by

theorem contractionCoeff_neg_zero_neg_one : contractionCoeff .neg 0 (-1) = -1 := by

theorem contractionCoeff_neg_neg_one_zero : contractionCoeff .neg (-1) 0 = 1 := by
\end{leancode}

\begin{dictionary}
The first identity holds summand by summand (no finiteness needed, since both sides are \lean{finsum}s of the same family).
The regression values by hand: $\gamma_-(0,-1)$ is the coefficient of $x^{0}z_b^{-1}$ in $\iota_{z_b,x}(x-z_b)^{-1} =
-\sum_{a\ge0}x^az_b^{-1-a}$, i.e.\ $-1$; and $\gamma_-(-1,0)$ is the coefficient of $x^0z_b^0$ in $\iota_{z_b,x}(x-z_b)^{0} = 1$.
\end{dictionary}

\begin{physics}
Technical: it identifies the abstract contraction with the explicit re-centring formula and fixes the two
orientation conventions by explicit values, the formal analogue of checking a sign convention on the simplest
correlator.
\end{physics}

\textbf{Proof on paper.} Reparametrise the summation by $k$; the support statements are those of
Proposition~\ref{l01:prop:expandvals}(b); the values are computed above. \qed

\begin{status}
Unconditional; axioms clean (checked for \lean{singleDeltaContract_eq_residueSubst} and both regression values).
\end{status}

\begin{warning}\label{l01:warn:binomialinterface}
\textbf{A corrected mistake: the negative orientation is not ``the positive binomial with a different index''.} The header
of \leanpath{FormalDist4DeltaContract.lean} records that an earlier interface computed the negative-orientation contraction
coefficient by reusing the positive-expansion binomial formula, and that it returned $0$ at $(r,k) = (-1,0)$, where the
correct value is $1$. The present interface derives every coefficient from \lean{singleDelta}, and
\lean{contractionCoeff_neg_neg_one_zero} is the regression certificate. The lesson for the reader: the two expansions of
$(x-z_b)^n$ are supported in \emph{opposite} half-planes of exponents (Proposition~\ref{l01:prop:expandvals}(b)), so no
reindexing of one formula produces the other.
\end{warning}

\begin{definition}[Double contraction and double residue]\label{l01:def:kernelcontract}
Consider a kernel and a payload
\begin{equation*}
  K\in\C[[z_b^{\pm1}]][[x_1^{\pm1},x_2^{\pm1},z_{o_1}^{\pm1},z_{o_2}^{\pm1}]],\qquad
  f\in W[[x_1^{\pm1},z_b^{\pm1},x_2^{\pm1}]]
\end{equation*}
(coordinates in this order). The \emph{double contraction} in the order $(1,2)$ is
\begin{equation*}
  \big(\Res_{x_1}\Res_{x_2}K\,f\big)_{r_1,\beta',r_2} = \sum_{s_1}\sum_{s_2}\sum_{\beta}
  K_{-s_1-1,\,-s_2-1,\,r_1,\,r_2}(\beta'-\beta)\;f_{s_1,\beta,s_2},
\end{equation*}
an element of $W[[z_{o_1}^{\pm1},z_b^{\pm1},z_{o_2}^{\pm1}]]$, and in the order $(2,1)$ with the outer two sums exchanged. For the
double delta kernel the corresponding \emph{double residue} is
$\sum_{k_1}\sum_{k_2}\gamma_{o_1}(r_1,k_1)\gamma_{o_2}(r_2,k_2)\,f_{r_1+k_1,\;\beta'-k_1-k_2,\;r_2+k_2}$, in either order.
\end{definition}

\noindent File \leanpath{VOA/\allowbreak Foundations/\allowbreak FormalDist4DeltaContract.lean}.
\begin{leancode}
noncomputable def kernelContractSummand (K : FormalDist₄ (FormalDist ℂ))
    (f : FormalDist₃ W) (outer₁ baseOut outer₂ source₁ source₂ baseIn : ℤ) : W :=
  K (-source₁ - 1) (-source₂ - 1) outer₁ outer₂ (baseOut - baseIn) •
    f source₁ baseIn source₂

noncomputable def kernelContract₁₂ (K : FormalDist₄ (FormalDist ℂ))
    (f : FormalDist₃ W) : FormalDist₃ W :=
  fun outer₁ baseOut outer₂ => ∑ᶠ source₁ : ℤ, ∑ᶠ source₂ : ℤ, ∑ᶠ baseIn : ℤ,
    kernelContractSummand K f outer₁ baseOut outer₂ source₁ source₂ baseIn

noncomputable def kernelContract₂₁ (K : FormalDist₄ (FormalDist ℂ))
    (f : FormalDist₃ W) : FormalDist₃ W :=
  fun outer₁ baseOut outer₂ => ∑ᶠ source₂ : ℤ, ∑ᶠ source₁ : ℤ, ∑ᶠ baseIn : ℤ,
    kernelContractSummand K f outer₁ baseOut outer₂ source₁ source₂ baseIn

noncomputable def doubleResidueSummand (o₁ o₂ : Orientation) (f : FormalDist₃ W)
    (outer₁ baseOut outer₂ k₁ k₂ : ℤ) : W :=
  (contractionCoeff o₁ outer₁ k₁ * contractionCoeff o₂ outer₂ k₂) •
    f (outer₁ + k₁) (baseOut - k₁ - k₂) (outer₂ + k₂)

def DoubleResidueFinite (o₁ o₂ : Orientation) (f : FormalDist₃ W) : Prop :=
  ∀ outer₁ baseOut outer₂,
    (Function.support fun k : ℤ × ℤ =>
      doubleResidueSummand o₁ o₂ f outer₁ baseOut outer₂ k.1 k.2).Finite

noncomputable def doubleResidue₁₂ (o₁ o₂ : Orientation) (f : FormalDist₃ W) :
    FormalDist₃ W :=
  fun outer₁ baseOut outer₂ => ∑ᶠ k₁ : ℤ, ∑ᶠ k₂ : ℤ,
    doubleResidueSummand o₁ o₂ f outer₁ baseOut outer₂ k₁ k₂
\end{leancode}

\begin{dictionary}
\textbf{Coordinates.} Payload \lean{f source₁ baseIn source₂}: the base variable is in the \emph{middle}. Output
\lean{(outer₁, baseOut, outer₂)}: again base in the middle. So on payloads and outputs (both \lean{FormalDist₃}),
the three-variable \lean{mulBySub₁₃} of Definition~\ref{l01:def:fd3} is $x_1-x_2$, respectively $z_{o_1}-z_{o_2}$.
\lean{doubleResidue₂₁} is \lean{doubleResidue₁₂} with the two sums exchanged (omitted).
\textbf{Junk values.} All five constructions are iterated \lean{finsum}s with no finiteness built in.
\lean{DoubleResidueFinite} is the joint-finiteness predicate for the double residue (cf.\ \lonewarn{l01:warn:fubini}).
\end{dictionary}

\begin{physics}
The double contraction is the pair of contour integrals $\oint_{x_1}\oint_{x_2}$ that converts a configuration
$a_1(x_1)\,b(z_b)\,a_2(x_2)$ into the relative coordinates of $x_1,x_2$ with respect to $z_b$; the two orders of integration
are the two orders in which the induced vertex operators $Y_W(a_1,z_{o_1})$, $Y_W(a_2,z_{o_2})$ act.
\end{physics}

\begin{status}
Definitions.
\end{status}

\begin{theorem}[Double contraction of the double delta]\label{l01:thm:kernelresidue}
For all orientations and every payload $f$: (a) on the residue graph the double kernel factorises,
$K_{o_1o_2}(-r_1-k_1-1,-r_2-k_2-1,r_1,r_2)(k_1+k_2) = \gamma_{o_1}(r_1,k_1)\gamma_{o_2}(r_2,k_2)$, and it vanishes off the total
antidiagonal; (b) the double contraction of $K_{o_1o_2}$ in either order is the double residue in the same order;
(c) if $f$ satisfies \lean{DoubleResidueFinite}, the two orders agree; (d) the outer-to-inner transport survives contraction:
$\Res\Res\big((z_{o_1}-z_{o_2})^NK\big)f = \Res\Res\big((x_1-x_2)^NK\big)f$ for every $f$.
\end{theorem}

\noindent File \leanpath{VOA/\allowbreak Foundations/\allowbreak FormalDist4DeltaContract.lean}.
\begin{leancode}
theorem doubleDelta_on_residue_graph (o₁ o₂ : Orientation)
    (outer₁ outer₂ k₁ k₂ : ℤ) :
    doubleDelta o₁ o₂
        (-outer₁ - k₁ - 1) (-outer₂ - k₂ - 1) outer₁ outer₂ (k₁ + k₂) =
      contractionCoeff o₁ outer₁ k₁ * contractionCoeff o₂ outer₂ k₂ := by

theorem doubleDelta_eq_zero_off_total_diag (o₁ o₂ : Orientation)
    (source₁ source₂ outer₁ outer₂ base : ℤ)
    (h : source₁ + source₂ + base ≠ -outer₁ - outer₂ - 2) :
    doubleDelta o₁ o₂ source₁ source₂ outer₁ outer₂ base = 0 := by

theorem kernelContract₁₂_doubleDelta (o₁ o₂ : Orientation) (f : FormalDist₃ W) :
    kernelContract₁₂ (doubleDelta o₁ o₂) f = doubleResidue₁₂ o₁ o₂ f := by

theorem kernelContract₂₁_doubleDelta (o₁ o₂ : Orientation) (f : FormalDist₃ W) :
    kernelContract₂₁ (doubleDelta o₁ o₂) f = doubleResidue₂₁ o₁ o₂ f := by

theorem doubleResidue_comm (o₁ o₂ : Orientation) (f : FormalDist₃ W)
    (hfin : DoubleResidueFinite o₁ o₂ f) :
    doubleResidue₁₂ o₁ o₂ f = doubleResidue₂₁ o₁ o₂ f := by

theorem kernelContract_doubleDelta_comm (o₁ o₂ : Orientation) (f : FormalDist₃ W)
    (hfin : DoubleResidueFinite o₁ o₂ f) :
    kernelContract₁₂ (doubleDelta o₁ o₂) f =
      kernelContract₂₁ (doubleDelta o₁ o₂) f := by

theorem kernelContract₁₂_doubleDelta_transport (o₁ o₂ : Orientation) (N : ℕ)
    (f : FormalDist₃ W) :
    kernelContract₁₂ (FormalDist₄.mulBySub₃₄Pow N (doubleDelta o₁ o₂)) f =
      kernelContract₁₂ (FormalDist₄.mulBySub₁₂Pow N (doubleDelta o₁ o₂)) f :=
\end{leancode}

\begin{dictionary}
\textbf{Hypotheses.} (a), (b), (d) have none. That (b) needs no finiteness is correct and not a junk-value accident: the
innermost sum over $\beta$ (\lean{baseIn}) has at most one nonzero term by the total-antidiagonal support, and the remaining
double sum is a bijective reindexing $(s_1,s_2)\mapsto(s_1-r_1,s_2-r_2)$ of the double-residue sum, and \lean{finsum} is
invariant under bijections in all cases. (c) needs \lean{DoubleResidueFinite}, exactly the joint finiteness of
\lonewarn{l01:warn:fubini}. (d) is Theorem~\ref{l01:thm:outerinner} with the equal kernels contracted against the same $f$; its
proof is the one-line term \lean{kernelContract₁₂_congr (mulBySub₃₄Pow_doubleDelta o₁ o₂ N) f}. The file also proves the
$(2,1)$-order analogue of (d). \textbf{Strength.} (d) is weaker than what is needed for locality, because the difference
operators still act on the kernel; moving them onto the output and onto the payload is the next theorem, and that is
where a finiteness hypothesis enters.
\end{dictionary}

\begin{physics}
(c) is the formal statement that the two contour integrals may be performed in either order when only finitely many
intermediate states contribute, which is the case for genuine field configurations (energy bounded below). (d) says that
before any integration is performed the separation of the relative coordinates equals the separation of the original
points.
\end{physics}

\textbf{Proof on paper.} (a) The convolution over \lean{split} has the single term $\text{split} = k_1$. (b) As explained in
the dictionary. (c) Theorem~\ref{l01:thm:fubini}. (d) Theorem~\ref{l01:thm:outerinner}. \qed

\begin{status}
Unconditional for (a), (b), (d); (c) conditional on \lean{DoubleResidueFinite}. Axioms clean (checked for
\lean{doubleDelta_on_residue_graph}, \lean{kernelContract₁₂_doubleDelta}, \lean{doubleResidue_comm},
\lean{kernelContract_doubleDelta_comm}, \lean{kernelContract₁₂_doubleDelta_transport}).
\end{status}

\subsection{Moving difference powers through the contraction; locality transport}

\begin{definition}[Shift-stable joint finiteness]\label{l01:def:shiftfinite}
A three-index family $F_{s_1,s_2,\beta}$ is \emph{triple-finite} if its support in $\Z^3$ is finite. For a kernel $K$, a payload
$f$ and integer offsets $d_1,d_2,d_3,d_4,e_1,e_3$, the \emph{shifted contraction summand} at output $(r_1,\beta',r_2)$ is
\begin{equation*}
  (s_1,s_2,\beta)\longmapsto K_{-s_1-1-d_1,\;-s_2-1-d_2,\;r_1-d_3,\;r_2-d_4}(\beta'-\beta)\cdot f_{s_1-e_1,\;\beta,\;s_2-e_3}.
\end{equation*}
The pair $(K,f)$ is \emph{shift-contract-finite} if for all offsets and all outputs this family is triple-finite.
\end{definition}

\noindent File \leanpath{VOA/\allowbreak Foundations/\allowbreak FormalDist4DeltaSupport.lean}.
\begin{leancode}
def TripleFinite (F : ℤ → ℤ → ℤ → W) : Prop :=
  (Function.support fun t : ℤ × (ℤ × ℤ) => F t.1 t.2.1 t.2.2).Finite

noncomputable def shiftedContractSummand (K : FormalDist₄ (FormalDist ℂ))
    (f : FormalDist₃ W) (d₁ d₂ d₃ d₄ e₁ e₃ : ℤ)
    (outer₁ baseOut outer₂ source₁ source₂ baseIn : ℤ) : W :=
  K (-source₁ - 1 - d₁) (-source₂ - 1 - d₂)
      (outer₁ - d₃) (outer₂ - d₄) (baseOut - baseIn) •
    f (source₁ - e₁) baseIn (source₂ - e₃)

def ShiftContractFinite (K : FormalDist₄ (FormalDist ℂ)) (f : FormalDist₃ W) : Prop :=
  ∀ d₁ d₂ d₃ d₄ e₁ e₃ outer₁ baseOut outer₂,
    TripleFinite (shiftedContractSummand K f d₁ d₂ d₃ d₄ e₁ e₃
      outer₁ baseOut outer₂)
\end{leancode}

\begin{dictionary}
\textbf{Why this shape.} Each application of a difference operator replaces a summand by a difference of two shifted
summands; distributing the subtraction through a nested \lean{finsum} is legitimate only if both pieces have finite
support (\lean{tripleFinsum_sub_distrib}). Requiring finiteness for \emph{all} offsets makes the predicate stable under
any number of such steps: the file proves that it survives \lean{FormalDist₄.mulBySub₃₄}, \lean{FormalDist₄.mulBySub₁₂} on the
kernel and \lean{mulBySub₁₃} on the payload, and their powers. The file header says plainly that this is ``stronger than any
single displayed identity, but exactly stable''.

\textbf{Could it be vacuous?} No. If $f$ has finite support, then for every offset the summand vanishes outside a
translate of $\operatorname{supp}f$, so $(K,f)$ is shift-contract-finite for \emph{every} $K$; so the predicate has many
nonzero instances. \textbf{Is it satisfied by the intended payloads?} That is the open question of
\lonewarn{l01:warn:shiftfinite}.
\end{dictionary}

\begin{physics}
The requirement that, on fixed external states, only finitely many intermediate modes contribute to a double contour
integral, uniformly under the finite shifts produced by multiplying by polynomials in the coordinates. In a CFT this
follows from the energy bound; formally it must be derived from lower truncation of the fields.
\end{physics}

\begin{status}
Definitions.
\end{status}

\begin{theorem}[Difference powers through contraction; locality transport]\label{l01:thm:transport}
Suppose $(K,f)$ is shift-contract-finite. Then for every $N\in\N$
\begin{align*}
  \Res\Res\big((z_{o_1}-z_{o_2})^NK\big)f &= (z_{o_1}-z_{o_2})^N\,\Res\Res\,Kf,\\
  \Res\Res\big((x_1-x_2)^NK\big)f &= \Res\Res\,K\big((x_1-x_2)^Nf\big).
\end{align*}
Consequently, for the double delta kernel $K = K_{o_1o_2}$ with $(K,f)$ shift-contract-finite,
$(z_{o_1}-z_{o_2})^N\Res\Res\,Kf = \Res\Res\,K(x_1-x_2)^Nf$, and if $(x_1-x_2)^Nf = 0$ then
$(z_{o_1}-z_{o_2})^N\Res\Res\,Kf = 0$.
\end{theorem}

\noindent File \leanpath{VOA/\allowbreak Foundations/\allowbreak FormalDist4DeltaSupport.lean}.
\begin{leancode}
theorem kernelContract₁₂_mulBySub₃₄Pow (K : FormalDist₄ (FormalDist ℂ))
    (f : FormalDist₃ W) (hfin : ShiftContractFinite K f) (N : ℕ) :
    kernelContract₁₂ (FormalDist₄.mulBySub₃₄Pow N K) f =
      mulBySub₁₃Pow N (kernelContract₁₂ K f) := by

theorem kernelContract₁₂_mulBySub₁₂Pow (K : FormalDist₄ (FormalDist ℂ))
    (f : FormalDist₃ W) (hfin : ShiftContractFinite K f) (N : ℕ) :
    kernelContract₁₂ (FormalDist₄.mulBySub₁₂Pow N K) f =
      kernelContract₁₂ K (mulBySub₁₃Pow N f) := by

theorem kernelContract₁₂_doubleDelta_locality_transport (o₁ o₂ : Orientation)
    (N : ℕ) (f : FormalDist₃ W)
    (hfin : ShiftContractFinite (doubleDelta o₁ o₂) f) :
    mulBySub₁₃Pow N (kernelContract₁₂ (doubleDelta o₁ o₂) f) =
      kernelContract₁₂ (doubleDelta o₁ o₂) (mulBySub₁₃Pow N f) := by

theorem kernelContract₁₂_doubleDelta_eq_zero_of_local (o₁ o₂ : Orientation)
    (N : ℕ) (f : FormalDist₃ W)
    (hfin : ShiftContractFinite (doubleDelta o₁ o₂) f)
    (hlocal : mulBySub₁₃Pow N f = 0) :
    mulBySub₁₃Pow N (kernelContract₁₂ (doubleDelta o₁ o₂) f) = 0 := by
\end{leancode}

\begin{dictionary}
\textbf{Decoding.} \lean{FormalDist₄.mulBySub₃₄Pow N K} is $(z_{o_1}-z_{o_2})^NK$ and \lean{FormalDist₄.mulBySub₁₂Pow N K} is
$(x_1-x_2)^NK$ on the kernel. The unqualified \lean{mulBySub₁₃Pow N} is the \emph{three-variable} operator of
Definition~\ref{l01:def:fd3}, acting on the output (where it is $(z_{o_1}-z_{o_2})^N$, first minus third coordinate) or on the
payload (where it is $(x_1-x_2)^N$). With these readings the four statements are the displayed ones.

\textbf{Every hypothesis.} \lean{hfin : ShiftContractFinite ...} is the support interface of
Definition~\ref{l01:def:shiftfinite}; \lean{hlocal} is locality of the payload in $(x_1,x_2)$. The first two theorems hold for an
arbitrary kernel $K$; the last two for the double delta, combining them with Theorem~\ref{l01:thm:kernelresidue}(d).

\textbf{Same, stronger, weaker?} As a formal statement, the last theorem is the locality-transport step of Li's proof
(the file header calls the target ``Li's local-system theorem (Proposition 3.2.9)'') with its finiteness
requirement made explicit. It is \emph{conditional}: it says nothing about payloads for which \lean{hfin} fails, and at this
commit no theorem establishes \lean{hfin} for payloads built from fields.
\end{dictionary}

\begin{physics}
If $a_1(x_1)$ and $a_2(x_2)$ are mutually local, then after re-centring both at $z_b$ the induced fields
$Y_W(a_1,z_{o_1})$ and $Y_W(a_2,z_{o_2})$ are mutually local in the relative coordinates: locality is preserved by the
passage from fields on $W$ to the vertex operators of the local system. Physically: the OPE of two chiral fields
defines new chiral fields that are again local with respect to each other, the operator-algebra version of the
statement that composite operators of a local theory are local.
\end{physics}

\textbf{Proof on paper.} For one factor $z_{o_1}-z_{o_2}$ on the kernel: the contraction of $(z_{o_1}-z_{o_2})K$ is a nested sum of
differences of two shifted summands; by \lean{hfin} both pieces have finite joint support, so the sum splits, and each
piece is the contraction of $K$ evaluated at a shifted output: together, $(z_{o_1}-z_{o_2})$ applied to the output. For one
factor $x_1-x_2$ on the kernel: shifting the kernel's source index by one is the same as shifting the payload's source index
by one after the bijective reindexing $s_i\mapsto s_i+1$ of the summation variable (residue pairing:
$\Res_x(xK)f = \Res_xK(xf)$); again \lean{hfin} legitimises splitting the difference. Induction on $N$ using the stability of
\lean{hfin}. The locality transport is the chain: output power $=$ (first theorem) contraction of the outer power on $K$
$=$ (Theorem~\ref{l01:thm:outerinner}) contraction of the inner power on $K$ $=$ (second theorem) contraction of $K$ against the
power on the payload. If the latter vanishes, the contraction of $K$ against $0$ is $0$. \qed

\begin{status}
\textbf{Conditional} on \lean{ShiftContractFinite (doubleDelta o₁ o₂) f} (and, for the last theorem, on payload locality, which
is the intended input). \texttt{\#print axioms} for all four $=$ \lean{[propext, Classical.choice, Quot.sound]}.
\end{status}

\begin{warning}\label{l01:warn:shiftfinite}
\textbf{The four-variable layer is not yet connected to the rest of the library.} When this chapter was first written, and still in the released source as far as point (1) is concerned (re-checked by a grep of the imports for this edition):
(1) no file outside \leanpath{VOA/\allowbreak Foundations/\allowbreak FormalDist4*.lean} imports \leanpath{FormalDist4.lean},
\leanpath{FormalDist4Delta.lean}, \leanpath{FormalDist4DeltaContract.lean} or \leanpath{FormalDist4DeltaSupport.lean}, except
the umbrella \leanpath{VOA/\allowbreak Foundations.lean}; in particular \leanpath{VOA/\allowbreak Core/\allowbreak LocalSystem.lean}
does not use them, although the header of \leanpath{FormalDist4.lean} then said the layer is ``needed in the proof of Li's
local-system theorem''. (The current header says more cautiously that its purpose is to ``support the double-residue
calculation in Li's local-system theorem''; a grep of the current tree confirms that only the umbrella imports the layer.) (2) The hypothesis \lean{ShiftContractFinite (doubleDelta o₁ o₂) f} has not been derived, for any
payload built from fields, from lower truncation (\lean{IsField}). The project's internal records of the local-system
route list exactly this support derivation as the next obligation, and claim neither it nor the induced locality of regular
fields. A rough count suggests that for the $(+,+)$ orientation the
summand is supported where $s_i\ge r_i$ and $\beta = \beta'-(s_1-r_1)-(s_2-r_2)$, so joint finiteness requires the payload to
vanish for very negative base exponent $\beta$, which is the kind of truncation fields provide; this is a heuristic, not a
proof. A referee should therefore treat Theorem~\ref{l01:thm:transport} as a correct conditional theorem whose hypothesis is
an open input, and the whole layer as infrastructure for future work.
\end{warning}

%======================================================================
\section{Verification record}\label{l01:sec:verification}
%======================================================================

\paragraph{What was run.} (This run dates from the first writing of the chapter, at an earlier commit, and was not re-run for this edition.) The tree at that commit was already built (the \texttt{.olean} files of all thirteen
modules are present and the working tree has no modifications under \texttt{VOA/}); we did \emph{not} re-run the aggregate
\texttt{lake build} for this chapter. We created a scratch file outside the repository importing
\lean{VOA.Foundations.FormalDist4DeltaSupport} and \lean{VOA.Foundations.IotaZExpansion} (which together import all
thirteen modules) and containing sixty \texttt{\#print axioms} commands, and ran it with the audit-safe invocation required
by the project (a bare \texttt{lake env lean} may drop \texttt{maxSynthPendingDepth}):
\begin{plaincode}
export PATH="$HOME/.elan/bin:$PATH"
lake env lean -DmaxSynthPendingDepth=3 Axioms.lean
\end{plaincode}

\paragraph{Result.} All sixty declarations report exactly the three standard axioms; there is no \lean{sorryAx} and no
\lean{VOA.pbw}. A representative excerpt of the output:
\begin{plaincode}
'VOA.local_iff_deltaSum' depends on axioms: [propext, Classical.choice, Quot.sound]
'VOA.expand_jump' depends on axioms: [propext, Classical.choice, Quot.sound]
'VOA.negPowJump_eq_deltaDeriv' depends on axioms: [propext, Classical.choice, Quot.sound]
'VOA.cbz_vandermonde' depends on axioms: [propext, Classical.choice, Quot.sound]
'VOA.expandPos_mul' depends on axioms: [propext, Classical.choice, Quot.sound]
'VOA.eq_zero_of_mulBySubPow_of_bddBelow' depends on axioms: [propext, Classical.choice, Quot.sound]
'VOA.subst_mulBySub' depends on axioms: [propext, Classical.choice, Quot.sound]
'VOA.localε_of_car' depends on axioms: [propext, Classical.choice, Quot.sound]
'VOA.mulBySub₂₃Pow_eq_sum' depends on axioms: [propext, Classical.choice, Quot.sound]
'VOA.FormalDist₄Delta.mulBySub₃₄Pow_doubleDelta' depends on axioms: [propext, Classical.choice, Quot.sound]
'VOA.FormalDist₄Delta.kernelContract₁₂_doubleDelta_eq_zero_of_local' depends on axioms: [propext,
 Classical.choice,
 Quot.sound]
\end{plaincode}
The full list of the sixty names checked is:
\begin{plaincode}
VOA.mode_eq_res, VOA.res_deriv, VOA.deriv_lpSMul, VOA.res_lpSMul_deriv, VOA.algSum_add,
VOA.algSum_smul, VOA.mulBySub_deltaDeriv, VOA.mulBySubPow_deltaDeriv_self,
VOA.mulBySubPow_deltaDeriv_eq_zero, VOA.eq_deltaDeriv_zero_of_mulBySub_eq_zero,
VOA.exists_deltaSum_of_mulBySubPow_eq_zero, VOA.local_iff_deltaSum,
VOA.local_of_bracketDist_eq_zero, VOA.bracketDistε_one, VOA.localε_one_iff, VOA.bracketDistε_swap,
VOA.localε_iff_deltaSum, VOA.localε_of_car, VOA.localε_neg_one_of_anticomm, VOA.gBinom_neg_one,
VOA.expand_jump, VOA.mulBySub_deltaForm, VOA.eq_zero_of_mulBySub_of_bddBelow,
VOA.eq_zero_of_mulBySubPow_of_bddBelow, VOA.gBinom_succ_pascal, VOA.cbz_pascal,
VOA.subst_summand_finite, VOA.subst_mulBySub, VOA.subst_mulBySubPow, VOA.mulBySub_expandPos,
VOA.mulBySub_expandNeg, VOA.negPowJump_eq_deltaDeriv, VOA.cbz_vandermonde, VOA.expandPos_mul,
VOA.expandNeg_mul, VOA.expandPos_zero, VOA.expandNeg_zero, VOA.mulBySubPow_expandPos,
VOA.mulBySubPow_expandPos_neg, VOA.mulBySubPow_negPowJump_eq_zero, VOA.negPowJump_deltaSum,
VOA.mulBySub₂₃_eq, VOA.mulBySub₂₃Pow_eq_sum, VOA.FormalDist₄.finsum_finsum_comm_of_joint_finite,
VOA.FormalDist₄.coeffFinsum₂_comm, VOA.FormalDist₄Delta.singleDelta_relation,
VOA.FormalDist₄Delta.mulBySub₃₄_doubleDelta, VOA.FormalDist₄Delta.mulBySub₃₄Pow_doubleDelta,
VOA.FormalDist₄Delta.singleDeltaContract_eq_residueSubst,
VOA.FormalDist₄Delta.doubleDelta_on_residue_graph,
VOA.FormalDist₄Delta.contractionCoeff_neg_zero_neg_one,
VOA.FormalDist₄Delta.contractionCoeff_neg_neg_one_zero,
VOA.FormalDist₄Delta.kernelContract₁₂_doubleDelta, VOA.FormalDist₄Delta.doubleResidue_comm,
VOA.FormalDist₄Delta.kernelContract_doubleDelta_comm,
VOA.FormalDist₄Delta.kernelContract₁₂_doubleDelta_transport,
VOA.FormalDist₄Delta.kernelContract₁₂_mulBySub₃₄Pow,
VOA.FormalDist₄Delta.kernelContract₁₂_mulBySub₁₂Pow,
VOA.FormalDist₄Delta.kernelContract₁₂_doubleDelta_locality_transport,
VOA.FormalDist₄Delta.kernelContract₁₂_doubleDelta_eq_zero_of_local
\end{plaincode}

\paragraph{Release audit 2026-10-27.} \texttt{Notes/\allowbreak Textbook/\allowbreak referee/\allowbreak RELEASE\_AUDIT\_FOUNDATIONS.md}.
For this chapter its relevant findings are: \lean{Local} is used verbatim in the standard sense (it checked the convention
of \lean{mulBySub}), and the structures \lean{PreVertexAlgebra}/\lean{BarePreVA} that use \lean{Local} as an axiom lack
translation covariance, so they are not vertex algebras in the standard sense (Chapter~\ref{ch:L03}); after the audit
the library gave them these names (formerly \lean{VertexAlgebra} and \lean{BareVA}) and reserved \lean{VertexAlgebra} for
the translation-covariant notion.

\paragraph{Reading the result honestly.} A clean \texttt{\#print axioms} does not report explicit hypotheses. The
theorems of this chapter with genuine hypotheses are: those carrying a Laurent-polynomial hypothesis
(\lean{deriv_lpSMul}, \lean{res_lpSMul_deriv}), a lower-truncation hypothesis (\lean{eq_zero_of_mulBySubPow_of_bddBelow},
\lean{subst_mulBySub}, \lean{subst_mulBySubPow}), a joint-finiteness hypothesis (\lean{finsum_finsum_comm_of_joint_finite},
\lean{coeffFinsum₂_comm}, \lean{doubleResidue_comm}, \lean{kernelContract_doubleDelta_comm}), and the shift-stable
support hypothesis \lean{ShiftContractFinite} (the four theorems of Theorem~\ref{l01:thm:transport}). The first three
kinds are meaningfulness or genuine mathematical conditions that are part of the textbook statements. The last is an
open input at this commit (\lonewarn{l01:warn:shiftfinite}).

\paragraph{Downstream use (checked by \texttt{grep} at this commit).} \lean{Local} is the locality axiom of
\lean{PreVertexAlgebra}, \lean{BarePreVA} and \lean{VertexOperatorAlgebra}; \lean{eq_zero_of_mulBySubPow_of_bddBelow} is used in nine files of \texttt{VOA/Core} (weak associativity, Goddard
uniqueness, reconstruction); \lean{subst_mulBySubPow} in \leanpath{Core/\allowbreak Skew.lean}; \lean{mulBySub₂₃Pow_eq_sum} in
\leanpath{Core/\allowbreak Dong.lean} and \leanpath{Core/\allowbreak DongSuper.lean}; \lean{expandPos_mul} in
\leanpath{Examples/\allowbreak LatticeA1.lean}; \lean{localε_of_car} in three free-field example files; \lean{LaurentPoly} in
\leanpath{Examples/\allowbreak VirasoroKacLaurent.lean}. Not used outside the thirteen files of this chapter: \lean{formalDelta}, \lean{delta₂},
\lean{lpSMul} and its two theorems, \lean{binomialExpand}, \lean{algSum}, \lean{negPowJump_eq_deltaDeriv}, and the whole
four-variable layer.

\paragraph{Docstring discrepancies found while writing this chapter.}
All but one of these have since been corrected in the source (a grep of the current tree finds none of (1), (3), (4), (5),
(6), and only half of (2)); they are kept here as a record of what the first edition reported.
(1) \leanpath{IotaZExpansion.lean} listed a lemma \lean{expandPos_mul_zero} that does not exist (Theorem~\ref{l01:thm:mult}).
(2) \leanpath{IteratedLaurent.lean} described \lean{IsBddBelow₂} as the class $\mathcal W((z_1))((z_2))$; it is the larger class
$\mathcal W[[z_1^{\pm1}]]((z_2))$ (Definition~\ref{l01:def:bddbelow}). The definition's docstring is now correct; the docstring
of \lean{eq_zero_of_mulBySubPow_of_bddBelow} still says $\mathcal W((z_1))((z_2))$.
(3) \leanpath{Expansion.lean} said the general-$n$ multiplicativity of the expansions is ``provided by
\lean{subst_mulBySubPow}''; the multiplicativity of the expansions is \lean{expandPos_mul}/\lean{expandNeg_mul} in
\leanpath{IotaZExpansion.lean}, while \lean{subst_mulBySubPow} is the multiplicativity of the substitution, a different
statement.
(4) The heuristic ``no finite-dimensional fermionic toy example'' in \leanpath{Locality.lean}
(\lonewarn{l01:warn:finitefermion}).
(5) \leanpath{FormalDist4.lean} said the layer is needed for Li's theorem, but it is not (yet) imported by
\leanpath{Core/\allowbreak LocalSystem.lean} (\lonewarn{l01:warn:shiftfinite}).
(6) \leanpath{FormalDist4Delta.lean} stated its central identity in variables $(z_3-z_0)^NK = (z_1-z_4)^NK$ that do not
match the Lean coordinates; the Lean statement is $(z_{o_1}-z_{o_2})^NK = (x_1-x_2)^NK$ (Theorem~\ref{l01:thm:outerinner}).
None of these affects the truth of any Lean statement.

%======================================================================
\section{Referee's checklist}\label{l01:sec:checklist}
%======================================================================

\begingroup\small
\begin{longtable}{>{\raggedright\arraybackslash}p{0.30\textwidth}>{\raggedright\arraybackslash}p{0.33\textwidth}>{\raggedright\arraybackslash}p{0.2\textwidth}>{\raggedright\arraybackslash}p{0.09\textwidth}}
\toprule
Lean declaration & Intended meaning & Hypotheses / status & Check\\
\midrule\endhead
\lean{FormalDist}, \lean{FormalDist₂} & $\loneW[[z^{\pm1}]]$, $\loneW[[z_1^{\pm1},z_2^{\pm1}]]$; no product & definition & Def.~\ref{l01:def:formaldist}, \ref{l01:def:formaldist2}\\
\lean{mode} & $f_{(m)} = f_{-m-1}$ & definition; \textbf{key convention} & Def.~\ref{l01:def:mode}\\
\lean{res}, \lean{deriv} & $\Res_z$, $\partial_z$ & definition & Def.~\ref{l01:def:mode}\\
\lean{mode_eq_res}, \lean{res_deriv} & $f_{(m)} = \Res z^mf$; $\Res\partial f = 0$ & none; unconditional & Prop.~\ref{l01:prop:moderes}\\
\lean{lpSMul} & Laurent polynomial times distribution & definition; junk value if $g$ infinite & Def.~\ref{l01:def:lpsmul}\\
\lean{deriv_lpSMul}, \lean{res_lpSMul_deriv} & Leibniz; integration by parts & $g$ Laurent polynomial & Thm.~\ref{l01:thm:leibniz}, \ref{l01:thm:byparts}\\
\lean{gBinom}, \lean{cbz}, \lean{negOnePow} & $\binom nk$, $n\in\Z$; $C(n,k) = 0$ for $k<0$; $(-1)^n$ & definition & Def.~\ref{l01:def:gbinom}\\
\lean{cbz_vandermonde} & $\sum_iC(m,i)C(n,k-i) = C(m+n,k)$, all $m,n,k\in\Z$ & none; unconditional & Thm.~\ref{l01:thm:vandermonde}\\
\lean{mulBySub}, \lean{mulBySubPow} & multiplication by $(z_1-z_2)^N$ & definition & Def.~\ref{l01:def:formaldist2}\\
\lean{deltaDeriv}, \lean{cCoeff} & $c(z_2)\partial_{z_2}^j\delta(z_1,z_2)$ (plain $\partial^j$); $\Res_{z_1}(z_1-z_2)^jf$ & definition & Def.~\ref{l01:def:deltaderiv}\\
\lean{local_iff_deltaSum} & $(z_1-z_2)^Nf = 0\iff f = \sum_{j<N}c^j\partial^j\delta$ (Gannon 5.1.2, Kac 2.3) & none; unconditional & Thm.~\ref{l01:thm:deltasum}\\
\lean{expandPos}, \lean{expandNeg}, \lean{deltaForm} & $\iota_{z_1,z_2}$, $\iota_{z_2,z_1}$ of $(z_1-z_2)^n$; $\delta(z_1,z_2)$ & definition & Def.~\ref{l01:def:expand}\\
\lean{expand_jump} & $\iota_{12}-\iota_{21}$ of $(z_1-z_2)^{-1}$ is $\delta$ & none; unconditional & Thm.~\ref{l01:thm:jump}\\
\lean{negPowJump_}\allowbreak\lean{eq_deltaDeriv} & jump of $(z_1-z_2)^{-k-1}$ is $\partial^{(k)}\delta$ & none; unconditional & Thm.~\ref{l01:thm:alljump}\\
\lean{expandPos_mul}, \lean{expandNeg_mul} & $\iota(z_1-z_2)^m\iota(z_1-z_2)^n = \iota(z_1-z_2)^{m+n}$ & none; \lean{convolPos} is the genuine product (paper) & Thm.~\ref{l01:thm:mult}\\
\lean{bracketDist}, \lean{Local}, \lean{normalOrdered} & $[a(z_1),b(z_2)]$; mutual locality; ${:}a(z_1)b(z_2){:}$ & definition & Def.~\ref{l01:def:bracket}\\
\lean{IsField}, \lean{nprodCoeff}, \lean{nprod} & field (lower truncation); same-variable ${:}a(z)b(z){:}$, the $(-1)$-st product & definition; \lean{∑ᶠ} finite because both arguments are fields & Def.~\ref{l01:def:field}\\
\lean{localε_iff_}\allowbreak\lean{deltaSum}, \lean{localε_of_car} & super-locality; CAR/CCR $\Rightarrow$ local & none; unconditional & Thm.~\ref{l01:thm:epsdeltasum}, \ref{l01:thm:car}\\
\lean{IsBddBelow₂} & lower truncation in $z_2$, uniform in $z_1$ & definition (weaker than $\loneW((z_1))((z_2))$) & Def.~\ref{l01:def:bddbelow}\\
\lean{eq_zero_of_}\allowbreak\lean{mulBySubPow_of_bddBelow} & $z_1-z_2$ not a zero-divisor on lower-truncated series & lower truncation; unconditional & Thm.~\ref{l01:thm:cancel}\\
\lean{subst}, \lean{subst_mulBySub(Pow)} & $f(z_0+z_2,z_2)$; $(z_1-z_2)\mapsto z_0$ & lower truncation; junk value otherwise & Def.~\ref{l01:def:subst}, Thm.~\ref{l01:thm:subst}\\
\lean{mulBySub₂₃Pow_}\allowbreak\lean{eq_sum} & $(z_2-z_3)^N = \sum\binom Ni(z_2-z_1)^i(z_1-z_3)^{N-i}$ & none; unconditional & Thm.~\ref{l01:thm:fd3binom}\\
\lean{finsum_finsum_}\allowbreak\lean{comm_of_joint_finite} & Fubini for algebraic sums & joint finite support (necessary) & Thm.~\ref{l01:thm:fubini}\\
\lean{mulBySub₃₄Pow_}\allowbreak\lean{doubleDelta} & $(z_{o_1}-z_{o_2})^NK = (x_1-x_2)^NK$ & none; unconditional & Thm.~\ref{l01:thm:outerinner}\\
\lean{kernelContract₁₂_}\allowbreak\lean{doubleDelta}, \lean{doubleResidue_comm} & contraction $=$ double residue; order exchange & none; \lean{DoubleResidueFinite} for exchange & Thm.~\ref{l01:thm:kernelresidue}\\
\lean{kernelContract₁₂_}\allowbreak\lean{doubleDelta_}\allowbreak\lean{eq_zero_of_local} & locality transported through double delta & \textbf{conditional} on \lean{ShiftContractFinite} (open input); layer unused & Thm.~\ref{l01:thm:transport}, p.~\pageref{l01:warn:shiftfinite}\\
\bottomrule
\end{longtable}
\endgroup

%======================================================================
\section{Notes and references}\label{l01:sec:notes}
%======================================================================

\textbf{Formal calculus.} The formal-distribution approach to vertex algebras, with the delta function, the two expansions
and the structure theorem for local distributions, is developed in V.~G.~Kac, \emph{Vertex Algebras for Beginners},
University Lecture Series 10, American Mathematical Society, second edition (1998), Chapters~1--2 (the structure theorem
is his Theorem~2.3; Dong's lemma is Lemma~3.2; Goddard uniqueness and reconstruction are in Chapter~4). The same material
in the language of fields and with geometric motivation is in E.~Frenkel and D.~Ben-Zvi, \emph{Vertex Algebras and
Algebraic Curves}, Mathematical Surveys and Monographs 88, American Mathematical Society, second edition (2004),
Chapter~1. The most thorough treatment of formal calculus, including the binomial convention, the three-term delta-function
identity and the substitution $z_1\mapsto z_0+z_2$, is J.~Lepowsky and H.~Li, \emph{Introduction to Vertex Operator Algebras
and Their Representations}, Progress in Mathematics 227, Birkh\"auser (2004), Chapter~2; locality, weak associativity and
local systems are in Chapters~3 and~5 of the same book. The original source of the formal-variable calculus in this
area is I.~Frenkel, J.~Lepowsky and A.~Meurman, \emph{Vertex Operator Algebras and the Monster}, Pure and Applied
Mathematics 134, Academic Press (1988), Chapters~2 and~8.

\textbf{The library's primary reference.} The docstrings of the files of this chapter follow T.~Gannon, \emph{Moonshine
beyond the Monster: The Bridge Connecting Algebra, Modular Forms and Physics}, Cambridge University Press (2006),
§5.1.2 (formal distributions, Paradoxes~5.1 and~5.2, the residue and integration by parts (5.1.3), the binomial
convention (5.1.4), Proposition~5.1.2, normal ordering (5.1.6)).

\textbf{Local systems.} H.~Li, ``Local systems of vertex operators, vertex superalgebras and modules'', \emph{Journal of
Pure and Applied Algebra} 109 (1996), 143--195 (arXiv:hep-th/9406185), is the source of the construction that the
four-variable layer of Section~\ref{l01:sec:four} is intended to formalise. Dong's lemma is due to C.~Dong; see Li's paper
and Kac's book for proofs.

\textbf{Physics.} Radial quantisation, the operator product expansion and the relation between OPEs and commutators
are explained in P.~Di~Francesco, P.~Mathieu and D.~S\'en\'echal, \emph{Conformal Field Theory}, Springer (1997),
Chapters~5--6, and in P.~Ginsparg, ``Applied conformal field theory'', in \emph{Fields, Strings and Critical Phenomena}
(Les Houches 1988), North-Holland (1990), arXiv:hep-th/9108028. The free-fermion and $\beta\gamma$/$bc$ systems mentioned in
Section~\ref{l01:sec:super} are standard; see J.~Polchinski, \emph{String Theory}, Vol.~1, Cambridge University Press (1998),
Chapter~2.

\textbf{Project records.} The trust model (no project axiom;
\lean{VOA.pbw}, formerly a project axiom, is now a theorem and is not used here) is explained in
Chapter~\ref{ch:01}. The status of the four-variable layer is the one stated in Warning~\ref{l01:warn:shiftfinite}.

\endgroup
